% ============================================================
% 拓展答案册（S5 试产臂汇编草案）—— 源＝拓展册 各件 main.tex（只读）；工具/答案抽册器.py 逐件抽册口径，同序同号逐片保持
%% 源件 md5 见 逐件汇总.json；值/详解逐字照源（钉值硬断言）；题面不重印
%% 双档：ansbook-true＝详解印本／ansbook-false＝纯值速查本（\ansbookdetail 件面开关）
% 编译：xelatex <壳> ×2，cwd＝本目录；sty＝qp-m3.sty 副本（md5 7c3930362be8a0a2bdf21bbf8ac16573 锁同）
% ============================================================
\documentclass[fontset=none]{ctexart}
\usepackage{qp-m3}
% —— 册面补钉：pifont（源件值域含 \ding{51} 勾形；qp-m3 不供给，照 ROUTING 承源件同精神挂载，未用件零副作用） ——
\usepackage{pifont}
\xeCJKDeclareCharClass{Default}{"2212}%
\setCJKfamilyfont{nscblk}[Path=C:/提示词/工作区/字替对照-0909/variantF/fonts/,BoldFont=NSC-w600.ttf]{NSC-w500.ttf}%
\xeCJKDeclareSubCJKBlock{nscblk}{"25B1}%
\newcommand{\pxparallelogram}{{\CJKfamily{nscblk}▱}}%
\xeCJKsetup{CJKglue={\hskip 0.10em plus 0.08em minus 0.05em}}%
\ansblockgrayfalse
\newif\ifansbookdetail
\ifdefined\ansbookdetail
  \ifnum\ansbookdetail=0 \ansbookdetailfalse\else\ansbookdetailtrue\fi
\else
  \ansbookdetailtrue
\fi
\renewcommand{\qpzhangming}{第二章\quad 平面解析几何}%
\renewcommand{\qpjianming}{拓展答案册}%

\begin{document}
\zhangtitle{拓展答案册}
\par\glueguard{1}\addvspace{2pt}\noindent\biaoqian{【纯答案册】}{\fangsong 题面不重印\quad 逐片印面号与正册同号同序}\par\addvspace{2pt}

\begin{multicols}{2}
\emergencystretch=1em

\par\Needspace*{3\baselineskip}\noindent\biaoqian{【上册】}\par

\begin{ansblock}[2章-拓-课时01-T1]
% ans:2章-拓-课时01-T1
\ansitem{T1}{\((1)\)见详解；\((2)\) 等号成立\(\Leftrightarrow C\)的坐标\(x\in [-3,2]\)；等号不成立\(\Leftrightarrow x<-3\)或\(x>2\)．}
\ifansbookdetail
\ansnote{详解}{\((1)\) 设\(C\)的坐标为\(x_{3}\)．因为 \(x_{1}-x_{2}=(x_{1}-x_{3})+(x_{3}-x_{2})\)，由绝对值三角不等式 \(d(A,B)=|x_{1}-x_{2}|=|(x_{1}-x_{3})+(x_{3}-x_{2})|\leqslant |x_{1}-x_{3}|+|x_{3}-x_{2}|=d(A,C)+d(B,C)\)，得证． \((2) |x_{1}-x_{2}|=|(x_{1}-x_{3})+(x_{3}-x_{2})|\)取等号\(\Leftrightarrow (x_{1}-x_{3})(x_{3}-x_{2})\geqslant 0\)，即\(x_{3}\)介于\(x_{1}\)与\(x_{2}\)之间（含端点）． 此处\(x_{1}=-3\)，\(x_{2}=2\)，故等号成立\(\Leftrightarrow -3\leqslant x\leqslant 2\)；等号不成立（即\(d(A,B)<d(A,C)+d(B,C)\)）\(\Leftrightarrow x<-3\)或\(x>2\)．}
\fi
\end{ansblock}

\begin{ansblock}[2章-拓-课时01-T2]
% ans:2章-拓-课时01-T2
\ansitem{T2}{\((1)\)见详解；\((2)\) 等号成立\(\Leftrightarrow C\)在以\(AB\)为对角线的矩形内（含边界），即\((x\_C-x_{1})(x\_C-x_{2})\leqslant 0\)且\((y\_C-y_{1})(y\_C-y_{2})\leqslant 0\)；等号不成立\(\Leftrightarrow C\)在该矩形外．}
\ifansbookdetail
\ansnote{详解}{\((1)\) 设\(C(x_{3},y_{3})\)．\(x_{1}-x_{2}=(x_{1}-x_{3})+(x_{3}-x_{2})\)，\(y_{1}-y_{2}=(y_{1}-y_{3})+(y_{3}-y_{2})\)，两次用绝对值三角不等式： \(d(A,B)=|x_{1}-x_{2}|+|y_{1}-y_{2}|\leqslant (|x_{1}-x_{3}|+|x_{3}-x_{2}|)+(|y_{1}-y_{3}|+|y_{3}-y_{2}|)=d(A,C)+d(B,C)\)，得证． \((2)\) 横、纵两个方向分别取等：第一个不等式取等\(\Leftrightarrow (x_{1}-x_{3})(x_{3}-x_{2})\geqslant 0\)，第二个取等\(\Leftrightarrow (y_{1}-y_{3})(y_{3}-y_{2})\geqslant 0\)， 故总等号成立\(\Leftrightarrow x_{3}\)介于\(x_{1}\)、\(x_{2}\)之间且\(y_{3}\)介于\(y_{1}\)、\(y_{2}\)之间，即\(C\)落在以\(AB\)为对角线（边平行于坐标轴）的矩形内部或边界上； 只要有一个方向不取等，就有\(d(A,B)<d(A,C)+d(B,C)\)，即\(C\)在该矩形外．}
\fi
\end{ansblock}

\begin{ansblock}[2章-拓-课时02-T1]
% ans:2章-拓-课时02-T1
\ansitem{T1}{\((1) \theta \in [0,\pi /2)\)时斜率由\(0\)递增并趋向\(+\infty \)；\(\theta \in (\pi /2,\pi )\)时斜率由\(-\infty \)递增趋向\(0\)；\((2)\) 不能，理由见详解．}
\ifansbookdetail
\ansnote{详解}{\((1) k=tan\theta \)在\([0,\pi /2)\)上单调递增，\(k\)从\(0\)增大到正无穷大；在\((\pi /2,\pi )\)上也单调递增，\(k\)从负无穷大增大到\(0\)． \((2)\) 不能．例如取\(\theta _{1}=45^{\circ}\)，\(\theta _{2}=120^{\circ}\)，\(\theta _{1}<\theta _{2}\)，但\(k_{1}=1>k_{2}=-\sqrt{3}\)，即倾斜角大者斜率反而小； 斜率关于倾斜角只在每一段\([0,\pi /2)\)与\((\pi /2,\pi )\)内分别递增，整体上不具有单调性．}
\fi
\end{ansblock}

\begin{ansblock}[2章-拓-课时02-T2]
% ans:2章-拓-课时02-T2
\ansitem{T2}{\((1) k\_AB=0\)（倾斜角\(0\)）；\(k\_BC=\sqrt{3}\)（倾斜角\(\pi /3\)）；\(k\_AC=\sqrt{3}/3\)（倾斜角\(\pi /6\)）；\((2) [\pi /6, \pi /3]\)．}
\ifansbookdetail
\ansnote{详解}{\((1) k\_AB=(1-1)/(1+1)=0\)，\(k\_BC=(\sqrt{3}+1-1)/(2-1)=\sqrt{3}\)，\(k\_AC=(\sqrt{3}+1-1)/(2+1)=\sqrt{3}/3\)． 因斜率等于倾斜角的正切值且倾斜角\(\alpha \in [0,\pi )\)，故\(AB\)、\(BC\)、\(AC\)的倾斜角依次为\(0\)、\(\pi /3\)、\(\pi /6\)． \((2)\) 当直线\(CD\)绕点\(C\)由\(CA\)逆时针旋转到\(CB\)时，\(CD\)与线段\(AB\)恒有交点（即\(D\)在\(AB\)上），此时\(k\_CD\)由\(k\_AC\)增大到\(k\_BC\)， \(k\_CD\in [\sqrt{3}/3, \sqrt{3}]\)，对应倾斜角\(\alpha \in [\pi /6, \pi /3]\)．}
\fi
\end{ansblock}

\begin{ansblock}[2章-拓-课时02-T3]
% ans:2章-拓-课时02-T3
\ansitem{T3}{\((1) 0^{\circ}\)；\((2) 90^{\circ}\)；\((3) 45^{\circ}\)．}
\ifansbookdetail
\ansnote{详解}{\((1)\) 纵坐标同为\(c\)，直线与\(x\)轴平行，倾斜角为\(0^{\circ}\)（\(k=0\)）； \((2)\) 横坐标同为\(a\)，直线垂直于\(x\)轴，斜率不存在，倾斜角为\(90^{\circ}\)； \((3) k=[(c+a)-(b+c)]/(a-b)=(a-b)/(a-b)=1\)（\(a\neq b\)），\(tan\alpha =1\)，\(\alpha \in [0^{\circ},180^{\circ})\)，故倾斜角为\(45^{\circ}\)．}
\fi
\end{ansblock}

\begin{ansblock}[2章-拓-课时02-T4]
% ans:2章-拓-课时02-T4
\ansitem{T4}{\([-1/6, 5/3]\)}
\ifansbookdetail
\ansnote{详解}{\((y+1)/(x+1)=[y-(-1)]/[x-(-1)]\) 的几何意义是过\(M(x,y)\)与定点\(N(-1,-1)\)两点的直线的斜率，其中点\(M\)在线段\(y=-2x+8\)（\(x\in [2,5]\)）上运动． 线段两端点为\((2,4)\)与\((5,-2)\)．记\(k(x)=(y+1)/(x+1)=(-2x+9)/(x+1)=-2+11/(x+1)\)，在\(x\in [2,5]\)上随\(x\)增大而单调递减， 故最大值在\(x=2\)处取得：\(k=(4+1)/(2+1)=5/3\)；最小值在\(x=5\)处取得：\(k=(-2+1)/(5+1)=-1/6\)． 所以\((y+1)/(x+1)\)的取值范围为\([-1/6, 5/3]\)．}
\fi
\end{ansblock}

\begin{ansblock}[2章-拓-课时04-T1]
% ans:2章-拓-课时04-T1
\ansitem{T1}{\((1) k\geqslant 0\)；\((2) S\)的最小值为\(16\)，此时直线\(l\)的方程为\(x-2y+8=0\)}
\ifansbookdetail
\ansnote{详解}{\((1)\) 将方程整理为\(k(x+4)+(-y+2)=0\)，令\(x+4=0\)且\(y=2\)，得直线恒过定点\((-4,2)\)（在第二象限）． 斜率\(k=tan\theta \)．过\((-4,2)\)的直线不经过第四象限，当且仅当它\("\)向右上或水平张开\("\)：由\(l_{1}\)（过\((-4,2)\)与原点方向的临界位置）逆时针转到竖直位置均不进入第四象限，即倾斜角\(\theta \in [0^{\circ},90^{\circ}]\)，故\(k\geqslant 0\)（\(k\)不存在时直线\(x=-4\)也不过第四象限，但它不在\(y=kx+b\)的参数范围内——题设方程本身取遍\(k\in R\)时不含竖直线，而\(x=-4\)无法由该方程表示，故不另计）； \(k<0\)时直线向右下倾斜，自第二象限的点\((-4,2)\)出发必穿过第四象限，不合．所以\(k\geqslant 0\)． \((2)\) 由题意\(A\)、\(B\)存在且\(a<0\)、\(b>0\)（竖直线\(x=-4\)不与\(y\)轴相交，此时\(k=0\)给\(y=2\)与\(x\)轴不相交，均已在端点情形排除，直线可设为截距式\(x/a+y/b=1\)，\(a<0\)，\(b>0\)）． \((-4,2)\)在直线上：\(-4/a+2/b=1\)．因\(a<0\)，\(b>0\)，\(-4/a>0\)，\(2/b>0\)， \(1=-4/a+2/b\geqslant 2\sqrt{}((-4/a)\cdot (2/b))=2\sqrt{}(-8/(ab))\)，两边平方得\(1\geqslant -32/(ab)\)，即\(-ab\geqslant 32\)，\(ab\leqslant -32\)． \(S=(1/2)|OA|\cdot |OB|=(1/2)(-a)b=-ab/2\geqslant 16\)， 当且仅当\(-4/a=2/b=1/2\)，即\(a=-8\)，\(b=4\)时取等号． 此时直线为\(x/(-8)+y/4=1\)，化简得\(x-2y+8=0\)． 所以\(S\)的最小值为\(16\)，此时直线\(l\)的方程为\(x-2y+8=0\)．}
\fi
\end{ansblock}

\begin{ansblock}[2章-拓-课时05-T6]
% ans:2章-拓-课时05-T6
\ansitem{T6}{\((1) 2x+y-6=0\)或\(8x+y-12=0\)；\((2) 8\)，\(4x+y-8=0\)}
\ifansbookdetail
\ansnote{详解}{设\(A(a,0)\)、\(B(0,b)\)，\(a>0\)，\(b>0\)，\(l\)：\(x/a+y/b=1\)，代入\(P(1,4)\)得\(1/a+4/b=1\)（\(*\)）． \((1)\) 又\((1/2)ab=9\)即\(ab=18\)；由\((*)\)去分母\(b+4a=ab=18\)，\(b=18-4a\)，\(a(18-4a)=18\)，\(2a^{2}-9a+9=0\)，\(a=3\)或\(3/2\)； 对应\(b=6\)或\(12\)，\(l\)为\(x/3+y/6=1\)或\(x/(3/2)+y/12=1\)，即\(2x+y-6=0\)或\(8x+y-12=0\)． \((2) S=(1/2)ab=(1/2)ab\cdot (1/a+4/b)^{2}=(1/2)(8+b/a+16a/b)\geqslant (1/2)(8+2\sqrt{b/a\cdot 16a/b})=8\)， 当且仅当\(b/a=16a/b\)即\(b=4a\)，代\((*)\)：\(1/a+1/a=1\)，\(a=2\)、\(b=8\)时取等，\(S\)最小值\(8\)，此时\(l\)：\(x/2+y/8=1\)即\(4x+y-8=0\)．}
\fi
\end{ansblock}

\begin{ansblock}[2章-拓-课时05-T7]
% ans:2章-拓-课时05-T7
\ansitem{T7}{\((1)\) 最小值\(8\)，\(x/4+y/2=1\)；\((2)\) 最小值\(4\)，\(x/3+y/3=1\)}
\ifansbookdetail
\ansnote{详解}{\((1)\) 设\(l\)：\(x/a+y/b=1\)（\(a>0\)，\(b>0\)），过\(P\)得\(2/a+1/b=1\)． \(1=2/a+1/b\geqslant 2\sqrt{}(2/(ab))\)，得\(ab\geqslant 8\)，当且仅当\(2/a=1/b\)即\(a=4\)、\(b=2\)时取等； \(|OA|\cdot |OB|=ab\)的最小值为\(8\)，此时\(l\)：\(x/4+y/2=1\)． \((2)\) 由\((1)\)得\(b=a/(a-2)\)（\(a>2\)）．\(A(a,0)\)、\(B(0,b)\)， \(|PA|=\sqrt{}((a-2)^{2}+1)\)，\(|PB|=\sqrt{}(4+(b-1)^{2})\)，而\(b-1=a/(a-2)-1=2/(a-2)\)， \(|PA|\cdot |PB|=\sqrt{}((a-2)^{2}+1)\cdot \sqrt{}(4+4/(a-2)^{2})=2\sqrt{}((a-2)^{2}+1/(a-2)^{2}+2)\geqslant 2\sqrt{2+2}=4\)， 当且仅当\((a-2)^{2}=1/(a-2)^{2}\)即\(a-2=1\)、\(a=3\)时取等，此时\(b=3\)； \(|PA|\cdot |PB|\)的最小值为\(4\)，此时\(l\)：\(x/3+y/3=1\)．}
\fi
\end{ansblock}

\begin{ansblock}[2章-拓-课时05-T8]
% ans:2章-拓-课时05-T8
\ansitem{T8}{\((1) y=(3/4)x\pm 3\)；\((2) m\neq 1\)时\(y=(1/(m-1))(x-1)\)，\(m=1\)时\(x=1\)；\((3) x+y-1=0\)或\(x-y-7=0\)或\(y=-(3/4)x\)}
\ifansbookdetail
\ansnote{详解}{\((1)\) 设\(l\)：\(y=(3/4)x+b\)，则\(y\)轴截距\(b\)，\(x\)轴截距\(-4b/3\)，面积\((1/2)|b\cdot (-4b/3)|=2b^{2}/3=6\)，\(b^{2}=9\)，\(b=\pm 3\)，\(l\)：\(y=(3/4)x\pm 3\)；回代验证：\(b=3\) 代回，面积\((1/2)\cdot |3\cdot (-4)|=6\) ；\(b=-3\) 代回，面积\((1/2)\cdot |(-3)\cdot 4|=6\) ，双向验证成立； \((2) m\neq 1\)时两点横、纵坐标均不等，斜率\(k=(1-0)/(m-1)=1/(m-1)\)，点斜式\(y=(1/(m-1))(x-1)\)；\(m=1\)时\(A\)、\(B\)同在直线\(x=1\)上，\(l\)：\(x=1\)； \((3)\) 设截距\(a\)、\(b\)．①过原点（\(a=b=0\)）：\(l\)过\((0,0)\)、\((4,-3)\)，\(y=-(3/4)x\)； ②不过原点，\(x/a+y/b=1\)，\(|a|=|b|\neq 0\)且\(4/a-3/b=1\)：\(a=b\)时\(1/a=1\)、\(a=1\)（\(x+y-1=0\)）；\(a=-b\)时\(4/a+3/a=1\)、\(a=7\)（\(x-y-7=0\)）． 共\(3\)条． （按源卷核对注：源答案\((1)\)写作\("y=(4/3)x\pm 3"\)，与其题面\("\)斜率\(3/4"\)不符——本卷按题面取修正值 \(y=(3/4)x\pm 3\)；面积回代双向验证已并入\((1)\)详解正文 ，详见亲算账④\("\)源误更正\("\)。）}
\fi
\end{ansblock}

\begin{ansblock}[2章-拓-课时10-02]
% ans:2章-拓-课时10-02
\ansitem{02}{\((1)x^{2}/4+y^{2}/3=1\)（\(x\neq \pm 2\)）；\((2)\)点\(M\)恒在定直线\(x=-4/3\)上}
\ifansbookdetail
\ansnote{详解}{\((1)\)重心\(P\)分中线为\(2:1\)，故 \(|PB|+|PC|=(2/3)\times 6=4>|BC|=2\)，\(P\)的轨迹是以\(B\)、\(C\)为焦点的椭圆（\(A\)与\(BC\)共线时退化，除去长轴端点），\(a=2\)、\(c=1\)、\(b=\sqrt{3}\)：\(x^{2}/4+y^{2}/3=1\)（\(x\neq \pm 2\)）。 \((2)\)设\(PQ\)：\(x=my-3\)，\(P(x_{1},y_{1})\)、\(Q(x_{2},y_{2})\)。联立 \((3m^{2}+4)y^{2}-18my+15=0\)，\(y_{1}+y_{2}=18m/(3m^{2}+4)\)、\(y_{1}y_{2}=15/(3m^{2}+4)\)，得 \(2my_{1}y_{2}=(5/3)(y_{1}+y_{2})\)。直线\(EP\)：\(y=y_{1}/(my_{1}-1)\cdot (x+2)\)；直线\(FQ\)：\(y=y_{2}/(my_{2}-5)\cdot (x-2)\)。联立解得 \(x=2(2my_{1}y_{2}-y_{2}-5y_{1})/(-y_{2}+5y_{1})\)，代入 \(2my_{1}y_{2}=(5/3)(y_{1}+y_{2})\)：\(x=2[(2/3)y_{2}-(10/3)y_{1}]/(-y_{2}+5y_{1})=(4/3)(y_{2}-5y_{1})/(-y_{2}+5y_{1})=-4/3\) 恒成立。故\(M\)恒在定直线 \(x=-4/3\) 上。}
\fi
\end{ansblock}

\begin{ansblock}[2章-拓-课时10-03]
% ans:2章-拓-课时10-03
\ansitem{03}{\((x-1/2)^{2}+y^{2}=1/4\)}
\ifansbookdetail
\ansnote{详解}{\(F(1,0)\)。\(OM\perp AB\) 且垂足\(M\)在过\(F\)的动直线\(AB\)上 \(\Rightarrow  \angle OMF=90^{\circ}\)（\(M=F\) 时直线\(AB\)为\(x=1\)亦合），故\(M\)在以\(OF\)为直径的圆上：圆心\((1/2,0)\)、半径\(1/2\)，方程 \((x-1/2)^{2}+y^{2}=1/4\)。}
\fi
\end{ansblock}

\begin{ansblock}[2章-拓-课时10-04]
% ans:2章-拓-课时10-04
\ansitem{04}{\(x^{2}+y^{2}=(2/3)b^{2}\)〔图注：题面示意照录自源〕}
\ifansbookdetail
\ansnote{详解}{设弦\(Q_{1}Q_{2}\)：\(mx+ny=1\)（不过原点可如此设）。齐次化：\(x^{2}/(2b^{2})+y^{2}/b^{2}-(mx+ny)^{2}=0\)，展开乘\(2b^{2}\)：\((2-2b^{2}n^{2})y^{2}+(1-2b^{2}m^{2})x^{2}-4mnb^{2}xy=0\)，两边除\(x^{2}\)得关于\(k=y/x\)的方程 \((2-2b^{2}n^{2})k^{2}-4mnb^{2}k+1-2b^{2}m^{2}=0\)，两根即\(OQ_{1}\)、\(OQ_{2}\)斜率，垂直 \(\Rightarrow  k_{1}k_{2}=(1-2b^{2}m^{2})/(2-2b^{2}n^{2})=-1 \Rightarrow  1-2b^{2}m^{2}=-2+2b^{2}n^{2} \Rightarrow  2b^{2}(m^{2}+n^{2})=3\)（\(\ast \)）。设\(D(x_{0},y_{0})\)，\(OD\perp Q_{1}Q_{2}\) 且\(D\)在弦上：直线\(Q_{1}Q_{2}\)即 \(x_{0}x+y_{0}y=x_{0}^{2}+y_{0}^{2}\)（过\(D\)垂直于\(OD\)），对照 \(mx+ny=1\) 得 \(m=x_{0}/(x_{0}^{2}+y_{0}^{2})\)、\(n=y_{0}/(x_{0}^{2}+y_{0}^{2})\)，代入（\(\ast \)）：\(2b^{2}(x_{0}^{2}+y_{0}^{2})/(x_{0}^{2}+y_{0}^{2})^{2}=3 \Rightarrow  x_{0}^{2}+y_{0}^{2}=(2/3)b^{2}\)。\(D\)轨迹为 \(x^{2}+y^{2}=(2/3)b^{2}\)。验算：\(a^{2}=2b^{2}\)、\(c^{2}=b^{2}\)，原点到\("\)斜距\("\)型结论 \(|OD|^{2}=a^{2}b^{2}/(a^{2}+b^{2})=2b^{4}/(3b^{2})=(2/3)b^{2}\) （椭圆切线对偶结构，此处仅作后台校核不作讲解）。}
\fi
\end{ansblock}

\begin{ansblock}[2章-拓-课时10-05]
% ans:2章-拓-课时10-05
\ansitem{05}{\(B\)}
\ifansbookdetail
\ansnote{详解}{设过\((0,1)\)的直线 \(y=kx+1\)（竖直直线\(x=0\)时\(A(0,2)\)、\(B(0,-2)\)，\(OP=OA+OB=(0,0)\)，\(P=(0,0)\)适合下文方程），代入圆：\((1+k^{2})x^{2}+2kx-3=0\)，\(x_{1}+x_{2}=-2k/(1+k^{2})\)，\(y_{1}+y_{2}=k(x_{1}+x_{2})+2=2/(1+k^{2})\)。由向量\(OA+\)向量\(OB=\)向量\(OP\) 得 \(P(x,y)=(x_{1}+x_{2}, y_{1}+y_{2})=(-2k/(1+k^{2}), 2/(1+k^{2}))\)。消\(k\)：\(y-1=(2-(1+k^{2}))/(1+k^{2})=(1-k^{2})/(1+k^{2})\)，于是 \(x^{2}+(y-1)^{2}=[4k^{2}+(1-k^{2})^{2}]/(1+k^{2})^{2}=(1+2k^{2}+k^{4})/(1+k^{2})^{2}=(1+k^{2})^{2}/(1+k^{2})^{2}=1\)，即 \(x^{2}+(y-1)^{2}=1\)。几何法（源【编注】口径）：设\(AB\)中点\(D\)，\(OP=2OD\)（向量），且\(OD\perp AB\)、\(D\)在以\(O(0,0)\)与\((0,1)\)为直径端点的圆 \(x^{2}+(y-1/2)^{2}=1/4\) 上，以\(O\)为中心\(2\)倍伸位即得 \(x^{2}+(y-1)^{2}=1\)。选\(B\)。}
\fi
\end{ansblock}

\begin{ansblock}[2章-拓-课时10-06]
% ans:2章-拓-课时10-06
\ansitem{06}{\(x^{2}/(1/2)+y^{2}=1\)（\(x\neq 0\)），即 \(2x^{2}+y^{2}=1\)（\(x\neq 0\)）}
\ifansbookdetail
\ansnote{详解}{\(l_{1}\)过定点\(A(0,1)\)、斜率\(k_{1}=k\)；\(l_{2}\)过定点\(B(0,-1)\)、斜率\(k_{2}=-2/k\)，\(k_{1}\cdot k_{2}=-2\)为定值。设\(P(x,y)\)，两斜率存在（\(k\neq 0\)）故 \(x\neq 0\)，\(k_{1}=(y-1)/x\)、\(k_{2}=(y+1)/x\)：\((y-1)(y+1)/x^{2}=-2 \Rightarrow  y^{2}-1=-2x^{2} \Rightarrow  2x^{2}+y^{2}=1\)（\(x\neq 0\)）。}
\fi
\end{ansblock}

\begin{ansblock}[2章-拓-课时10-07]
% ans:2章-拓-课时10-07
\ansitem{07}{\((y+1)^{2}/8-(x+1)^{2}/8=1\)}
\ifansbookdetail
\ansnote{详解}{设\(A(a,a)\)、\(B(b,b)\)（\(b>a\)）、\(M(x,y)\)。\(P\)、\(A\)、\(M\)共线：\((x+2)/(y-2)=(a+2)/(a-1) \Rightarrow  a=2(x+y)/(x-y+4)\)（\(y\neq 2\)）；\(Q\)、\(B\)、\(M\)共线：\(x/(y-2)=b/(b-2) \Rightarrow  b=2x/(x-y+2)\)。\(|AB|=\sqrt{2}(b-a)=\sqrt{2} \Rightarrow  b=a+1\)。代入整理：\(2x/(x-y+2)=2(x+y)/(x-y+4)+1 \Rightarrow  x^{2}-y^{2}+2x-2y+8=0\)。\(y=2\)时交点\(M\)与\(P\)或\(Q\)重合，两点坐标均满足该方程，故方程即轨迹方程；配方 \((y+1)^{2}/8-(x+1)^{2}/8=1\)，轨迹为中心\((-1,-1)\)、焦点在直线\(x=-1\)上的双曲线。}
\fi
\end{ansblock}

\begin{ansblock}[2章-拓-课时10-08]
% ans:2章-拓-课时10-08
\ansitem{08}{\((1)x^{2}-y^{2}=4\)（\(x\neq \pm 2\)）；\((2)\)存在，方向向量\((1,0)\)}
\ifansbookdetail
\ansnote{详解}{\((1)A(-2,0)\)、\(B(2,0)\)。设\(M(m,n)\)、\(N(m,-n)\)（\(m\neq 0\)，非直径），\(m^{2}+n^{2}=4\)。\(AM\)：\(y=n/(m+2)\cdot (x+2)\)；\(BN\)：\(y=-n/(m-2)\cdot (x-2)\)。两式相乘：\(y^{2}=-n^{2}/(m^{2}-4)\cdot (x^{2}-4)=(-n^{2}/-n^{2})(x^{2}-4)=x^{2}-4\)，即 \(x^{2}-y^{2}=4\)；\(m\neq 0\)排除\(M\)、\(N\)在\(x\)轴上（\(x\neq \pm 2\)）。 \((2)\)①\(l\)倾斜角\(90^{\circ}\)：设\(x=t\)（\(|t|>2\)），\(C\)、\(D=(t,\pm \sqrt{t^{2}-4})\)，\(AC\cdot AD=(t+2)^{2}-(t^{2}-4)=8t+8=0 \rightarrow  t=-1\)与\(|t|>2\)矛盾，无解。②\(l\)不竖直：设\(y=kx+t\)，联立 \((1-k^{2})x^{2}-2ktx-t^{2}-4=0\)（\(1-k^{2}\neq 0\)，\(\Delta >0\)），\(x_{1}+x_{2}=2kt/(1-k^{2})\)、\(x_{1}x_{2}=(-t^{2}-4)/(1-k^{2})\)。\(AC\perp AD\)：\((x_{1}+2)(x_{2}+2)+(kx_{1}+t)(kx_{2}+t)=0 \Rightarrow  (1+k^{2})x_{1}x_{2}+(2+kt)(x_{1}+x_{2})+t^{2}+4=0\)，代入化简得 \(k(t-2k)=0\)：\(k=0\)（恒合，验证\(\Delta =4(t^{2}+4)>0\)且\(1-0\neq 0\) ）或 \(t=2k\)（\(l\)：\(y=k(x+2)\)恒过\(A(-2,0)\)，舍）。综上存在，\(k=0\)，一个方向向量为\((1,0)\)。}
\fi
\end{ansblock}

\begin{ansblock}[2章-拓-课时10-09]
% ans:2章-拓-课时10-09
\ansitem{09}{\((1)(x-1)^{2}+y^{2}=1\)；\((2)(x-1/2)^{2}+(y-1/2)^{2}=3/2\)}
\ifansbookdetail
\ansnote{详解}{\((1)\)设\(P(x_{0},y_{0})\)（\(x_{0}^{2}+y_{0}^{2}=4\)）、\(AP\)中点\(M(x,y)\)：\(x=(2+x_{0})/2\)、\(y=y_{0}/2 \Rightarrow  x_{0}=2x-2\)、\(y_{0}=2y\)，代入：\((2x-2)^{2}+(2y)^{2}=4\)，即 \((x-1)^{2}+y^{2}=1\)。 \((2)\)设\(PQ\)中点\(N(x,y)\)。\(\angle PBQ=90^{\circ} \Rightarrow  Rt\triangle PBQ\)斜边中线 \(|BN|=|PN|=|PQ|/2\)；又\(ON\perp PQ\)（垂径定理），\(|OP|^{2}=|ON|^{2}+|PN|^{2} \Rightarrow  |PN|^{2}=4-(x^{2}+y^{2})\)。于是 \(|BN|^{2}=|PN|^{2}\)：\((x-1)^{2}+(y-1)^{2}=4-x^{2}-y^{2}\)，展开整理 \(2x^{2}+2y^{2}-2x-2y-2=0\)，即 \((x-1/2)^{2}+(y-1/2)^{2}=3/2\)。}
\fi
\end{ansblock}

\begin{ansblock}[2章-拓-课时10-10]
% ans:2章-拓-课时10-10
\ansitem{10}{\(D\)}
\ifansbookdetail
\ansnote{详解}{双曲线 \(c=\sqrt{4+12}=4\)，圆心\(D(-4,0)\)即左焦点，右焦点\(F_{2}(4,0)\)；圆半径\(r=1\)。\(|PN|\leqslant |PD|+r\)（\(N\)取\(PD\)延长线与圆交点时等号），故 \(-|PN|\geqslant -|PD|-1\)；又 \(|PM|\geqslant |PF_{2}|-|MF_{2}|=|PF_{2}|-\sqrt{13}\)（\(P\)、\(M\)、\(F_{2}\)共线时等号）。两式相加：\(|PM|-|PN|\geqslant |PF_{2}|-|PD|-\sqrt{13}-1\)。\(P\)在左支：\(|PF_{2}|-|PD|=2a=4\)（双曲线定义），得 \(|PM|-|PN|\geqslant 4-\sqrt{13}-1=3-\sqrt{13}\)。选\(D\)。}
\fi
\end{ansblock}

\begin{ansblock}[2章-拓-课时10-11]
% ans:2章-拓-课时10-11
\ansitem{11}{\(B\)}
\ifansbookdetail
\ansnote{详解}{\(l_{1}\perp l_{2}\)（\(m\cdot 1+(-1)\cdot m=0\)），\(l_{1}\)恒过定点\(S(3,1)\)、\(l_{2}\)恒过定点\(T(1,3)\)（\(m=0\)与一般值均合），\(\angle SPT=90^{\circ} \Rightarrow  P\)在以\(ST\)为直径的圆上：圆心\(M(2,2)\)、半径 \(r_{2}=|ST|/2=\sqrt{2}\)，即 \((x-2)^{2}+(y-2)^{2}=2\)。设\(AB\)中点\(D\)：\(|CD|=\sqrt{2^{2}-1^{2}}=\sqrt{3} \Rightarrow  D\)在以\(C(-1,-1)\)为心、\(\sqrt{3}\)为半径的圆上。\(PA+PB=2PD\)，\(|PD|\geqslant |MC|-r_{2}-\sqrt{3}=\sqrt{18}-\sqrt{2}-\sqrt{3}=3\sqrt{2}-\sqrt{2}-\sqrt{3}=2\sqrt{2}-\sqrt{3}\)（\(P\)、\(D\)在连心线上时等号可达），故 \(|PA+PB|_{2=}2|PD|\geqslant 4\sqrt{2}-2\sqrt{3}\)。选\(B\)。}
\fi
\end{ansblock}

\begin{ansblock}[2章-拓-课时10-12]
% ans:2章-拓-课时10-12
\ansitem{12}{\((3\sqrt{2}-\sqrt{6})/2\)}
\ifansbookdetail
\ansnote{详解}{\(M(1,1)\)，\(P(2,2)\)，\(|MP|=\sqrt{2}\)。设\(MP\)中点\(Q(3/2,3/2)\)。\(AB\perp CD\) 于 \(P \Rightarrow  Rt\triangle APC\)中 \(N\)为斜边\(AC\)中点：\(|PN|=|AN|=|CN|\)。垂径：\(AN^{2}+MN^{2}=AM^{2}=4\)，故 \(PN^{2}+MN^{2}=4\)。向量：\(PN=NQ+QP\)、\(MN=NQ+QM\)，\(QP=-QM\) 且 \(|QM|=|QP|=\sqrt{2}/2\)，相加：\(PN^{2}+MN^{2}=2NQ^{2}+2QM^{2}=2NQ^{2}+1=4 \Rightarrow  NQ^{2}=3/2\)，\(N\)在以\(Q\)为心、\(\sqrt{3/2}=\sqrt{6}/2\)为半径的圆上。\(|ON|\geqslant |OQ|-\sqrt{6}/2=(3\sqrt{2})/2-\sqrt{6}/2=(3\sqrt{2}-\sqrt{6})/2\)。验算存在性：\(N=OQ\)连线与圆交点确对应某对垂直弦（\(|PN|^{2}+|MN|^{2}=4\) 与 \(PN=MN\) 型中点可反解），等号可达。答 \((3\sqrt{2}-\sqrt{6})/2\)。}
\fi
\end{ansblock}

\begin{ansblock}[2章-拓-课时10-13]
% ans:2章-拓-课时10-13
\ansitem{13}{\(k=1\)：\(x=a/2\)（线段\(OA\)的垂直平分线）；\(k\neq 1\)（\(k>0\)）：\((k^{2}-1)x^{2}+(k^{2}-1)y^{2}-2k^{2}ax+k^{2}a^{2}=0\)，即圆心\((k^{2}a/(k^{2}-1),0)\)、半径\(|ka/(k^{2}-1)|\)的圆（阿波罗尼斯圆）}
\ifansbookdetail
\ansnote{详解}{设\(M(x,y)\)为曲线上任意一点，依题意 \(|MO|/|MA|=k\)（\(k>0\)）：\(\sqrt{x^{2}+y^{2}}/\sqrt{}((x-a)^{2}+y^{2})=k\)，两边平方并展开：\(x^{2}+y^{2}=k^{2}[(x-a)^{2}+y^{2}]\)，整理得 \((k^{2}-1)x^{2}+(k^{2}-1)y^{2}-2k^{2}ax+k^{2}a^{2}=0\)（化简步步同解，两性保住了）。分类： ①\(k=1\)：方程化为 \(-2ax+a^{2}=0\)，即 \(x=a/2\)——过\(OA\)中点\((a/2,0)\)且垂直于\(x\)轴的直线，即线段\(OA\)的垂直平分线。 ②\(k\neq 1\)：两边除以 \(k^{2}-1\)：\(x^{2}+y^{2}-[2k^{2}a/(k^{2}-1)]x+k^{2}a^{2}/(k^{2}-1)=0\)。配方：\((x-k^{2}a/(k^{2}-1))^{2}+y^{2}=k^{4}a^{2}/(k^{2}-1)^{2}-k^{2}a^{2}/(k^{2}-1)=k^{2}a^{2}[k^{2}-(k^{2}-1)]/(k^{2}-1)^{2}=k^{2}a^{2}/(k^{2}-1)^{2}>0\) 恒成立，故恒表示圆，圆心\((k^{2}a/(k^{2}-1),0)\)、半径\(|ka/(k^{2}-1)|\)（无需再验 \(D^{2}+E^{2}-4F>0\)，配方已直接给出正的半径平方）。 验算（\(k=1/2\)、\(a=2\) 特值）：圆心\((-2/3,0)\)、半径\(4/3\)，圆与\(x\)轴交点 \((2/3,0)\)：\(|MO|/|MA|=(2/3)/(4/3)=1/2\)；\((-2,0)\)：\(2/4=1/2\)。}
\fi
\end{ansblock}

\begin{ansblock}[2章-拓-课时10-14]
% ans:2章-拓-课时10-14
\ansitem{14}{\(x^{2}+y^{2}-4x-4y-1=0\)（即圆\(C_{2}\)本身，\((x-2)^{2}+(y-2)^{2}=9\)）}
\ifansbookdetail
\ansnote{详解}{记 \(C_{1}\)：\(x^{2}+y^{2}-1=0\)，\(C_{2}\)：\(x^{2}+y^{2}-4x-4y-1=0\)。①两圆确有两个交点：\(C_{1}\)圆心\((0,0)\)半径\(1\)，\(C_{2}\)圆心\((2,2)\)半径\(3\)，圆心距\(2\sqrt{2}\)，满足 \(3-1<2\sqrt{2}<3+1\) 。②过两交点的一切圆可表为 \(C_{1}+\lambda C_{2}=0\)（\(\lambda \neq -1\)；\(\lambda =-1\) 时退化为公共弦所在直线 \(x+y=0\)）——注意该系不含圆\(C_{2}\)本身。③代点\((2,-1)\)：\(C_{1}\)值\(=4+1-1=4\)，\(C_{2}\)值\(=4+1-8+4-1=0\)，系内圆过\((2,-)\)须 \(4+\lambda \cdot 0=0\)，无解！④而点\((2,-1)\)恰在\(C_{2}\)上（上式\(=0\)），\(C_{2}\)本身即\("\)过两交点且过\((2,-1)"\)的圆；两相交点加线外一点定一圆，故唯一，所求圆即 \(C_{2}\)：\(x^{2}+y^{2}-4x-4y-1=0\)。}
\fi
\end{ansblock}

\begin{ansblock}[2章-拓-课时10-15]
% ans:2章-拓-课时10-15
\ansitem{15}{\(x^{2}-9y^{2}=1\)（\(y\neq 0\)）}
\ifansbookdetail
\ansnote{详解}{设\(P(x_{1},y_{1})\)、\(M(x,y)\)。双曲线中 \(a=3\)、\(b=1\)、\(c=\sqrt{10}\)，\(F_{1}(-\sqrt{10},0)\)、\(F_{2}(\sqrt{10},0)\)。\(P\)为三角形顶点，\(y_{1}\neq 0\)。重心坐标公式：\(x=(x_{1}-\sqrt{10}+\sqrt{10})/3=x_{1}/3\)，\(y=y_{1}/3\)，即 \(x_{1}=3x\)、\(y_{1}=3y\)。\(P\)在双曲线上：\((3x)^{2}/9-(3y)^{2}=1\)，即 \(x^{2}-9y^{2}=1\)；由 \(y_{1}\neq 0\) 得 \(y\neq 0\)。故重心\(M\)的轨迹方程为 \(x^{2}-9y^{2}=1\)（\(y\neq 0\)）。}
\fi
\end{ansblock}

\begin{ansblock}[2章-拓-课时11-01]
% ans:2章-拓-课时11-01
\ansitem{01}{\(B\)}
\ifansbookdetail
\ansnote{详解}{关系式即\("\)到两定点\((1,0)\)、\((-1,0)\)距离和为常数\(4"\)，\(4>|F_{1}F_{2}|=2\)，\(P\)轨迹为椭圆：焦点在\(x\)轴、\(c=1\)、\(2a=4\Rightarrow a=2\)、\(b^{2}=a^{2}-c^{2}=3\)。标准方程 \(x^{2}/4+y^{2}/3=1\)。选\(B\)。}
\fi
\end{ansblock}

\begin{ansblock}[2章-拓-课时11-03]
% ans:2章-拓-课时11-03
\ansitem{03}{\(D\)}
\ifansbookdetail
\ansnote{详解}{圆\(M\)圆心\(M(-1,0)\)半径\(1\)；圆\(M_{2}\)圆心\((1,0)\)半径\(3\)。\(|MM_{2}|=2=3-1\)，两定圆内切于点\(A(-2,0)\)（左端点）。设动圆半径\(r_{1}\)：外切 \(\Rightarrow  |M_{1}M|=r_{1}+1\)；内切（动圆含于大圆内） \(\Rightarrow  |M_{1}M_{2}|=3-r_{1}\)。相加：\(|M_{1}M|+|M_{1}M_{2}|=4>|MM_{2}|=2\)，\(M_{1}\)轨迹以\(M\)、\(M_{2}\)为焦点的椭圆：\(a=2\)、\(c=1\)、\(b^{2}=3\)，\(x^{2}/4+y^{2}/3=1\)。退化剔除：\(r_{1}=0\) 时 \(M_{1}\) 与切点\(A(-2,0)\)重合（\("\)动圆\("\)退化为点，圆\(M_{1}\)不存在），故去 \(x=-2\)（右端点\(x=2\)处 \(r_{1}=\)… 检验：\(M_{1}(2,0)\)：\(|M_{1}M|=3=r_{1}+1\rightarrow r_{1}=2\)；\(|M_{1}M_{2}|=1=3-2\)  合法保留）。选\(D\)。}
\fi
\end{ansblock}

\begin{ansblock}[2章-拓-课时11-04]
% ans:2章-拓-课时11-04
\ansitem{04}{\((1)33/5\)；\((2)17/5\)}
\ifansbookdetail
\ansnote{详解}{\(a=5\)、\(b=3\)、\(c=4\)，\(F_{1}(-4,0)\)、\(F_{2}(4,0)\)。\(P\)横坐标\(2\)：\(4/25+y^{2}/9=1 \Rightarrow  y^{2}=9\cdot 21/25=189/25\)。\((1)|PF_{1}|=\sqrt{}((2+4)^{2}+y^{2})=\sqrt{36+189/25}=\sqrt{1089/25}=33/5\)。\((2)\)定义：\(|PF_{2}|=2a-|PF_{1}|=10-33/5=17/5\)。（直算复核：\(|PF_{2}|=\sqrt{}((2-4)^{2}+189/25)=\sqrt{4+189/25}=\sqrt{289/25}=17/5\)  两路一致）}
\fi
\end{ansblock}

\begin{ansblock}[2章-拓-课时11-05]
% ans:2章-拓-课时11-05
\ansitem{05}{最小值\(\sqrt{14}/4\)，最大值\(11\)}
\ifansbookdetail
\ansnote{详解}{焦点在\(y\)轴，\(a=6\)、\(b=2\)。设\(P(x_{0},y_{0})\)：\(x_{0}^{2}/4+y_{0}^{2}/36=1 \Rightarrow  x_{0}^{2}=4(1-y_{0}^{2}/36)=4-y_{0}^{2}/9\)，\(y_{0}\in [-6,6]\)。\(|PA|^{2}=x_{0}^{2}+(y_{0}-5)^{2}=4-y_{0}^{2}/9+y_{0}^{2}-10y_{0}+25=(8/9)y_{0}^{2}-10y_{0}+29=(8/9)(y_{0}-45/8)^{2}+7/8\)。顶点\(y_{0}=45/8\in [-6,6]\)，最小值\(|PA|^{2}=7/8 \Rightarrow  |PA|min=\sqrt{7/8}=\sqrt{14}/4\)；端点比较：\(y_{0}=-6\) 时 \((8/9)\cdot 36+60+29=121 \Rightarrow  |PA|=11\)；\(y_{0}=6\) 时 \(32-60+29=1 \Rightarrow  |PA|=1\)。最大\(11\)、最小\(\sqrt{14}/4\)。}
\fi
\end{ansblock}

\begin{ansblock}[2章-拓-课时11-06]
% ans:2章-拓-课时11-06
\ansitem{06}{\(B\)}
\ifansbookdetail
\ansnote{详解}{圆心\(C(0,1)\)、半径\(1\)：\(|MN|\leqslant |MC|+1\)（\(N\)在\(MC\)延长线与圆交点时等号）。设\(M(x_{0},y_{0})\)，\(x_{0}^{2}=18-2y_{0}^{2}\)（\(y_{0}\in [-3,3]\)）。\(|MC|^{2}=x_{0}^{2}+(y_{0}-1)^{2}=18-2y_{0}^{2}+y_{0}^{2}-2y_{0}+1=-(y_{0}+1)^{2}+20\leqslant 20\)，\(y_{0}=-1\in [-3,3]\) 可取，\(|MC|max=2\sqrt{5}\)。\(|MN|max=2\sqrt{5}+1\)。选\(B\)。}
\fi
\end{ansblock}

\begin{ansblock}[2章-拓-课时11-07]
% ans:2章-拓-课时11-07
\ansitem{07}{\(AC\)}
\ifansbookdetail
\ansnote{详解}{\(a^{2}=5\)、\(b^{2}=3\)、\(c^{2}=2\)：\(a=\sqrt{5}\)、\(c=\sqrt{2}\)。\(A\)对\(B\)错：定义距离和\(=2a=2\sqrt{5}\)。\(C\)对\(D\)错：\(|PF_{1}|=2a-|PF_{2}|\)，要最大须\(|PF_{2}|\)最小\(=a-c\)（右顶点），故\(max=2a-(a-c)=a+c=\sqrt{5}+\sqrt{2}\)。（\(|PF_{2}|\in [a-c,a+c]\)由拓\(4\)型距离直算可得，未引\(2.5.2\)性质表。）}
\fi
\end{ansblock}

\begin{ansblock}[2章-拓-课时11-08]
% ans:2章-拓-课时11-08
\ansitem{08}{\(10+\sqrt{13}\)}
\ifansbookdetail
\ansnote{详解}{标准形\(x^{2}/25+y^{2}/9=1\)：\(a=5\)、\(c=4\)，\(F_{1}(-4,0)\)、右焦点\(F_{2}(4,0)\)。定义转化：\(|PF_{1}|=10-|PF_{2}|\)，\(|PA|+|PF_{1}|=10+(|PA|-|PF_{2}|)\leqslant 10+|AF_{2}|=10+\sqrt{}((1-4)^{2}+2^{2})=10+\sqrt{13}\)（三角不等式，\(|PA|-|PF_{2}|\leqslant |AF_{2}|\) 当\(F_{2}\)在线段\(PA\)上取等）。等号可取性：\(A(1,2)\)在椭圆内（\(1/25+4/9=109/225<1\)），自\(A\)出发过\(F_{2}\)的射线必与椭圆交于一点\(P\)（\(F_{2}\)亦椭圆内点，交点在\(F_{2}\)外侧），此时 \(|PA|-|PF_{2}|=|AF_{2}|=\sqrt{13}\)。最大值\(10+\sqrt{13}\)。}
\fi
\end{ansblock}

\begin{ansblock}[2章-拓-课时11-09]
% ans:2章-拓-课时11-09
\ansitem{09}{\((3,4)\)或\((3,-4)\)或\((-3,4)\)或\((-3,-4)\)}
\ifansbookdetail
\ansnote{详解}{\(c^{2}=45-20=25\)：\(F_{1}(-5,0)\)、\(F_{2}(5,0)\)。设\(P(x,y)\)：垂直 \(\Rightarrow \) 向量\(PF_{1}\cdot \)向量\(PF_{2}=(-5-x,-y)\cdot (5-x,-y)=(x^{2}-25)+y^{2}=0\)，即\(P\)在圆\(x^{2}+y^{2}=25\)上；联立 \(x^{2}/45+y^{2}/20=1\)：代\(x^{2}=25-y^{2}\)：\((25-y^{2})/45+y^{2}/20=1\)，\(\times 180\)：\(100-4y^{2}+9y^{2}=180 \Rightarrow  y^{2}=16\)、\(x^{2}=9\)。\(P=(\pm 3,\pm 4)\)四解。验算任一组：\(9/45+16/20=1/5+4/5=1\)；\((-5-3)(5-3)+(-4)(-4)?\) 取\((3,4)\)：\((x^{2}-25)+y^{2}=9-25+16=0\)。}
\fi
\end{ansblock}

\begin{ansblock}[2章-拓-课时11-10]
% ans:2章-拓-课时11-10
\ansitem{10}{\((1)x^{2}/10+y^{2}/5=1\)；\((2)y^{2}/(1/4)+x^{2}/(1/5)=1\)（即\(5x^{2}+4y^{2}=1\)）}
\ifansbookdetail
\ansnote{详解}{\((1)\)共焦：\(c^{2}=9-4=5\)，设 \(x^{2}/a^{2}+y^{2}/(a^{2}-5)=1\)，代\(Q\)：\(8/a^{2}+1/(a^{2}-5)=1\)，令\(u=a^{2}\)：\(8(u-5)+u=u(u-5) \Rightarrow  u^{2}-14u+40=0 \Rightarrow  u=10\)或\(4\)（\(u=4<5\)舍，\(b^{2}<0\)）：\(a^{2}=10\)、\(b^{2}=5\)，\(x^{2}/10+y^{2}/5=1\)。验算\(Q\)：\(8/10+1/5=1\) 。（源法二共焦系\(x^{2}/(9+\lambda )+y^{2}/(4+\lambda )=1\)代点\(\lambda =1\)同果。） \((2)\)焦点轴不定，设 \(mx^{2}+ny^{2}=1\)（\(m>0\)、\(n>0\)、\(m\neq n\)）。代\(P(1/3,1/3)\)：\((m+n)/9=1\)；代\(Q(0,-1/2)\)：\(n/4=1 \Rightarrow  n=4\)、\(m=5\)：\(5x^{2}+4y^{2}=1\)，标准形 \(y^{2}/(1/4)+x^{2}/(1/5)=1\)（焦点在\(y\)轴）。验算：\(P\)：\(5/9+4/9=1\)。}
\fi
\end{ansblock}

\begin{ansblock}[2章-拓-课时11-11]
% ans:2章-拓-课时11-11
\ansitem{11}{\(C\)}
\ifansbookdetail
\ansnote{详解}{由\(S\_\)圆\(=S\_\)环恒成立，半椭球体积\(=\)柱\(-\)锥\(=\pi b^{2}a-(1/3)\pi b^{2}a=(2/3)\pi b^{2}a\)。短轴长\(2 \Rightarrow  b=1\)；长轴\(4 \Rightarrow \) 半长轴\(a=2\)。\(V\)椭球\(=2\cdot (2/3)\pi \cdot 1^{2}\cdot 2=8\pi /3\)。选\(C\)。}
\fi
\end{ansblock}

\begin{ansblock}[2章-拓-课时11-12]
% ans:2章-拓-课时11-12
\ansitem{12}{\(B\)}
\ifansbookdetail
\ansnote{详解}{绕\(y\)轴旋转：截面为半径\(|x|\)的圆盘。构造底面半径\(2\)、高\(3\)的圆柱，挖去以下底圆心为顶点、以上底为底的圆锥。高\(h\)处（\(0\leqslant h\leqslant 3\)）：圆锥截面半径\(r=(2/3)h\)，环面积\(=4\pi -(4h^{2}\pi )/9\)；橄榄体同高处截面\(=\pi x^{2}\)，由 \(x^{2}/4+h^{2}/9=1\) 得 \(x^{2}=4(1-h^{2}/9)\)，截面\(=4\pi -(4h^{2}\pi )/9\)，两处恒等。半橄榄\(=\)柱\(-\)锥\(=\pi \cdot 2^{2}\cdot 3-(1/3)\pi \cdot 2^{2}\cdot 3=12\pi -4\pi =8\pi \)，\(V=2\times 8\pi =16\pi \)。选\(B\)。}
\fi
\end{ansblock}

\begin{ansblock}[2章-拓-课时11-13]
% ans:2章-拓-课时11-13
\ansitem{13}{\(4.75m\)}
\ifansbookdetail
\ansnote{详解}{接触点\(P\)处绳两直段和\(=30\)：\(|PA|+|PC|=30\)（\(A\)、\(C\)为两桅杆顶），\(P\)在以\(A\)、\(C\)为焦点的椭圆上：\(2a=30\)、\(a=15\)；两焦点距离\(=\)两桅杆顶间距\(=15\)（两顶等高\(7.5m\)、水平相距\(15m\)），\(2c=15\)、\(c=7.5\)；\(b^{2}=a^{2}-c^{2}=225-56.25=168.75=675/4\)。以\(AC\)中点为原点建系：\(x^{2}/225+y^{2}/(675/4)=1\)。甲板在杆顶下方\(7.5m\)：\(y\_P=-7.5\)，代得 \(x\_P^{2}=225(1-56.25/168.75)=225\cdot (2/3)=150\)，\(P\)偏\(AB\)侧取 \(x\_P=-\sqrt{150}\)。\(P\)到桅杆\(AB\)（所在竖直线\(x=-7.5\)）的距离\(=|x\_P|-7.5=\sqrt{150}-7.5\approx 12.25-7.5=4.75(m)\)。}
\fi
\end{ansblock}

\begin{ansblock}[2章-拓-课时11-14]
% ans:2章-拓-课时11-14
\ansitem{14}{\(C\)}
\ifansbookdetail
\ansnote{详解}{\(l_{1}\)恒过\((-1,0)\)、\(l_{2}\)恒过\((1,0)\)（\(m=0\)与\(m\neq 0\)均合），且\(l_{1}\perp l_{2}\)（系数\(m\cdot 1+(-1)\cdot m=0\)），交点\(Q\)满足\(\angle (-1,0)Q(1,0)=90^{\circ}\)或重合端点，即\(Q\)在圆\(x^{2}+y^{2}=1\)上（含\((\pm 1,0)\)退化界点——\(m\)致两线共线时无交，逐\(m\)检验交点总存在，轨迹即整圆）。椭圆\(c^{2}=3\)：\(F_{1}(-\sqrt{3},0)\)、\(F_{2}(\sqrt{3},0)\)。设\(Q(cos\theta ,sin\theta )\)：\(|QF_{1}|=\sqrt{}((cos\theta +\sqrt{3})^{2}+sin^{2}\theta )=\sqrt{4+2cos\theta }\)，\(|QF_{2}|=\sqrt{4-2cos\theta }\)，乘积\(=\sqrt{16-12cos^{2}\theta }\in [\sqrt{4},\sqrt{16}]=[2,4]\)（\(cos^{2}\theta \in [0,1]\)）。选\(C\)。}
\fi
\end{ansblock}

\begin{ansblock}[2章-拓-课时11-15]
% ans:2章-拓-课时11-15
\ansitem{15}{\(b^{2}\cdot tan(\alpha /2)\)}
\ifansbookdetail
\ansnote{详解}{\(|PF_{1}|+|PF_{2}|=2a\)、\(|F_{1}F_{2}|=2c\)。余弦定理：\(cos\alpha =(|PF_{1}|^{2}+|PF_{2}|^{2}-4c^{2})/(2|PF_{1}||PF_{2}|)=[(2a)^{2}-2|PF_{1}||PF_{2}|-4c^{2}]/(2|PF_{1}||PF_{2}|)=(4b^{2}-2|PF_{1}||PF_{2}|)/(2|PF_{1}||PF_{2}|)=(2b^{2})/(|PF_{1}||PF_{2}|)-1\)，解出 \(|PF_{1}||PF_{2}|=2b^{2}/(1+cos\alpha )\)。\(S=\)½\(|PF_{1}||PF_{2}|sin\alpha =b^{2}\cdot sin\alpha /(1+cos\alpha )=b^{2}\cdot [2sin(\alpha /2)cos(\alpha /2)]/[2cos^{2}(\alpha /2)]=b^{2}\cdot tan(\alpha /2)\)。}
\fi
\end{ansblock}

\begin{ansblock}[2章-拓-课时11-16]
% ans:2章-拓-课时11-16
\ansitem{16}{\(B\)}
\ifansbookdetail
\ansnote{详解}{内心到三边距离\(=\)内切圆半径\(r\)；\(I(x_{0},2)\)到\(x\)轴（边\(F_{1}F_{2}\)所在直线）距离为\(2\)，故 \(r=2\)。面积\(=\)½\(r\times \)周长：\(4b=\)½\(\cdot 2\cdot (|MF_{1}|+|MF_{2}|+|F_{1}F_{2}|)=2a+2c \Rightarrow  a+c=2b\)。平方：\((a+c)^{2}=4b^{2}=4(a^{2}-c^{2})=4(a-c)(a+c) \Rightarrow  a+c=4(a-c) \Rightarrow  5c=3a\)，所求\(=(2a)/(2c)=a/c=5/3\)。选\(B\)。验算：\(a/c=5/3\)时取\(a=5k\)、\(c=3k\)、\(b=4k\)：\(a+c=8k=2b=8k\)  面积式闭合。}
\fi
\end{ansblock}

\begin{ansblock}[2章-拓-课时11-17]
% ans:2章-拓-课时11-17
\ansitem{17}{\((1)3<m<7\)或\(7<m<11\)；\((2)7<m<11\)；\((3)3<m<7\)}
\ifansbookdetail
\ansnote{详解}{\((1)\)椭圆标准方程要求两分母同正且不相等（相等即圆）：\(m-3>0\)、\(11-m>0\)、\(m-3\neq 11-m \Rightarrow  3<m<11\) 且 \(m\neq 7\)。\((2)\)焦点在\(x\)轴\(\Leftrightarrow x^{2}\)项分母大：\(m-3>11-m>0 \Rightarrow  m>7\) 且 \(m<11\)。\((3)\)焦点在\(y\)轴\(\Leftrightarrow y^{2}\)项分母大：\(11-m>m-3>0 \Rightarrow  3<m<7\)。}
\fi
\end{ansblock}

\begin{ansblock}[2章-拓-课时12-01]
% ans:2章-拓-课时12-01
\ansitem{01}{\((1)(\sqrt{5},0)\)、\((-\sqrt{5},0)\)；\((2)(0,\sqrt{6})\)、\((0,-\sqrt{6})\)；\((3)(0,-3)\)、\((0,3)\)；\((4)(\sqrt{2}/2,0)\)、\((-\sqrt{2}/2,0)\)}
\ifansbookdetail
\ansnote{详解}{化标准式后比较分母定焦点所在轴，按 \(c^{2}=a^{2}-b^{2}\) 求 \(c\)。 \((1)a^{2}=9\)、\(b^{2}=4\)，\(c^{2}=5\)，焦点在 \(x\) 轴：\((\pm \sqrt{5},0)\)； \((2)a^{2}=9\)、\(b^{2}=3\)，\(c^{2}=6\)，焦点在 \(y\) 轴：\((0,\pm \sqrt{6})\)； \((3)\)化为 \(x^{2}/7+y^{2}/16=1\)，\(a^{2}=16\)、\(b^{2}=7\)，\(c^{2}=9\)，焦点在 \(y\) 轴：\((0,\pm 3)\)； \((4)\)化为 \(x^{2}+y^{2}/(1/2)=1\)，\(a^{2}=1\)、\(b^{2}=1/2\)，\(c^{2}=1/2\)，\(c=\sqrt{2}/2\)，焦点在 \(x\) 轴：\((\pm \sqrt{2}/2,0)\)．}
\fi
\end{ansblock}

\begin{ansblock}[2章-拓-课时12-02]
% ans:2章-拓-课时12-02
\ansitem{02}{\(\sqrt{5}-\sqrt{2}\)}
\ifansbookdetail
\ansnote{详解}{以 \(O\) 为原点、\(F_{1}F_{2}\) 为 \(x\) 轴，则 \(F_{1}(-c,0)\)、\(F_{2}(c,0)\)，等腰直角三角形顶点 \(P(0,c)\)（\(|OP|=\)½\(|F_{1}F_{2}|=c\)）． 斜边 \(F_{1}F_{2}\) 外的两腰 \(PF_{1}\)：\(y=x+c\)（\(x\in [-c,0]\)）、\(PF_{2}\)：\(x+y=c\)（\(x\in [0,c]\)）．由对称性，\(M\)、\(N\) 分别在两腰上且关于 \(y\) 轴对称，\(MN\) 水平，\(|MN|=(1/3)|F_{1}F_{2}|=2c/3\)，故 \(N\) 的横坐标为 \(c/3\)、\(M\) 的横坐标为 \(-c/3\)， 代入相应腰的方程得 \(N(c/3, 2c/3)\)、\(M(-c/3, 2c/3)\)． 由椭圆定义：\(2a=|NF_{1}|+|NF_{2}|=\sqrt{}((c/3+c)^{2}+(2c/3)^{2})+\sqrt{}((c/3-c)^{2}+(2c/3)^{2})=\sqrt{16c^{2}/9+4c^{2}/9}+\sqrt{4c^{2}/9+4c^{2}/9}=(2\sqrt{5}/3)c+(2\sqrt{2}/3)c\)， 故 \(e=c/a=3/(\sqrt{5}+\sqrt{2})=3(\sqrt{5}-\sqrt{2})/(5-2)=\sqrt{5}-\sqrt{2}\)． （数值核验：\(\sqrt{5}-\sqrt{2}\approx 2.236-1.414=0.822\in (0,1)\) ）}
\fi
\end{ansblock}

\begin{ansblock}[2章-拓-课时12-03]
% ans:2章-拓-课时12-03
\ansitem{03}{\(1/2\) 或 \(3/2\)}
\ifansbookdetail
\ansnote{详解}{设 \(|PF_{1}|=4t\)、\(|F_{1}F_{2}|=2c=3t\)、\(|PF_{2}|=2t\)（\(t>0\)），则 \(c=3t/2\)。 椭圆支：\(2a=|PF_{1}|+|PF_{2}|=6t \Rightarrow  a=3t\)，\(e=c/a=(3t/2)/(3t)=1/2\) （且 \(2a>2c\)：\(6t>3t\) ，可成椭圆）； 双曲线支：\(2a=||PF_{1}|-|PF_{2}||=2t \Rightarrow  a=t\)，\(e=c/a=(3t/2)/t=3/2\) （\(2a<2c\)：\(2t<3t\) ，可成双曲线）．}
\fi
\end{ansblock}

\begin{ansblock}[2章-拓-课时12-04]
% ans:2章-拓-课时12-04
\ansitem{04}{\((1)\)证明见详解；\((2)x^{2}/144+y^{2}/25=1\)}
\ifansbookdetail
\ansnote{详解}{\((1)A_{1}(-a,0)\)、\(A_{2}(a,0)\)、\(B_{1}(0,-b)\)、\(B_{2}(0,b)\)，两对角线 \(A_{1}A_{2}\) 与 \(B_{1}B_{2}\) 同为对方的中垂线（互相垂直平分），故 \(A_{1}B_{1}A_{2}B_{2}\) 是平行四边形且对角线垂直，即为菱形． \((2)\)菱形面积\(=\)½\(\cdot |A_{1}A_{2}|\cdot |B_{1}B_{2}|=\)½\(\cdot 2a\cdot 2b=2ab=120 \Rightarrow  ab=60\)；相邻顶点 \(A_{2}(a,0)\)、\(B_{2}(0,b)\) 的距离即边长 \(\sqrt{a^{2}+b^{2}}=13 \Rightarrow  a^{2}+b^{2}=169\)． 联立 \(a^{2}+b^{2}=169\) 与 \(ab=60\)：\(a^{2}\)、\(b^{2}\) 是 \(u^{2}-169u+3600=0\) 的两根，\(\Delta =169^{2}-4\cdot 3600=28561-14400=14161=119^{2}\)， \(u=(169\pm 119)/2=144\) 或 \(25\)，由 \(a>b\) 得 \(a^{2}=144\)、\(b^{2}=25\)，椭圆 \(x^{2}/144+y^{2}/25=1\)．}
\fi
\end{ansblock}

\begin{ansblock}[2章-拓-课时12-05]
% ans:2章-拓-课时12-05
\ansitem{05}{\(8\)}
\ifansbookdetail
\ansnote{详解}{\(P\)、\(Q\) 关于 \(O\) 对称，\(F_{1}\)、\(F_{2}\) 关于 \(O\) 对称，故 \(PF_{1}QF_{2}\) 的对角线 \(PQ\) 与 \(F_{1}F_{2}\) 互相平分（是平行四边形）；又 \(|PQ|=|F_{1}F_{2}|\)，对角线相等的平行四边形为矩形，于是 \(\angle F_{1}PF_{2}=90^{\circ}\)． \(a^{2}=16\)、\(b^{2}=4 \Rightarrow  c^{2}=12\)，\(|F_{1}F_{2}|^{2}=4c^{2}=48\)．设 \(|PF_{1}|=m\)、\(|PF_{2}|=n\)，则 \(m+n=2a=8\)，\(m^{2}+n^{2}=48\)， \(\Rightarrow  2mn=(m+n)^{2}-(m^{2}+n^{2})=64-48=16 \Rightarrow  mn=8\)． 矩形面积\(=mn=8\)．}
\fi
\end{ansblock}

\begin{ansblock}[2章-拓-课时12-06]
% ans:2章-拓-课时12-06
\ansitem{06}{\(D\)}
\ifansbookdetail
\ansnote{详解}{\(a^{2}=25\)、\(b^{2}=25/2\)，\(c^{2}=a^{2}-b^{2}=25/2\)，\(e^{2}=c^{2}/a^{2}=1/2\)，故 \(b=c=5\sqrt{2}/2\)． 设 \(P(x,y)\)，\(x^{2}\in [0,25]\)．由焦半径 \(|PF_{1}|=a+ex\)、\(|PF_{2}|=a-ex\) 得 \(|PF_{1}|^{2}+|PF_{2}|^{2}=(a+ex)^{2}+(a-ex)^{2}=2a^{2}+2e^{2}x^{2}=50+x^{2}\)，\(|PF_{1}|\cdot |PF_{2}|=a^{2}-e^{2}x^{2}=25-x^{2}/2\)，且 \(|F_{1}F_{2}|^{2}=4c^{2}=50\)， 故 \(cos\angle F_{1}PF_{2}=(|PF_{1}|^{2}+|PF_{2}|^{2}-|F_{1}F_{2}|^{2})/(2|PF_{1}|\cdot |PF_{2}|)=(50+x^{2}-50)/[2(25-x^{2}/2)]=x^{2}/(50-x^{2})\geqslant 0\)， 分母 \(50-x^{2}\geqslant 25>0\) 恒正，等号当且仅当 \(x=0\)，即 \(P\) 为短轴端点 \((0,\pm 5\sqrt{2}/2)\)． \(cos\angle F_{1}PF_{2}\geqslant 0\) 说明 \(\angle F_{1}PF_{2}\leqslant 90^{\circ}\)，且 \(cos\) 在 \([0,\pi ]\) 上随角单调递减，故 \(cos\) 取最小值 \(0\) 时顶角最大： \(x=0\) 时 \(|PF_{1}|=|PF_{2}|=\sqrt{b^{2}+c^{2}}=\sqrt{25/2+25/2}=5=a\)，\(|F_{1}F_{2}|=2c=5\sqrt{2}\)，\(5^{2}+5^{2}=50=(5\sqrt{2})^{2} \Rightarrow  \angle F_{1}PF_{2}=90^{\circ}\)； \(x^{2}=25\)（长轴端点，三点共线）时 \(cos=25/(50-25)=1\)，顶角 \(0^{\circ}\)，与图形一致． 所以 \(\angle F_{1}PF_{2}\) 的最大值为 \(90^{\circ}\)．故选：\(D\)．}
\fi
\end{ansblock}

\begin{ansblock}[2章-拓-课时12-07]
% ans:2章-拓-课时12-07
\ansitem{07}{\([-2\sqrt{6}/3, 2\sqrt{6}/3]\)}
\ifansbookdetail
\ansnote{详解}{\(c^{2}=4-1=3\)，\(F_{1}(-\sqrt{3},0)\)、\(F_{2}(\sqrt{3},0)\)．\(P\) 不在 \(x\) 轴上的线段 \(F_{1}F_{2}\) 内（椭圆上点 \(|y|\) 可取\(0\)但此时 \(x=\pm 2\)，在 \(F_{1}F_{2}\) 之外），故 \(\angle F_{1}PF_{2}\geqslant 90^{\circ} \Leftrightarrow \) 向量 \(PF_{1}\cdot PF_{2}\leqslant 0 \Leftrightarrow  (-\sqrt{3}-x)(\sqrt{3}-x)+y^{2}\leqslant 0 \Leftrightarrow  x^{2}+y^{2}-3\leqslant 0\)，即 \(P\) 在以原点为圆心、\(\sqrt{3}\) 为半径的圆内（含界）． 以 \(y^{2}=1-x^{2}/4\) 代入：\(x^{2}+1-x^{2}/4-3\leqslant 0 \Rightarrow  (3/4)x^{2}\leqslant 2 \Rightarrow  x^{2}\leqslant 8/3 \Rightarrow  |x|\leqslant 2\sqrt{6}/3\)， 又 \(2\sqrt{6}/3\approx 1.633<2\)  在椭圆 \(x\) 范围内，故 \(x\in [-2\sqrt{6}/3, 2\sqrt{6}/3]\)．}
\fi
\end{ansblock}

\begin{ansblock}[2章-拓-课时12-08]
% ans:2章-拓-课时12-08
\ansitem{08}{\(A\)}
\ifansbookdetail
\ansnote{详解}{\(a^{2}=18\)、\(b^{2}=9\)，短轴端点 \(A_{1}(0,3)\)、\(A_{2}(0,-3)\)．设 \(M(x_{0},y_{0})\)（\(x_{0}\neq 0\)），\(N(x_{1},y_{1})\)． \(k_{MA_{1}}=(y_{0}-3)/x_{0}\)，\(NA_{1}\perp MA_{1} \Rightarrow  k_{NA_{1}}=-x_{0}/(y_{0}-3)\)，\(NA_{1}\)：\(y=-x_{0}/(y_{0}-3)\cdot x+3\)； 同理 \(NA_{2}\)：\(y=-x_{0}/(y_{0}+3)\cdot x-3\)． 两式相减消 \(y\)：\(0=-x_{0}x[1/(y_{0}-3)-1/(y_{0}+3)]+6 \Rightarrow  x_{0}x\cdot 6/(y_{0}^{2}-9)=6 \Rightarrow  x_{1}=(y_{0}^{2}-9)/x_{0}\)． 由 \(M\) 在椭圆上：\(x_{0}^{2}/18+y_{0}^{2}/9=1 \Rightarrow  y_{0}^{2}-9=-x_{0}^{2}/2\)，故 \(x_{1}=-x_{0}/2\)． 两三角形共底 \(A_{1}A_{2}\)（在 \(y\) 轴上），面积比等于高的比\(=|x_{0}|/|x_{1}|=2\)．故选：\(A\)．}
\fi
\end{ansblock}

\begin{ansblock}[2章-拓-课时12-09]
% ans:2章-拓-课时12-09
\ansitem{09}{\(B\)}
\ifansbookdetail
\ansnote{详解}{向量\(PM\cdot \)向量\(AM=0 \Rightarrow  PM\perp AM\)，故 \(\triangle PMA\) 是以 \(M\) 为直角顶点的直角三角形， \(|PM|^{2}=|PA|^{2}-|AM|^{2}=|PA|^{2}-1\)，即 \(|PM|=\sqrt{|PA|^{2}-1}\)（\(|PA|>1\) 时 \(M\) 即以 \(A\) 为圆心、\(1\) 为半径的圆上自 \(P\) 所引切线的切点，必存在）， 所以 \(|PM|\) 随 \(|PA|\) 单调增大，只需求 \(|PA|\) 的最小值． 椭圆 \(a^{2}=25\)、\(b^{2}=16 \Rightarrow  c^{2}=a^{2}-b^{2}=9\)、\(c=3\)，右焦点为 \((3,0)\)，恰与点 \(A\) 重合，故 \(|PA|\) 即焦半径 \(|PF_{2}|\)， 其最小值为 \(a-c=5-3=2\)（\(P\) 为右顶点 \((5,0)\)）． 于是 \(|PM|\geqslant \sqrt{2^{2}-1}=\sqrt{3}\)，当 \(P\) 为右顶点时取等（此时 \(|PA|=2>1\)，切点 \(M\) 存在，\(|PM|=\sqrt{3}\)）．故选：\(B\)．}
\fi
\end{ansblock}

\begin{ansblock}[2章-拓-课时12-10]
% ans:2章-拓-课时12-10
\ansitem{10}{\([7,9)\)}
\ifansbookdetail
\ansnote{详解}{\(x=\sqrt{a-y^{2}} \Leftrightarrow  x^{2}+y^{2}=a\) 且 \(x\geqslant 0\)（须 \(a>0\)），表示以原点为圆心、\(\sqrt{a}\) 为半径的右半圆（含 \(x\) 轴上的两端点）． 椭圆 \(a_{1}^{2}=9\)、\(b_{1}^{2}=7\)，即 \(a_{1}=3\)、\(b_{1}=\sqrt{7}\)．椭圆上任一点 \(P(x,y)\) 到原点的距离平方： \(|OP|^{2}=x^{2}+y^{2}\)，由 \(x^{2}=9(1-y^{2}/7)\) 得 \(|OP|^{2}=9-(2/7)y^{2}\)，\(y^{2}\in [0,7]\)，故 \(|OP|^{2}\in [7,9]\)； 其中 \(|OP|^{2}=7\) 仅在上下顶点 \((0,\pm \sqrt{7})\)，\(|OP|^{2}=9\) 仅在左右顶点 \((\pm 3,0)\)，介于两者之间的每个值对应 \(x=\pm \)两个、\(y=\pm \)两个共四个点． 按 \(a=|OP|^{2}\) 逐段数交点（半圆只取 \(x\geqslant 0\)）： \(a<7\)：\(|OP|^{2}\) 在椭圆上最小为 \(7\)，半圆落在椭圆内部，无交点； \(a=7\)：交点 \((0,\sqrt{7})\)、\((0,-\sqrt{7})\)，恰 \(2\) 个 ； \(7<a<9\)：由 \(9-(2/7)y^{2}=a\) 得 \(y^{2}=7(9-a)/2\in (0,7)\)，两个 \(y\) 值各配 \(x=\sqrt{}(9(a-7)/2)>0\) 一个交点，共 \(2\) 个 ； \(a=9\)：仅 \(y=0\)、\(x=3\)，\(1\) 个交点（左顶点 \(x=-3\) 不在半圆上）； \(a>9\)：无交点． 综上 \(a\in [7,9)\)．}
\fi
\end{ansblock}

\begin{ansblock}[2章-拓-课时12-11]
% ans:2章-拓-课时12-11
\ansitem{11}{\(6\)}
\ifansbookdetail
\ansnote{详解}{\(a^{2}=4\)、\(b^{2}=3 \Rightarrow  c^{2}=1\)，故 \(F(-1,0)\)、\(O(0,0)\)．设 \(P(x,y)\)，\(x\in [-2,2]\)， 向量\(OP\cdot \)向量\(FP=(x,y)\cdot (x+1,y)=x^{2}+x+y^{2}\)． 以 \(y^{2}=3(1-x^{2}/4)=3-(3/4)x^{2}\) 代入：原式\(=x^{2}+x+3-(3/4)x^{2}=(1/4)x^{2}+x+3=(1/4)(x+2)^{2}+2\)． \(x\in [-2,2]\) 上该二次函数自对称轴 \(x=-2\) 起单调递增，故 \(x=2\)（\(P\) 为右顶点）时取最大值 \((1/4)\cdot 16+2=6\)．}
\fi
\end{ansblock}

\begin{ansblock}[2章-拓-课时12-12]
% ans:2章-拓-课时12-12
\ansitem{12}{\([-3/4, -3/8]\)}
\ifansbookdetail
\ansnote{详解}{\(a=2\)，\(A_{1}(-2,0)\)、\(A_{2}(2,0)\)．设 \(P(m,n)\)（\(m\neq \pm 2\)，两斜率均存在且非零），则 \(n^{2}=3(1-m^{2}/4)=3(4-m^{2})/4\)， \(k_{PA_{1}}\cdot k_{PA_{2}}=[n/(m+2)]\cdot [n/(m-2)]=n^{2}/(m^{2}-4)=[3(4-m^{2})/4]/(m^{2}-4)=-3/4\)， 即第三定义的斜率积定值 \(-b^{2}/a^{2}=-3/4\)． 由 \(1\leqslant k_{PA_{1}}\leqslant 2\) 且 \(k_{PA_{2}}=-3/(4k_{PA_{1}})\)；\(k_{PA_{1}}>0\) 时 \(k\mapsto -3/(4k)\) 单调递增， \(k_{PA_{1}}=1 \Rightarrow  k_{PA_{2}}=-3/4\)；\(k_{PA_{1}}=2 \Rightarrow  k_{PA_{2}}=-3/8\)， 故 \(k_{PA_{2}}\in [-3/4, -3/8]\)．}
\fi
\end{ansblock}

\begin{ansblock}[2章-拓-课时12-13]
% ans:2章-拓-课时12-13
\ansitem{13}{\(1/2+2\sqrt{21}/3\)；\((\pm 2\sqrt{5}/3, -2/3)\)}
\ifansbookdetail
\ansnote{详解}{椭圆化为 \(x^{2}/4+y^{2}=1\)，设 \(Q(2cos\alpha , sin\alpha )\)；圆心 \(C(0,2)\)、半径 \(1/2\)． 三角不等式 \(|PQ|\leqslant |PC|+|CQ|=1/2+|CQ|\)，等号当 \(P\) 为半径 \(CQ\) 延长线与圆的交点时成立，故只需求 \(|CQ|\) 的最大值． \(|CQ|^{2}=4cos^{2}\alpha +(sin\alpha -2)^{2}=4(1-sin^{2}\alpha )+sin^{2}\alpha -4sin\alpha +4=8-4sin\alpha -3sin^{2}\alpha =-3(sin\alpha +2/3)^{2}+28/3\leqslant 28/3\)， 等号当 \(sin\alpha =-2/3\)（可取）时成立，此时 \(cos\alpha =\pm \sqrt{1-4/9}=\pm \sqrt{5}/3\)， \(|CQ|max=\sqrt{28/3}=2\sqrt{7}/\sqrt{3}=2\sqrt{21}/3\)，相应 \(Q(\pm 2\sqrt{5}/3, -2/3)\)． 故 \(|PQ|max=1/2+2\sqrt{21}/3\)，相应点 \(Q\) 为 \((\pm 2\sqrt{5}/3, -2/3)\)．}
\fi
\end{ansblock}

\begin{ansblock}[2章-拓-课时12-14]
% ans:2章-拓-课时12-14
\ansitem{14}{\(B\)}
\ifansbookdetail
\ansnote{详解}{①向径即焦半径，\(|PF|=a-ex\)（设焦点为右焦点、\(x\in [-a,a]\)），故取值范围为 \([a-c, a+c]\)，①正确． ②比值为 \((a-c)/(a+c)=(1-e)/(1+e)\)，对 \(e\) 求导得 \([(1-e)/(1+e)]'=-2/(1+e)^{2}<0\)，故比值越大 \(\Leftrightarrow  e\) 越小 \(\Leftrightarrow \) 轨道越圆，②错误． ③以焦点 \(F(c,0)\) 为顶点看两半弧扫过的面积：左半弧（\(x\leqslant 0\) 段）扫过面积\(=\)半个椭圆面积\(+\triangle F B_{1}B_{2}\)（\(B_{1}\)、\(B_{2}\) 为短轴端点）\(=\pi ab/2+bc\)；右半弧扫过面积\(=\pi ab/2-bc\)．由面积守恒，扫过面积与用时成正比，左半弧用时长，③正确． ④近地点向径短、单位时间扫过同面积须弧长长，故速度最大；远地点速度最小，④所说相反，错误． 综上正确的是①③．故选：\(B\)．}
\fi
\end{ansblock}

\begin{ansblock}[2章-拓-课时12-15]
% ans:2章-拓-课时12-15
\ansitem{15}{\(5\sqrt{5}-5\)；\((\sqrt{5}+1)/2\)}
\ifansbookdetail
\ansnote{详解}{第一空：\(10>m>0\) 时焦点在 \(x\) 轴上，\(e^{2}=1-m/10\)．又 \(e=(\sqrt{5}-1)/2 \Rightarrow  e^{2}=(6-2\sqrt{5})/4=(3-\sqrt{5})/2\)， 故 \(m/10=1-(3-\sqrt{5})/2=(\sqrt{5}-1)/2\)，\(m=5(\sqrt{5}-1)=5\sqrt{5}-5\)． 第二空：设 \(\triangle PF_{1}F_{2}\) 的内切圆半径为 \(r\)，则 \(M\) 到三边距离均为 \(r\)． \(S\triangle PF_{1}F_{2}=(1/2)r(|PF_{1}|+|PF_{2}|+|F_{1}F_{2}|)=(1/2)r(2a+2c)=(a+c)r\)； \(S\triangle MF_{1}F_{2}=(1/2)r|F_{1}F_{2}|=cr\)． \(\triangle MF_{1}F_{2}\) 与 \(\triangle PF_{1}F_{2}\) 同底 \(F_{1}F_{2}\)，面积比等于高之比；又 \(P\)、\(M\)、\(N\) 共线且 \(N\) 在 \(F_{1}F_{2}\) 上，故高之比\(=|MN|/|PN|\)， 于是 \(|MN|/|PN|=cr/[(a+c)r]=c/(a+c)\)， \(|PM|=|PN|-|MN|=[1-c/(a+c)]|PN|=a/(a+c)\cdot |PN|\)， 所以 \(|PM|/|MN|=a/c=1/e=2/(\sqrt{5}-1)=(\sqrt{5}+1)/2\)．}
\fi
\end{ansblock}

\begin{ansblock}[2章-拓-课时12-16]
% ans:2章-拓-课时12-16
\ansitem{16}{\(AC\)}
\ifansbookdetail
\ansnote{详解}{\(A\)：圆写作 \(x^{2}/R^{2}+y^{2}/R^{2}=1\)，即 \(a^{2}=b^{2}=R^{2}\)，圆上点满足 \(x_{0}^{2}+y_{0}^{2}=R^{2}\)， \(R'=R^{4}((x_{0}^{2}+y_{0}^{2})/R^{4})^(3/2)=R^{4}(R^{2}/R^{4})^(3/2)=R^{4}/R^{3}=R\)，与点的位置无关，\(A\) 正确． \(B\)、\(C\)：记 \(f=x_{0}^{2}/a^{4}+y_{0}^{2}/b^{4}\)，令 \(u=x_{0}^{2}/a^{2}\in [0,1]\)，则 \(y_{0}^{2}/b^{2}=1-u\)，\(f=u/a^{2}+(1-u)/b^{2}=1/b^{2}+u(1/a^{2}-1/b^{2})\)， 因 \(a>b>0\) 有 \(1/a^{2}-1/b^{2}<0\)，故 \(f\) 关于 \(u\) 单调递减：\(u=0\)（短轴端点 \((0,\pm b)\)）时 \(f\) 最大\(=1/b^{2}\)，\(u=1\)（长轴端点 \((\pm a,0)\)）时 \(f\) 最小\(=1/a^{2}\)． 于是 \(R\_max=a^{2}b^{2}(1/b^{2})^(3/2)=a^{2}/b\)（在 \((0,\pm b)\) 处），\(R\_min=a^{2}b^{2}(1/a^{2})^(3/2)=b^{2}/a\)（在 \((\pm a,0)\) 处），\(R\in [b^{2}/a, a^{2}/b]\)． 故 \(B\)（最大值 \(a\)）错误，\(C\) 正确． \(D\)：\(b^{2}=1\)，\(R=a^{2}(1/(4a^{4})+y_{0}^{2})^(3/2)\)，其中 \(y_{0}^{2}=1-1/(4a^{2})\)．令 \(w=1/a^{2}\in (0,1)\)，则 \(1/(4a^{4})=w^{2}/4\)， \(R=(1/w)(1-w/4+w^{2}/4)^(3/2)\)，取对数 \(ln R=-ln w+(3/2)ln(1-w/4+w^{2}/4)\)，记 \(g(w)=1-w/4+w^{2}/4\)（\(g'=-1/4+w/2\)；\(g\) 的判别式 \(1/16-1<0\)，故 \(g>0\) 恒成立）， \(d(ln R)/dw=-1/w+(3/2)g'/g=[-g(w)+(3/2)w\cdot g'(w)]/[w\cdot g(w)]=(w^{2}/2-w/8-1)/[w\cdot g(w)]\)， \(w\in (0,1)\) 时 \(w^{2}/2<1/2\)，分子\(<-1/2-w/8<0\)，分母\(>0\)，故 \(d(ln R)/dw<0\)， 即 \(R\) 关于 \(w\) 单调递减，而 \(w=1/a^{2}\) 随 \(a\) 增大而减小，所以曲率半径 \(R\) 随 \(a\) 增大而增大，\(D\) 错误． 故选：\(AC\)．}
\fi
\end{ansblock}

\begin{ansblock}[2章-拓-课时12-17]
% ans:2章-拓-课时12-17
\ansitem{17}{\(BCD\)}
\ifansbookdetail
\ansnote{详解}{\(A\)：封闭图形由上半圆（半径 \(2\)）与下半椭圆（横半轴 \(2\)、竖半轴 \(3\)）拼成，横坐标范围 \(-2\leqslant x\leqslant 2\)．按 \(x\in \{-2,-1,0,1,2\}\) 逐列数严格内部整数点： \(x=\pm 2\) 时 \(y=0\) 在边界上，无内部整数点； \(x=\pm 1\) 时：上半需 \(y>0\) 且 \(1+y^{2}<4 \Rightarrow  y=1\)；下半需 \(y\leqslant 0\) 且 \(1/4+y^{2}/9<1 \Rightarrow  y^{2}<27/4 \Rightarrow  y=0,-1,-2\)（\(y=-3\) 时 \(1/4+1>1\) 在形外）； 即 \((-1,1)\)、\((-1,0)\)、\((-1,-1)\)、\((-1,-2)\) 与 \((1,1)\)、\((1,0)\)、\((1,-1)\)、\((1,-2)\)，共 \(8\) 个； \(x=0\) 时：上需 \(y>0\)、\(y^{2}<4 \Rightarrow  y=1\)；下需 \(y^{2}<9 \Rightarrow  y=0,-1,-2\)，共 \((0,1)\)、\((0,0)\)、\((0,-1)\)、\((0,-2)\)，\(4\) 个． 合计 \(8+4=12\) 个整数点，\(A\) 说 \(11\) 个，错误． \(B\)：\(y>0\) 段在半圆上，到原点距离恒为 \(2\)；\(y\leqslant 0\) 段设 \(P(2cos\theta , 3sin\theta )\)，\(\theta \in [\pi ,2\pi ]\)， \(|OP|^{2}=4cos^{2}\theta +9sin^{2}\theta =9-5cos^{2}\theta \in [4,9]\)（\(\theta =\pi ,2\pi \) 时取 \(4\)，\(\theta =3\pi /2\) 时取 \(9\)），即 \(|OP|\in [2,3]\)； 故曲线上到原点距离的最大值 \(3\)、最小值 \(2\)，和为 \(5\)，\(B\) 正确． \(C\)：椭圆 \(x^{2}/4+y^{2}/9=1\) 中 \(a^{2}=9\)（在 \(y\) 项）、\(b^{2}=4\)，\(c^{2}=5\)，焦点恰为 \(A(0,-\sqrt{5})\)、\(B(0,\sqrt{5})\)，故 \(P\) 在半椭圆上时 \(|PA|+|PB|=2a=6\)，\(|AB|=2\sqrt{5}\)． \(|PA|\cdot |PB|\leqslant [(|PA|+|PB|)/2]^{2}=9\)，等号当 \(|PA|=|PB|\)，即 \(P\) 在 \(x\) 轴上（\(P(\pm 2,0)\)，属 \(y\leqslant 0\) 段）取到． \(cos\angle APB=(|PA|^{2}+|PB|^{2}-|AB|^{2})/(2|PA||PB|)=[(|PA|+|PB|)^{2}-2|PA||PB|-20]/(2|PA||PB|)=16/(2|PA||PB|)-1=8/(|PA||PB|)-1\geqslant 8/9-1=-1/9\)， 故 \(cos\angle APB\) 的最小值为 \(-1/9\)，\(C\) 正确． \(D\)：由蒙日圆半径公式 \(r^{2}=a^{2}+b^{2}=9+4=13\)（亦如源取两条特殊切线 \(x=2\)、\(y=3\)，交点 \((2,3)\) 在蒙日圆上，\(r^{2}=4+9=13\)），蒙日圆为 \(x^{2}+y^{2}=13\)，\(D\) 正确． 故选：\(BCD\)．}
\fi
\end{ansblock}

\begin{ansblock}[2章-拓-课时12-18]
% ans:2章-拓-课时12-18
\ansitem{18}{\((1) x_{0}x/a^{2}+y_{0}y/b^{2}=1\)；\((2) x_{0}x/a^{2}+y_{0}y/b^{2}=1\)；\((3)\) 定值为 \(-b^{2}/a^{2}\)，证明见详解}
\ifansbookdetail
\ansnote{详解}{\((1)\) 类比①：圆方程 \(x^{2}/R^{2}+y^{2}/R^{2}=1\) 中 \(r^{2}\) 同时充当 \(a^{2}\)、\(b^{2}\)，切线 \(x_{0}x+y_{0}y=r^{2}\) 即 \(x_{0}x/R^{2}+y_{0}y/R^{2}=1\)，故把 \(r^{2}\) 对应替换为 \(a^{2}\)、\(b^{2}\)，猜想过 \(M(x_{0},y_{0})\) 的切线方程为 \(x_{0}x/a^{2}+y_{0}y/b^{2}=1\)． \((2)\) 设 \(A(x_{1},y_{1})\)、\(B(x_{2},y_{2})\)．由 \((1)\)，过椭圆上点 \(A\) 的切线 \(l_{1}\) 为 \(x_{1}x/a^{2}+y_{1}y/b^{2}=1\)；\(l_{1}\) 过 \(M(x_{0},y_{0})\)，故 \(x_{1}x_{0}/a^{2}+y_{1}y_{0}/b^{2}=1\)．同理 \(x_{2}x_{0}/a^{2}+y_{2}y_{0}/b^{2}=1\)． 两式说明 \(A\)、\(B\) 的坐标都满足同一方程 \(x_{0}x/a^{2}+y_{0}y/b^{2}=1\)，而两点确定一条直线，故直线 \(AB\) 的方程为 \(x_{0}x/a^{2}+y_{0}y/b^{2}=1\)． \((3)\) 由 \((2)\)，\(AB\)：\(x_{0}x/a^{2}+y_{0}y/b^{2}=1\)，因 \(M\) 不在坐标轴上，\(x_{0}y_{0}\neq 0\)，故 \(k_{AB}=-(b^{2}x_{0})/(a^{2}y_{0})\)；又 \(k_{OM}=y_{0}/x_{0}\)， 所以 \(k_{AB}\cdot k_{OM}=-(b^{2}x_{0})/(a^{2}y_{0})\cdot y_{0}/x_{0}=-b^{2}/a^{2}\)，为定值（与切点位置无关）． 设 \(AB\) 的中点为 \(P((x_{1}+x_{2})/2, (y_{1}+y_{2})/2)\)．\(A\)、\(B\) 均在椭圆上：\(x_{1}^{2}/a^{2}+y_{1}^{2}/b^{2}=1\)，\(x_{2}^{2}/a^{2}+y_{2}^{2}/b^{2}=1\)，两式相减得 \((x_{2}^{2}-x_{1}^{2})/a^{2}+(y_{2}^{2}-y_{1}^{2})/b^{2}=0\)，即 \((x_{2}+x_{1})(x_{2}-x_{1})/a^{2}+(y_{2}+y_{1})(y_{2}-y_{1})/b^{2}=0\)， 于是 \(k_{AB}\cdot k_{OP}=[(y_{2}-y_{1})/(x_{2}-x_{1})]\cdot [(y_{2}+y_{1})/(x_{2}+x_{1})]=-b^{2}/a^{2}\)． 结合已证的 \(k_{AB}\cdot k_{OM}=-b^{2}/a^{2}\)，且 \(k_{AB}\neq 0\)（\(M\) 不在坐标轴上使 \(x_{0}y_{0}\neq 0\)），得 \(k_{OP}=k_{OM}\)， 故 \(O\)、\(P\)、\(M\) 三点共线，即直线 \(OM\) 过弦 \(AB\) 的中点 \(P\)，所以 \(OM\) 平分线段 \(AB\)．}
\fi
\end{ansblock}

\begin{ansblock}[2章-拓-课时12-19]
% ans:2章-拓-课时12-19
\ansitem{19}{\((1) x^{2}/6+y^{2}/3=1\)；\((2) x^{2}+y^{2}=9\)}
\ifansbookdetail
\ansnote{详解}{\((1) e^{2}=c^{2}/a^{2}=1/2\)，又 \(a^{2}=b^{2}+c^{2}\)，故 \(a^{2}=2b^{2}\)．点 \((-2,1)\) 在椭圆上：\(4/a^{2}+1/b^{2}=1\)，即 \(4/(2b^{2})+1/b^{2}=3/b^{2}=1\)，得 \(b^{2}=3\)、\(a^{2}=6\)， 椭圆 \(C\)：\(x^{2}/6+y^{2}/3=1\)． \((2)\)（ⅰ）两切线中有一条斜率不存在时：该切线为 \(x=\pm \sqrt{6}\)，另一条与之垂直且与椭圆相切，只能是 \(y=\pm \sqrt{3}\)，交点 \(P\) 为 \((\pm \sqrt{6}, \pm \sqrt{3})\)（四组），均满足 \(6+3=9\)． （ⅱ）两切线斜率均存在时，设切线为 \(y-n=k(x-m)\)，与 \(x^{2}/6+y^{2}/3=1\)（即 \(x^{2}+2y^{2}=6\)）联立，消 \(y\) 得 \((1+2k^{2})x^{2}-4k(km-n)x+2(km-n)^{2}-6=0\)．相切 \(\Rightarrow  \Delta =16k^{2}(km-n)^{2}-4(1+2k^{2})[2(km-n)^{2}-6]=0\)， 展开约去同类项后得 \((m^{2}-6)k^{2}-2mnk+n^{2}-3=0\)（\(m^{2}\neq 6\) 时为关于 \(k\) 的二次方程）． 设两切线斜率为 \(k_{1}\)、\(k_{2}\)，则 \(k_{1}k_{2}=(n^{2}-3)/(m^{2}-6)\)，垂直 \(\Rightarrow  (n^{2}-3)/(m^{2}-6)=-1\)，即 \(m^{2}+n^{2}=9\)（\(m\neq \pm \sqrt{6}\)）． 结合（ⅰ）知斜率不存在的情形也在同一圆上，且该圆上每一点到原点的距离 \(3\) 大于椭圆上点的最大距离 \(\sqrt{6}\)，故圆上点均在椭圆外、两切线总存在， 所以动点 \(P\) 的轨迹方程为 \(x^{2}+y^{2}=9\)（即以 \(a^{2}+b^{2}=9\) 为半径平方的蒙日圆）．}
\fi
\end{ansblock}

\begin{ansblock}[2章-拓-课时12-20]
% ans:2章-拓-课时12-20
\ansitem{20}{\((0, \sqrt{6}/3)\)}
\ifansbookdetail
\ansnote{详解}{先定临界位置：对平面内一点 \(P\)（在椭圆外），过 \(P\) 的两条切线的夹角是 \(P\) 对椭圆上两点的张角 \(\angle APB\) 的最大值（\(A\)、\(B\) 沿切线方向时张角最大），而两切线夹角为直角 \(\Leftrightarrow  P\) 在蒙日圆上、为钝角 \(\Leftrightarrow  P\) 在蒙日圆内、为锐角 \(\Leftrightarrow  P\) 在蒙日圆外． 因此\("\)对椭圆上任意 \(A\)、\(B\) 与直线上任意 \(P\)，\(\angle APB\) 恒为锐角\("\Leftrightarrow \) 直线 \(3x+4y-10=0\) 上每一点都在蒙日圆外，即直线与蒙日圆相离． 椭圆 \(x^{2}/a^{2}+y^{2}=1\) 的 \(a^{2}=a^{2}\)、\(b^{2}=1\)，蒙日圆为 \(x^{2}+y^{2}=a^{2}+1\)，圆心 \(O(0,0)\)、半径 \(\sqrt{a^{2}+1}\)． 相离 \(\Leftrightarrow  d=|-10|/\sqrt{9+16}=2>\sqrt{a^{2}+1} \Leftrightarrow  a^{2}+1<4 \Leftrightarrow  a^{2}<3\)，结合 \(a>1\) 得 \(1<a<\sqrt{3}\)，即 \(a^{2}\in (1,3)\)． \(e^{2}=1-b^{2}/a^{2}=1-1/a^{2}\) 在 \(a^{2}\in (1,3)\) 上随 \(a^{2}\) 递增：\(a^{2}\rightarrow 1^{+}\) 时 \(e^{2}\rightarrow 0^{+}\)，\(a^{2}\rightarrow 3^{-}\) 时 \(e^{2}\rightarrow 2/3\)， 故 \(e^{2}\in (0, 2/3)\)，\(e\in (0, \sqrt{6}/3)\)．}
\fi
\end{ansblock}

\begin{ansblock}[2章-拓-课时12-21]
% ans:2章-拓-课时12-21
\ansitem{21}{\(\sqrt{6}/3\)}
\ifansbookdetail
\ansnote{详解}{设 \(P(m,n)\)，过 \(P\) 的切线 \(y-n=k(x-m)\) 与 \(x^{2}/a^{2}+y^{2}=1\)（即 \(x^{2}+a^{2}y^{2}=a^{2}\)）联立，消 \(y\) 得 \((1+a^{2}k^{2})x^{2}+2a^{2}k(n-km)x+a^{2}(n-km)^{2}-a^{2}=0\)，相切 \(\Rightarrow  \Delta =0\)，化简得 \((a^{2}-m^{2})k^{2}+2mnk+(1-n^{2})=0\)． 当两切线斜率均存在时，\(k_{1}k_{2}=(1-n^{2})/(a^{2}-m^{2})\)，\(PA\perp PB \Leftrightarrow  (1-n^{2})/(a^{2}-m^{2})=-1 \Leftrightarrow  m^{2}+n^{2}=a^{2}+1\)， 即满足 \(PA\perp PB\) 的点落在以 \(O\) 为圆心、\(\sqrt{a^{2}+1}\) 为半径的圆上（蒙日圆）． \("\)恰好存在一点 \(P\) 使 \(PA\perp PB" \Leftrightarrow \) 直线 \(ax+y-4=0\) 与该圆恰有一个公共点，即相切： \(d=|-4|/\sqrt{a^{2}+1}=\sqrt{a^{2}+1} \Rightarrow  a^{2}+1=4 \Rightarrow  a^{2}=3\)（\(a>1\) 故 \(a=\sqrt{3}\)），此时唯一点 \(P\) 为 \(O\) 到直线的垂足 \((\sqrt{3}, 1)\)（\(\sqrt{3}\cdot \sqrt{3}+1=4\) ，\(3+1=4=a^{2}+1\) ）． 该点 \(m^{2}=a^{2}\) 使上面关于 \(k\) 的方程退化，其两切线为 \(x=\sqrt{3}\)（竖直，切于右顶点）与 \(y=1\)（水平，切于上顶点），确实互相垂直，故取等情形成立、结论不受影响． 于是 \(a^{2}=3\)、\(b^{2}=1\)，\(c^{2}=2\)，\(e=c/a=\sqrt{2}/\sqrt{3}=\sqrt{6}/3\)．}
\fi
\end{ansblock}

\begin{ansblock}[2章-拓-课时12-22]
% ans:2章-拓-课时12-22
\ansitem{22}{\((1,11)\)}
\ifansbookdetail
\ansnote{详解}{两切线夹角为钝角 \(\Leftrightarrow  M\) 在蒙日圆内部（不含边界）；又要能引出两条切线，须 \(M\) 在椭圆外． 蒙日圆半径平方 \(r^{2}=a^{2}+b^{2}=16+9=25\)，即 \(M\) 在圆 \(x^{2}+y^{2}=25\) 内、椭圆外（开区域，两端边界均不含）． 圆心 \(O\) 到 \(l\)：\(x+y-6\sqrt{2}=0\) 的距离 \(d_{0}=6\sqrt{2}/\sqrt{2}=6>5\)，故直线 \(l\) 在蒙日圆外，圆上点到 \(l\) 的距离取值恰为 \((6-5, 6+5)=(1,11)\) 的上下确界． 对任一 \(v\in (1,11)\)，与 \(l\) 平行且相距 \(v\) 的直线与开圆盘相交成一条弦，弦的端点趋近圆上点；因椭圆上点到原点距离最大为 \(a=4<5\)，圆附近、弦上的点必在椭圆外，故该直线上的可取点确实存在，\(v\) 可达； 反过来，区域内的点到 \(l\) 的距离不可能达到 \(1\) 或 \(11\)（那需在圆外或圆上）． 故动点 \(M\) 到直线 \(l\) 的距离的取值范围是 \((1,11)\)（开区间，不含端点）。}
\fi
\end{ansblock}

\begin{ansblock}[2章-拓-课时12-23]
% ans:2章-拓-课时12-23
\ansitem{23}{\([2/3, \sqrt{2}/2]\)}
\ifansbookdetail
\ansnote{详解}{设 \(M(x_{0},y_{0})\)（\(x_{0}^{2}+y_{0}^{2}=3\)）、\(A(x_{1},y_{1})\)、\(B(x_{2},y_{2})\)． 由切线代换式（见 \(12-\)拓\(18\)）：过椭圆 \(x^{2}/2+y^{2}=1\) 上点 \((x_{1},y_{1})\) 的切线为 \(x_{1}x/2+y_{1}y=1\)，切线 \(MA\) 过 \(M \Rightarrow  x_{1}x_{0}/2+y_{1}y_{0}=1\)，同理 \(x_{2}x_{0}/2+y_{2}y_{0}=1\)， 故切点弦 \(AB\)：\(x_{0}x/2+y_{0}y=1\)． （ⅰ）\(y_{0}\neq 0\) 时，由直线解出 \(y=(1-x_{0}x/2)/y_{0}\) 代回 \(x^{2}+2y^{2}=2\)，整理并用 \(x_{0}^{2}=3-y_{0}^{2}\) 化简得 \((y_{0}^{2}+3)x^{2}-4x_{0}x+(4-4y_{0}^{2})=0\)，\(\therefore  x_{1}+x_{2}=4x_{0}/(y_{0}^{2}+3)\)，\(x_{1}x_{2}=(4-4y_{0}^{2})/(y_{0}^{2}+3)\)， \((x_{1}-x_{2})^{2}=(x_{1}+x_{2})^{2}-4x_{1}x_{2}=[16x_{0}^{2}-16(1-y_{0}^{2})(y_{0}^{2}+3)]/(y_{0}^{2}+3)^{2}=16y_{0}^{2}(y_{0}^{2}+1)/(y_{0}^{2}+3)^{2}\)， 又 \(y_{1}-y_{2}=-[x_{0}/(2y_{0})](x_{1}-x_{2})\)，故 \(|AB|=|x_{1}-x_{2}|\cdot \sqrt{}(1+x_{0}^{2}/(4y_{0}^{2}))=[4|y_{0}|\sqrt{y_{0}^{2}+1}/(y_{0}^{2}+3)]\cdot [\sqrt{4y_{0}^{2}+x_{0}^{2}}/(2|y_{0}|)]\)， 而 \(4y_{0}^{2}+x_{0}^{2}=4y_{0}^{2}+3-y_{0}^{2}=3(y_{0}^{2}+1)\)，所以 \(|AB|=2\sqrt{3}(y_{0}^{2}+1)/(y_{0}^{2}+3)\)． 原点 \(O\) 到 \(AB\) 的距离 \(d=1/\sqrt{x_{0}^{2}/4+y_{0}^{2}}=1/\sqrt{}(3(1+y_{0}^{2})/4)=2/[\sqrt{3}\cdot \sqrt{y_{0}^{2}+1}]\)， \(S=(1/2)|AB|\cdot d=(1/2)\cdot [2\sqrt{3}(y_{0}^{2}+1)/(y_{0}^{2}+3)]\cdot [2/(\sqrt{3}\sqrt{y_{0}^{2}+1})]=2\sqrt{y_{0}^{2}+1}/(y_{0}^{2}+3)\)． （ⅱ）\(y_{0}=0\) 时 \(x_{0}^{2}=3\)，\(AB\) 为 \(x=2/x_{0}=\pm 2\sqrt{3}/3\)，切点纵坐标满足 \(y^{2}=1-x^{2}/2=1-(2/\sqrt{3})^{2}/2=1-2/3=1/3\)， \(|AB|=2\sqrt{3}/3\)，\(d=2\sqrt{3}/3\)，\(S=(1/2)\cdot (2\sqrt{3}/3)\cdot (2\sqrt{3}/3)=2/3\)，与 \((\)ⅰ\()\) 式在 \(y_{0}=0\) 时的值一致，故通式覆盖全部情形． 令 \(t=\sqrt{y_{0}^{2}+1}\)，由 \(y_{0}^{2}\in [0,3]\) 得 \(t\in [1,2]\)，则 \(y_{0}^{2}+3=t^{2}+2\)，\(S=2t/(t^{2}+2)=2/(t+2/t)\)． \(t+2/t\) 在 \([1,2]\) 上：由对勾函数性质，在 \(t=\sqrt{2}\)（\(\in [1,2]\)）处取最小 \(2\sqrt{2}\)，在端点 \(t=1\)、\(t=2\) 处取最大 \(3\)， 故 \(t+2/t\in [2\sqrt{2}, 3]\)，\(S\in [2/3, 2/(2\sqrt{2})]=[2/3, \sqrt{2}/2]\)． 所以 \(\triangle AOB\) 面积的取值范围为 \([2/3, \sqrt{2}/2]\)．}
\fi
\end{ansblock}

\begin{ansblock}[2章-拓-课时12-24]
% ans:2章-拓-课时12-24
\ansitem{24}{\(C\)}
\ifansbookdetail
\ansnote{详解}{设过 \(Q(0,r)\) 的切线为 \(y=kx+r\)，与 \(x^{2}/12+y^{2}/4=1\)（即 \(x^{2}+3y^{2}=12\)）联立：\((1+3k^{2})x^{2}+6krx+3r^{2}-12=0\)， \(\Delta =36k^{2}r^{2}-4(1+3k^{2})(3r^{2}-12)=-12r^{2}+48+144k^{2}=0 \Rightarrow  k^{2}=(r^{2}-4)/12\)． 两条切线的斜率为该式的两根 \(k=\pm \sqrt{}((r^{2}-4)/12)\)，互相垂直 \(\Rightarrow \) 其积 \(-(r^{2}-4)/12=-1 \Rightarrow  r^{2}=16\)，圆 \(C\)：\(x^{2}+y^{2}=16\)（半径 \(4\)）． 直线 \(l\)：\(m(x-\sqrt{3})+(y-1)=0\) 恒过定点 \(P(\sqrt{3},1)\)，且 \(|OP|=\sqrt{3+1}=2<4\)，\(P\) 在圆内． 弦长 \(|AB|=2\sqrt{16-d^{2}}\)（\(d\) 为圆心 \(O\) 到 \(l\) 的距离），要使 \(|AB|\) 最小须 \(d\) 最大；\(l\) 恒过定点 \(P\)，故 \(d\leqslant |OP|=2\)，等号当 \(l\perp OP\) 时取得（此时 \(m\) 使 \(l\) 的方向与 \(OP\) 垂直，可取）． 所以 \(|AB|min=2\sqrt{16-4}=2\sqrt{12}=4\sqrt{3}\)．故选：\(C\)．}
\fi
\end{ansblock}

\begin{ansblock}[2章-拓-课时12-25]
% ans:2章-拓-课时12-25
\ansitem{25}{\(1\)}
\ifansbookdetail
\ansnote{详解}{\(a^{2}=4\)、\(b^{2}=1\)、\(c^{2}=3\)；蒙日圆半径平方 \(r^{2}=a^{2}+b^{2}=5\)（过四顶点作切线所得矩形的外接圆，对角线一半即 \(\sqrt{5}\)），圆方程 \(x^{2}+y^{2}=5\)，记圆心 \(O\)、半径 \(r=\sqrt{5}\)． 设 \(m=|PF_{1}|\cdot |PF_{2}|\)．由 \(|PF_{1}|+|PF_{2}|=2a=4\) 得 \(|PF_{1}|^{2}+|PF_{2}|^{2}=16-2m\)； 在 \(\triangle PF_{1}F_{2}\) 中 \(O\) 为 \(F_{1}F_{2}\) 的中点，由中线（阿波罗尼奥斯）定理 \(|PF_{1}|^{2}+|PF_{2}|^{2}=2(|PO|^{2}+|OF_{1}|^{2})=2|PO|^{2}+2c^{2}=2|PO|^{2}+6\)， 故 \(16-2m=2|PO|^{2}+6 \Rightarrow  |PO|^{2}=5-m\)． 又 \(|PF_{1}|\cdot |PF_{2}|=a^{2}-e^{2}x\_P^{2}=4-(3/4)x\_P^{2}\)，\(x\_P^{2}\in [0,4]\)，故 \(m\in [1,4]\)，从而 \(|PO|^{2}=5-m\in [1,4]<r^{2}=5\)，即 \(P\) 在蒙日圆内（与图相符，直线 \(l\) 与圆总交于两点）． 过圆内一点 \(P\) 的弦 \(MN\) 被 \(P\) 分成两段，由相交弦定理（或直接 \((r-|PO|)(r+|PO|)\)）得 \(|PM|\cdot |PN|=r^{2}-|PO|^{2}=5-(5-m)=m\)， 所以 \((|PM|\cdot |PN|)/(|PF_{1}|\cdot |PF_{2}|)=m/m=1\)．}
\fi
\end{ansblock}

\begin{ansblock}[2章-拓-课时12-26]
% ans:2章-拓-课时12-26
\ansitem{26}{\((1) 17/4\)；\((2)\) 一般结论：过圆 \(x^{2}+y^{2}=a^{2}+b^{2}\) 上任意一点 \(Q(m,n)\) 作椭圆 \(x^{2}/a^{2}+y^{2}/b^{2}=1 (a>b>0)\) 的两条切线，则这两条切线互相垂直；证明见详解}
\ifansbookdetail
\ansnote{详解}{\((1)\) 设 \(P(x_{0},y_{0})\)，则 \(x_{0}^{2}+y_{0}^{2}=34\) ①．椭圆 \(a^{2}=25\)、\(b^{2}=9 \Rightarrow  c^{2}=16\)，右焦点 \(F(4,0)\)． \(F\) 在线段 \(OP\) 的垂直平分线上 \(\Leftrightarrow  |FP|=|FO|=4\)，即 \((x_{0}-4)^{2}+y_{0}^{2}=16\) ②， ②展开为 \(x_{0}^{2}+y_{0}^{2}-8x_{0}+16=16\)，即 \(x_{0}^{2}+y_{0}^{2}-8x_{0}=0\)，代①：\(34-8x_{0}=0 \Rightarrow  x_{0}=34/8=17/4\)． （存在性核验：\(x_{0}=17/4<\sqrt{34}\)，\(y_{0}^{2}=34-289/16=255/16>0\) ）故点 \(P\) 的横坐标为 \(17/4\)． \((2)\) 一般结论：过圆 \(x^{2}+y^{2}=a^{2}+b^{2}\) 上任意一点 \(Q(m,n)\) 作椭圆 \(x^{2}/a^{2}+y^{2}/b^{2}=1 (a>b>0)\) 的两条切线，则这两条切线互相垂直（即题中①②所述\("\)圆 \(x^{2}+y^{2}=a^{2}+b^{2}"\)正是该椭圆的蒙日圆，①中 \(5^{2}+3^{2}=34\) 与 \((1)\) 的圆 \(O\) 同）。 证明：设 \(Q(m,n)\) 满足 \(m^{2}+n^{2}=a^{2}+b^{2}\)． （ⅰ）若过 \(Q\) 的切线中有一条斜率不存在，则该切线为 \(x=\pm a\)，于是 \(m=\pm a\)，由 \(m^{2}+n^{2}=a^{2}+b^{2}\) 得 \(n^{2}=b^{2}\)，即 \(n=\pm b\)，\(Q(\pm a,\pm b)\)； 此时直线 \(y=\pm b\) 也过 \(Q\) 且与椭圆相切（切于上、下顶点），它与 \(x=\pm a\) 垂直，故两切线互相垂直（再过 \(Q\) 不能有其他切线，因椭圆的切线必在 \(x=\pm a\)、\(y=\pm b\) 围成的矩形外或为其边）． （ⅱ）若过 \(Q\) 的两条切线斜率均存在，设切线为 \(y-n=k(x-m)\)，与 \(b^{2}x^{2}+a^{2}y^{2}=a^{2}b^{2}\) 联立，消 \(y\) 得 \((b^{2}+a^{2}k^{2})x^{2}+2a^{2}k(n-km)x+a^{2}(n-km)^{2}-a^{2}b^{2}=0\)， 相切 \(\Rightarrow  \Delta =4a^{4}k^{2}(n-km)^{2}-4(b^{2}+a^{2}k^{2})[a^{2}(n-km)^{2}-a^{2}b^{2}]=0\)，两边除 \(4a^{2}\) 并整理得 \((m^{2}-a^{2})k^{2}-2mnk+(n^{2}-b^{2})=0\)（ⅱ 中两根即两切线斜率；若 \(m^{2}=a^{2}\) 则由 \(a^{2}+b^{2}=m^{2}+n^{2}\) 得 \(n^{2}=b^{2}\)，即情形（ⅰ））． 故 \(k_{1}k_{2}=(n^{2}-b^{2})/(m^{2}-a^{2})\)；又 \(m^{2}+n^{2}=a^{2}+b^{2} \Rightarrow  m^{2}-a^{2}=b^{2}-n^{2}=-(n^{2}-b^{2})\)， 所以 \(k_{1}k_{2}=(n^{2}-b^{2})/[-(n^{2}-b^{2})]=-1\)，两切线互相垂直． 综合（ⅰ）（ⅱ），一般结论成立．}
\fi
\end{ansblock}

\begin{ansblock}[2章-拓-课时12-28]
% ans:2章-拓-课时12-28
\ansitem{28}{\(B\)}
\ifansbookdetail
\ansnote{详解}{\(\angle F_{1}PF_{2}=90^{\circ} \Leftrightarrow  P\) 在以 \(F_{1}F_{2}\) 为直径、即以原点 \(O\) 为圆心、\(c\) 为半径的圆上，故题设等价于该圆与椭圆有公共点． 椭圆上点到中心 \(O\) 的距离：\(|OP|^{2}=x^{2}+y^{2}=a^{2}-(a^{2}-b^{2})x^{2}/a^{2}\)（用 \(y^{2}=b^{2}(1-x^{2}/a^{2})\)），\(x^{2}\in [0,a^{2}]\) 时随 \(x^{2}\) 递减，故 \(|OP|^{2}\in [b^{2},a^{2}]\)，即 \(|OP|\in [b,a]\)． 于是圆 \(x^{2}+y^{2}=c^{2}\) 与椭圆有交点 \(\Leftrightarrow  b\leqslant c\leqslant a\)；而 \(c<a\) 恒成立（\(c^{2}=a^{2}-b^{2}<a^{2}\)），故条件即 \(b\leqslant c\)． \(b^{2}\leqslant c^{2} \Rightarrow  a^{2}-c^{2}\leqslant c^{2} \Rightarrow  a^{2}\leqslant 2c^{2} \Rightarrow  e^{2}\geqslant 1/2\)，结合 \(0<e<1\) 得 \(e\in [\sqrt{2}/2, 1)\)．故选：\(B\)．}
\fi
\end{ansblock}

\begin{ansblock}[2章-拓-课时12-29]
% ans:2章-拓-课时12-29
\ansitem{29}{错误}
\ifansbookdetail
\ansnote{详解}{椭圆的离心率 \(e=\sqrt{1-b^{2}/a^{2}}\) 只含 \(a\) 与 \(b\) 的比值这一个信息，而确定一个椭圆需要 \(a\)、\(b\) 两个量（还需定位焦点在哪个轴上），故 \(e\) 相同只保证形状相似，不保证大小与方位相同． 反例：\(x^{2}/4+y^{2}/3=1\) 与 \(x^{2}/8+y^{2}/6=1\)，两者的 \(b^{2}/a^{2}\) 同为 \(3/4\)，\(e\) 均为 \(1/2\)，但长轴长分别为 \(4\) 与 \(4\sqrt{2}\)，显然不是同一个椭圆；再把第二个改写为 \(y^{2}/8+x^{2}/6=1\)，则离心率仍为 \(1/2\) 而焦点改在 \(y\) 轴上． 故\("\)离心率相同的椭圆是同一个椭圆\("\)错误（正确说法是：离心率相同的椭圆形状相同\(/\)相似，长短轴之比相同）．}
\fi
\end{ansblock}

\begin{ansblock}[2章-拓-课时12-30]
% ans:2章-拓-课时12-30
\ansitem{30}{\(D\)}
\ifansbookdetail
\ansnote{详解}{\(2a=8\sqrt{2} \Rightarrow  a=4\sqrt{2}\)、\(a^{2}=32\)；\(e=\sqrt{2}/2 \Rightarrow  c=a\cdot e=4\)、\(c^{2}=16\)，\(b^{2}=32-16=16\)，椭圆为 \(x^{2}/32+y^{2}/16=1\)，\(F_{2}(4,0)\)． 由定义 \(|MF_{1}|=2a-|MF_{2}|\)，故 \(|MF_{1}|+|MB|=2a+(|MB|-|MF_{2}|)=8\sqrt{2}+(|MB|-|MF_{2}|)\)． 对任意 \(M\)，由三角形两边之差：\(|MB|-|MF_{2}|\geqslant -|MF_{2}|+|MB|\geqslant -|BF_{2}|\)（即 \(||MB|-|MF_{2}||\leqslant |BF_{2}|\)）， 而 \(|BF_{2}|=\sqrt{}((3-4)^{2}+(2-0)^{2})=\sqrt{5}\)，故 \(|MF_{1}|+|MB|\geqslant 8\sqrt{2}-\sqrt{5}\)． 等号条件：\(|MF_{2}|-|MB|=|BF_{2}|\)，即 \(B\) 在线段 \(MF_{2}\) 上，亦即 \(M\) 在射线 \(F_{2}B\)（自 \(F_{2}\) 过 \(B\) 向外）上；因 \(B(3,2)\) 满足 \(9/32+4/16=17/32<1\) 在椭圆内，该射线必与椭圆交于一点 \(M\)，故等号可取到． 所以最小值为 \(8\sqrt{2}-\sqrt{5}\)．故选：\(D\)．}
\fi
\end{ansblock}

\begin{ansblock}[2章-拓-课时12-31]
% ans:2章-拓-课时12-31
\ansitem{31}{\(C\)}
\ifansbookdetail
\ansnote{详解}{设 \(M(2t+6, t)\)（在直线上，\(x_{0}=2y_{0}+6\)），\(A(x_{1},y_{1})\)、\(B(x_{2},y_{2})\)． 由题给切线公式，切线 \(MA\)：\(x_{1}x/3+y_{1}y/2=1\)，它过 \(M \Rightarrow  x_{1}(2t+6)/3+y_{1}t/2=1\)；同理 \(x_{2}(2t+6)/3+y_{2}t/2=1\)． 故 \(A\)、\(B\) 的坐标同满足方程 \((2t+6)x/3+ty/2=1\)，此即直线 \(AB\) 的方程． 两边乘 \(6\)：\(2(2t+6)x+3ty=6 \Rightarrow  (4x+3y)t+(12x-6)=0\) 对一切 \(t\) 成立 \(\Rightarrow  4x+3y=0\) 且 \(12x-6=0\)， 解得 \(x=1/2\)、\(y=-2/3\)（满足 \(4\cdot (1/2)+3\cdot (-2/3)=2-2=0\) ）． （可行性核验：直线上点 \(M(2t+6,t)\) 到椭圆的定位 \((2t+6)^{2}/3+t^{2}/2=(11/6)t^{2}+8t+12\)，其最小值在 \(t=-24/11\) 处为 \(36/11>1\)，故 \(M\) 恒在椭圆外，两切线总存在．） 所以直线 \(AB\) 恒过定点 \((1/2, -2/3)\)．故选：\(C\)．}
\fi
\end{ansblock}

\begin{ansblock}[2章-拓-课时12-32]
% ans:2章-拓-课时12-32
\ansitem{32}{\((1) x^{2}/4+y^{2}=1\)；\((2)\) 证明见详解，定点为 \((0, -3/5)\)}
\ifansbookdetail
\ansnote{详解}{\((1) 2a=4 \Rightarrow  a=2\)；\(e=c/a=\sqrt{3}/2 \Rightarrow  c=\sqrt{3}\)，\(b^{2}=a^{2}-c^{2}=1\)，椭圆 \(E\)：\(x^{2}/4+y^{2}=1\)． \((2)\) 先设直线 \(PQ\) 斜率存在，方程 \(y=kx+m\)（\(m\neq 1\)，否则直线过 \(N(0,1)\)，\(P\)、\(Q\) 中有一点与 \(N\) 重合，\(\angle PNQ\) 无意义）． 与 \(x^{2}+4y^{2}=4\) 联立：\((1+4k^{2})x^{2}+8kmx+4(m^{2}-1)=0\)，须 \(\Delta =64k^{2}m^{2}-16(1+4k^{2})(m^{2}-1)=16[(1+4k^{2})-m^{2}]>0\)， 且 \(x_{1}+x_{2}=-8km/(1+4k^{2})\)，\(x_{1}x_{2}=4(m^{2}-1)/(1+4k^{2})\)． \(\angle PNQ=90^{\circ} \Leftrightarrow \) 向量\(NP\cdot NQ=0\)，其中 \(NP=(x_{1}, y_{1}-1)\)、\(NQ=(x_{2}, y_{2}-1)\)： \(x_{1}x_{2}+(kx_{1}+m-1)(kx_{2}+m-1)=(1+k^{2})x_{1}x_{2}+k(m-1)(x_{1}+x_{2})+(m-1)^{2}=0\)， 代韦达并乘 \((1+4k^{2})\)：\((1+k^{2})\cdot 4(m^{2}-1)+k(m-1)(-8km)+(m-1)^{2}(1+4k^{2})=0\)， 提取 \((m-1)\)（\(m\neq 1\)）：\(4(1+k^{2})(m+1)-8k^{2}m+(m-1)(1+4k^{2})=0\)， 展开：\(4+4k^{2}+4m+4k^{2}m-8k^{2}m+m+4k^{2}m-1-4k^{2}=0 \Rightarrow  3+5m=0 \Rightarrow  m=-3/5\)（满足 \(\Delta >0\)：\(1+4k^{2}>9/25\) 恒成立）． 故 \(PQ\)：\(y=kx-3/5\) 过定点 \((0, -3/5)\)． 再设直线 \(PQ\) 斜率不存在（\(x=s\)）：\(P(s,u)\)、\(Q(s,-u)\)（\(u^{2}=1-s^{2}/4\)），则 \(NP\cdot NQ=s^{2}+(u-1)(-u-1)=s^{2}+1-u^{2}=s^{2}+s^{2}/4=(5/4)s^{2}\)， 要为 \(0\) 须 \(s=0\)，此时 \(P(0,\pm 1)\)，其中 \((0,1)\) 即 \(N\)，\(\angle PNQ\) 不存在，故此情形不合题意． 综上，直线 \(PQ\) 恒过定点 \((0, -3/5)\)．}
\fi
\end{ansblock}

\begin{ansblock}[2章-拓-课时12-33]
% ans:2章-拓-课时12-33
\ansitem{33}{\((\)Ⅰ\() x^{2}/4+y^{2}=1\)；\((\)Ⅱ\() 1/5\)；\((\)Ⅲ\()\) 证明见详解，定点为 \((4/3, 0)\)}
\ifansbookdetail
\ansnote{详解}{\((\)Ⅰ\() e^{2}=c^{2}/a^{2}=3/4 \Rightarrow  a^{2}=4b^{2}\)（由 \(c^{2}=a^{2}-b^{2}\)）；点 \((1, \sqrt{3}/2)\) 代入：\(1/a^{2}+3/(4b^{2})=1\)，即 \(1/(4b^{2})+3/(4b^{2})=1/b^{2}=1 \Rightarrow  b^{2}=1\)、\(a^{2}=4\)，\(E\)：\(x^{2}/4+y^{2}=1\)． \((\)Ⅱ\() A(-2,0)\)、\(B(2,0)\)，设 \(P(3,t)\)（\(t\neq 0\)），\(k_{1}=t/(3+2)=t/5\)、\(k_{2}=t/(3-2)=t\)，故 \(k_{1}/k_{2}=1/5\)． \((\)Ⅲ\() PA\)：\(y=(t/5)(x+2)\)，与 \(x^{2}+4y^{2}=4\) 联立得 \((4t^{2}+25)x^{2}+16t^{2}x+16t^{2}-100=0\)． 其一根为 \(x\_A=-2\)，由韦达 \(x\_A+x\_C=-16t^{2}/(4t^{2}+25) \Rightarrow  x\_C=-16t^{2}/(4t^{2}+25)+2=(50-8t^{2})/(4t^{2}+25)\)，\(y\_C=(t/5)(x\_C+2)=20t/(4t^{2}+25)\)． \(PB\)：\(y=t(x-2)\)，联立得 \((1+4t^{2})x^{2}-16t^{2}x+16t^{2}-4=0\)，一根为 \(x\_B=2\)，\(x\_B+x\_D=16t^{2}/(1+4t^{2}) \Rightarrow  x\_D=(8t^{2}-2)/(1+4t^{2})\)，\(y\_D=t(x\_D-2)=-4t/(1+4t^{2})\)． 直线 \(CD\) 令 \(y=0\) 求横截距：\(x_{0}=x\_C-y\_C(x\_D-x\_C)/(y\_D-y\_C)\)，代上述四式化简得 \(x_{0}=4/3\)（与 \(t\) 无关）； （独立核验：直接取 \(t=1\)、\(t=2\) 分别算得 \(C\)、\(D\) 坐标，两点连线与 \(x\) 轴交点均为 \((4/3,0)\) ） 故直线 \(CD\) 恒过定点 \((4/3, 0)\)．}
\fi
\end{ansblock}

\begin{ansblock}[2章-拓-课时12-34]
% ans:2章-拓-课时12-34
\ansitem{34}{\((1) x^{2}/4+y^{2}/8=1\)；\((2)\) 存在，\(N(4,0)\)}
\ifansbookdetail
\ansnote{详解}{\((1) a^{2}\) 在 \(y\) 项下，故长轴在 \(y\) 轴、短轴在 \(x\) 轴，短轴右端点 \(A(b,0)\)．\(M(1,0)\) 为 \(OA\) 中点 \(\Rightarrow  b=2\)． \(e^{2}=c^{2}/a^{2}=1/2\) 且 \(c^{2}=a^{2}-b^{2} \Rightarrow  a^{2}-4=a^{2}/2 \Rightarrow  a^{2}=8\)，故 \(\Gamma \)：\(x^{2}/4+y^{2}/8=1\)． \((2)\) 因 \(M\)、\(N\) 均在 \(x\) 轴上，\(\angle PNM=\angle QNM\) 表示射线 \(NP\)、\(NQ\) 关于 \(x\) 轴对称，即 \(k_{PN}+k_{QN}=0\)（\(k_{PN}=k_{QN}\) 会使 \(P\)、\(Q\)、\(N\) 共线，此时两角互补不相等，舍）． 设 \(N(x_{0},0)\)．当 \(PQ\perp x\) 轴时，\(PQ\)：\(x=1\)，\(P\)、\(Q\) 关于 \(x\) 轴对称，对任意 \(x_{0}\) 恒有 \(\angle PNM=\angle QNM\)，不产生约束． 当 \(PQ\) 不垂直 \(x\) 轴时，设 \(PQ\)：\(y=k(x-1)\)，代入 \(2x^{2}+y^{2}=8\) 得 \((k^{2}+2)x^{2}-2k^{2}x+k^{2}-8=0\)（\(\Delta =4k^{4}-4(k^{2}+2)(k^{2}-8)=24k^{2}+64>0\) 恒成立）， \(x_{1}+x_{2}=2k^{2}/(k^{2}+2)\)，\(x_{1}x_{2}=(k^{2}-8)/(k^{2}+2)\)． \(k_{PN}+k_{QN}=k(x_{1}-1)/(x_{1}-x_{0})+k(x_{2}-1)/(x_{2}-x_{0})=k[2x_{1}x_{2}-(x_{0}+1)(x_{1}+x_{2})+2x_{0}]/[(x_{1}-x_{0})(x_{2}-x_{0})]\)， 其中分子方括号\(=[2(k^{2}-8)-(x_{0}+1)\cdot 2k^{2}+2x_{0}(k^{2}+2)]/(k^{2}+2)=(4x_{0}-16)/(k^{2}+2)=4(x_{0}-4)/(k^{2}+2)\)， 故 \(k_{PN}+k_{QN}=4k(x_{0}-4)/\{(k^{2}+2)(x_{1}-x_{0})(x_{2}-x_{0})\}\)． （分母不为 \(0\)：椭圆上点横坐标 \(|x|\leqslant 2\)，取 \(x_{0}=4\) 时 \(x\_i-x_{0}\neq 0\) ） 要对一切实数 \(k\) 成立，须 \(x_{0}-4=0\)，即 \(x_{0}=4\)． 综上，\(x\) 轴上存在定点 \(N(4,0)\)，使 \(\angle PNM=\angle QNM\) 恒成立．}
\fi
\end{ansblock}

\begin{ansblock}[2章-拓-课时12-35]
% ans:2章-拓-课时12-35
\ansitem{35}{\((1) x^{2}/4+y^{2}=1\)；\((2)\) 存在，\(k=\sqrt{3}/2\) 或 \(k=\pm \sqrt{11}/2\)}
\ifansbookdetail
\ansnote{详解}{\((1) 2b=2 \Rightarrow  b=1\)；\(e=c/a=\sqrt{3}/2\) 且 \(a^{2}=b^{2}+c^{2} \Rightarrow  c^{2}=3a^{2}/4 \Rightarrow  a^{2}/4=1 \Rightarrow  a^{2}=4\)、\(c=\sqrt{3}\)，\(C\)：\(x^{2}/4+y^{2}=1\)． \((2)\) 若 \(PAMN\) 为平行四边形，则 \(PA\parallel MN\) 且 \(|PA|=|MN|\)．由 \(PA\parallel MN\) 且过 \(A(2,0)\)，\(PA\)：\(y=k(x-2)\)，与 \(x=3\) 交于 \(P(3,k)\)，故 \(|PA|=\sqrt{1+k^{2}}\)． 联立 \(y=kx+\sqrt{3}\) 与 \(x^{2}+4y^{2}=4\)：\((4k^{2}+1)x^{2}+8\sqrt{3}kx+8=0\)，\(\Delta =192k^{2}-32(4k^{2}+1)=64k^{2}-32>0 \Leftrightarrow  k^{2}>1/2\)， \(x_{1}+x_{2}=-8\sqrt{3}k/(4k^{2}+1)\)，\(x_{1}x_{2}=8/(4k^{2}+1)\)， \(|MN|=\sqrt{1+k^{2}}|x_{1}-x_{2}|=\sqrt{1+k^{2}}\cdot \sqrt{64k^{2}-32}/(4k^{2}+1)\)． \(|PA|=|MN| \Rightarrow  \sqrt{64k^{2}-32}=4k^{2}+1\)，平方：\(64k^{2}-32=16k^{4}+8k^{2}+1 \Rightarrow  16k^{4}-56k^{2}+33=0\)， 即 \((4k^{2}-3)(4k^{2}-11)=0 \Rightarrow  k^{2}=3/4\) 或 \(k^{2}=11/4\)，得 \(k=\pm \sqrt{3}/2\)、\(k=\pm \sqrt{11}/2\)（均满足 \(k^{2}>1/2\) ）． 逐一检验四边形是否退化：\(k=-\sqrt{3}/2\) 时直线 \(MN\) 为 \(y=-(\sqrt{3}/2)x+\sqrt{3}\)，令 \(x=2\) 得 \(y=0\)，即直线 \(MN\) 过 \(A(2,0)\)，此时 \(M\)、\(N\) 中有一点与 \(A\) 重合（解 \((4k^{2}+1)x^{2}+8\sqrt{3}kx+8=0\) 得 \(x=1\) 或 \(x=2\)），\(A\)、\(M\)、\(N\) 三点共线，不能构成平行四边形，舍去； 其余三值：\(k=\sqrt{3}/2\) 时 \(MN\) 不过 \(A\)（\(x=2\) 处 \(y=2\sqrt{3}\neq 0\)），\(k=\pm \sqrt{11}/2\) 时 \(x=2\) 处 \(y=\pm \sqrt{11}+\sqrt{3}\neq 0\)，且 \(|PA|=|MN|\)、\(PA\parallel MN\) 同时成立，命名 \(M\)、\(N\) 使 向量\(PA=\)向量\(NM\) 即得平行四边形，均符合． 故 \(k=\sqrt{3}/2\) 或 \(k=\pm \sqrt{11}/2\)．}
\fi
\end{ansblock}

\begin{ansblock}[2章-拓-课时12-36]
% ans:2章-拓-课时12-36
\ansitem{36}{\((1) x^{2}/16+y^{2}/15=1\)；\((2) 8+\sqrt{13}\)；\((3)\) 不存在}
\ifansbookdetail
\ansnote{详解}{\((1)\) 圆心 \(M(-1,0)\)、半径 \(8\)，\(N(1,0)\) 在圆内（\(|MN|=2<8\)）．\(P\) 在 \(QN\) 的垂直平分线上 \(\Rightarrow  |PN|=|PQ|\)；又 \(P\) 在半径 \(MQ\) 上 \(\Rightarrow  |PM|+|PQ|=8\)， 故 \(|PM|+|PN|=|PM|+|PQ|=8>|MN|=2\)，由椭圆定义，\(P\) 的轨迹是以 \(M\)、\(N\) 为焦点、长轴长 \(8\) 的椭圆：\(2a=8\)、\(2c=2 \Rightarrow  a=4\)、\(c=1\)、\(b^{2}=15\)， \(E\)：\(x^{2}/16+y^{2}/15=1\)． \((2) N\) 为右焦点，\(M\) 为左焦点，\(|SN|+|SM|=2a=8 \Rightarrow  |SN|=8-|SM|\)， \(|SN|+|SB|=8+(|SB|-|SM|)\leqslant 8+|MB|\)（三角形两边之差小于第三边）， \(|MB|=\sqrt{}((2+1)^{2}+2^{2})=\sqrt{13}\)； 等号当 \(M\) 在线段 \(SB\) 上，即 \(S\) 在射线 \(BM\) 自 \(M\) 向外的部分上——\(M(-1,0)\) 在椭圆内（\(1/16<1\)），该射线与椭圆有唯一交点 \(S\)，故等号可取到，最大值为 \(8+\sqrt{13}\)． \((3)\) 不存在．理由：\(|MN|=2\)，若存在 \(T\) 使 \(2/|MN|=1/|TM|+1/|TN|\)，即 \(1=(|TM|+|TN|)/(|TM|\cdot |TN|)\)， 由定义 \(|TM|+|TN|=8\)，得 \(|TM|\cdot |TN|=8\)，故 \(|TM|\)、\(|TN|\) 是 \(u^{2}-8u+8=0\) 的两根，\(u=4\pm 2\sqrt{2}\)； 另一方面，\(T\) 在椭圆上，焦半径 \(|TM|=a+ex\_T\)（\(e=1/4\)）\(\in [a-c, a+c]=[3,5]\)，同理 \(|TN|\in [3,5]\)； 但 \(4-2\sqrt{2}\approx 1.17<3\)，\(4+2\sqrt{2}\approx 6.83>5\)，均落在 \([3,5]\) 之外（亦可从乘积看：\(u\in [3,5]\) 时 \(u(8-u)\) 在 \([3,4]\) 上递增、\([4,5]\) 上递减，取值 \([15,16]\)，恒大于 \(8\)）， 故满足条件的点 \(T\) 不存在．}
\fi
\end{ansblock}

\begin{ansblock}[2章-拓-课时12-37]
% ans:2章-拓-课时12-37
\ansitem{37}{定值为 \(-1\)}
\ifansbookdetail
\ansnote{详解}{设 \(P(x_{1},y_{1})\)（\(x_{1}>0\)、\(y_{1}>0\)）、\(B(x_{2},y_{2})\)，由 \(P\)、\(A\) 关于原点对称得 \(A(-x_{1},-y_{1})\)，又 \(C(x_{1},0)\)． \(A\)、\(B\)、\(P\) 均在椭圆上，即 \(x^{2}+2y^{2}=4\)，故 \(y\_i^{2}=2-x\_i^{2}/2\)（\(i=1,2\)）． 因 \(A\)、\(C\)、\(B\) 三点共线（\(B\) 在 \(AC\) 的延长线上），\(k_{AB}=k_{AC}\)，而 \(k_{AB}\cdot k_{PB}=[(y_{2}+y_{1})/(x_{2}+x_{1})]\cdot [(y_{2}-y_{1})/(x_{2}-x_{1})]=(y_{2}^{2}-y_{1}^{2})/(x_{2}^{2}-x_{1}^{2})=[(2-x_{2}^{2}/2)-(2-x_{1}^{2}/2)]/(x_{2}^{2}-x_{1}^{2})=-1/2\)， （此即第三定义：椭圆上两点与关于原点对称的两点连线斜率之积\(=(b^{2}/a^{2})-1=2/4-1=-1/2\) 的变式，此处 \(x^{2}/4+y^{2}/2=1\) 的 \(a^{2}=4\)、\(b^{2}=2\)） 又 \(k_{AC}=(0+y_{1})/(x_{1}+x_{1})=y_{1}/(2x_{1})\)，\(k_{PA}=(y_{1}+y_{1})/(x_{1}+x_{1})=y_{1}/x_{1}\)，故 \(k_{PA}=2k_{AC}=2k_{AB}\)， 所以 \(k_{PA}\cdot k_{PB}=2k_{AB}\cdot k_{PB}=2\times (-1/2)=-1\)，为定值（与 \(P\) 的位置无关）．}
\fi
\end{ansblock}

\begin{ansblock}[2章-拓-课时12-38]
% ans:2章-拓-课时12-38
\ansitem{38}{\((1) x^{2}/18+y^{2}/9=1\)；\((2) k_{1}+k_{2}\) 为定值 \(0\)，证明见详解}
\ifansbookdetail
\ansnote{详解}{\((1) c=3\)，\(A(0,b)\)．向量\(AF=(3,-b)\)，由向量\(AF=3\)向量\(FB\) 得向量\(FB=(1,-b/3)\)，故 \(B=F+\)向量\(FB=(4,-b/3)\)． \(B\) 在椭圆上：\(16/a^{2}+(b^{2}/9)/b^{2}=1 \Rightarrow  16/a^{2}+1/9=1 \Rightarrow  a^{2}=18\)，\(b^{2}=a^{2}-c^{2}=18-9=9\)， \(E\)：\(x^{2}/18+y^{2}/9=1\)（即 \(x^{2}+2y^{2}=18\)）． \((2)\) 当 \(l\) 不为 \(x\) 轴时，设 \(l\)：\(x=my+3\)（此式可表示除 \(x\) 轴外一切过 \(F\) 的直线，其中 \(m=0\) 即垂直于 \(x\) 轴的直线 \(x=3\)），设 \(A(x_{1},y_{1})\)、\(B(x_{2},y_{2})\)， 代入 \(x^{2}+2y^{2}=18\)：\((m^{2}+2)y^{2}+6my-9=0\)（\(\Delta =36m^{2}+36(m^{2}+2)>0\) 恒成立），\(y_{1}+y_{2}=-6m/(m^{2}+2)\)，\(y_{1}y_{2}=-9/(m^{2}+2)\)． \(k_{1}+k_{2}=y_{1}/(x_{1}-6)+y_{2}/(x_{2}-6)=y_{1}/(my_{1}-3)+y_{2}/(my_{2}-3)=[2my_{1}y_{2}-3(y_{1}+y_{2})]/[(my_{1}-3)(my_{2}-3)]\)， 分子\(=2m\cdot (-9/(m^{2}+2))-3\cdot (-6m/(m^{2}+2))=(-18m+18m)/(m^{2}+2)=0\)， 分母\(=m^{2}y_{1}y_{2}-3m(y_{1}+y_{2})+9=-9m^{2}/(m^{2}+2)+18m^{2}/(m^{2}+2)+9=9m^{2}/(m^{2}+2)+9=(18m^{2}+18)/(m^{2}+2)\neq 0\)， 故 \(k_{1}+k_{2}=0\)． 当 \(l\) 为 \(x\) 轴（\(y=0\)）时，\(A\)、\(B\) 为左右顶点 \((\pm 3\sqrt{2},0)\)，\(k_{1}=k_{2}=0\)，仍有 \(k_{1}+k_{2}=0\)． （另核验：\(P(6,0)\) 在椭圆外（\(36/18=2>1\)），且 \(6\notin [-3\sqrt{2},3\sqrt{2}]\)，故 \(x\_i-6\neq 0\)，斜率总有意义．） 所以 \(k_{1}+k_{2}\) 为定值 \(0\)．}
\fi
\end{ansblock}

\begin{ansblock}[2章-拓-课时12-39]
% ans:2章-拓-课时12-39
\ansitem{39}{\((1) x^{2}/4+y^{2}=1\)；\((2)\) 定值为 \(2\)，证明见详解}
\ifansbookdetail
\ansnote{详解}{\((1) A(2,0)\) 在 \(x\) 轴上、\(B(0,1)\) 在 \(y\) 轴上，故 \(a=2\)、\(b=1\)（\(a>b>0\) 满足），\(C\)：\(x^{2}/4+y^{2}=1\)． \((2)\) 设 \(P(x_{0},y_{0})\)，第三象限故 \(x_{0}\in (-2,0)\)、\(y_{0}\in (-1,0)\)，且 \(x_{0}^{2}+4y_{0}^{2}=4\)． 直线 \(PA\)：\(y=y_{0}/(x_{0}-2)\cdot (x-2)\)，令 \(x=0\) 得 \(M(0, -2y_{0}/(x_{0}-2))\)；直线 \(PB\)：\(y-1=(y_{0}-1)/x_{0}\cdot x\)，令 \(y=0\) 得 \(N(-x_{0}/(y_{0}-1), 0)\)． 因 \(x_{0}-2<0\)、\(-2y_{0}>0\)，\(M\) 的纵坐标\(<0\)；\(y_{0}-1<0\)、\(-x_{0}>0\)，\(N\) 的横坐标\(<0\)，即 \(M\) 在 \(y\) 轴负半轴、\(N\) 在 \(x\) 轴负半轴． 四边形 \(ABNM\) 的顶点依次为 \(A(2,0)\)、\(B(0,1)\)、\(N(\)负\(,0)\)、\(M(0,\)负\()\)，两条对角线 \(AN\) 在 \(x\) 轴上、\(BM\) 在 \(y\) 轴上，互相垂直且相交于原点， 故 \(S=(1/2)|AN|\cdot |BM|=(1/2)\cdot |2+x_{0}/(y_{0}-1)|\cdot |1+2y_{0}/(x_{0}-2)| =(1/2)\cdot |(x_{0}+2y_{0}-2)/(y_{0}-1)|\cdot |(x_{0}+2y_{0}-2)/(x_{0}-2)|=(1/2)\cdot |(x_{0}+2y_{0}-2)^{2}/[(y_{0}-1)(x_{0}-2)]|\)． 分子展开并用 \(x_{0}^{2}+4y_{0}^{2}=4\)：\((x_{0}+2y_{0}-2)^{2}=x_{0}^{2}+4y_{0}^{2}+4+4x_{0}y_{0}-4x_{0}-8y_{0}=8+4x_{0}y_{0}-4x_{0}-8y_{0}=4(x_{0}y_{0}-x_{0}-2y_{0}+2)=4(y_{0}-1)(x_{0}-2)\)， 所以 \(S=(1/2)\cdot |4|=2\)，与 \(P\) 的位置无关，即四边形 \(ABNM\) 的面积为定值 \(2\)．}
\fi
\end{ansblock}

\begin{ansblock}[2章-拓-课时12-40]
% ans:2章-拓-课时12-40
\ansitem{40}{\((1) x^{2}/4+y^{2}=1\)；\((2)\) 证明见详解，定直线为 \(x=4\)}
\ifansbookdetail
\ansnote{详解}{\((1) 2b=2 \Rightarrow  b=1\)；\(e=c/a=\sqrt{3}/2\) 且 \(a^{2}=b^{2}+c^{2} \Rightarrow  a^{2}=4\)、\(c=\sqrt{3}\)，\(E\)：\(x^{2}/4+y^{2}=1\)． \((2) A(-2,0)\)、\(B(2,0)\)，设 \(C(x_{1},y_{1})\)、\(D(x_{2},y_{2})\)．过 \(T\) 的直线不能为 \(x\) 轴（否则 \(C\)、\(D\) 即两顶点，\(AC\)、\(BD\) 均退化为 \(x\) 轴），故设为 \(x=my+1\)（\(m=0\) 时为直线 \(x=1\)）． 代入 \(x^{2}+4y^{2}=4\)：\((m^{2}+4)y^{2}+2my-3=0\)，\(\Delta =4m^{2}+12(m^{2}+4)>0\) 恒成立，\(y_{1}+y_{2}=-2m/(m^{2}+4)\)，\(y_{1}y_{2}=-3/(m^{2}+4)\)① 直线 \(AC\)：\(y=y_{1}/(x_{1}+2)\cdot (x+2)\)；直线 \(BD\)：\(y=y_{2}/(x_{2}-2)\cdot (x-2)\)。设交点为 \(Q(x,y)\)，两式相除得 \((x+2)/(x-2)=y_{2}(x_{1}+2)/[y_{1}(x_{2}-2)]=y_{2}(my_{1}+3)/[y_{1}(my_{2}-1)]=(my_{1}y_{2}+3y_{2})/(my_{1}y_{2}-y_{1})\)． 由①：\(my_{1}y_{2}=-3m/(m^{2}+4)=(3/2)\cdot (-2m/(m^{2}+4))=(3/2)(y_{1}+y_{2})\)， 故分子\(=(3/2)(y_{1}+y_{2})+3y_{2}=(3/2)(y_{1}+3y_{2})\)，分母\(=(3/2)(y_{1}+y_{2})-y_{1}=(1/2)(y_{1}+3y_{2})\)， （\(y_{1}+3y_{2}=0\) 不可能：若 \(y_{1}=-3y_{2}\)，则 \(y_{1}+y_{2}=-2y_{2}=-2m/(m^{2}+4)\)、\(y_{1}y_{2}=-3y_{2}^{2}=-3/(m^{2}+4)\)，两式相除得 \(y_{2}=m\)，代回得 \(m^{2}=m^{2}+4\)，矛盾 ） 所以 \((x+2)/(x-2)=3 \Rightarrow  x+2=3x-6 \Rightarrow  x=4\)， 即直线 \(AC\) 与 \(BD\) 的交点横坐标恒为 \(4\)，交点恒在定直线 \(x=4\) 上．}
\fi
\end{ansblock}

\begin{ansblock}[2章-拓-课时12-41]
% ans:2章-拓-课时12-41
\ansitem{41}{\((1) x^{2}/4+y^{2}/3=1\)；\((2)\) 证明见详解，定直线为 \(x=4\)}
\ifansbookdetail
\ansnote{详解}{\((1) t=-\sqrt{3}/3\) 时 \(l\)：\(x=-(\sqrt{3}/3)y+1\)，令 \(x=0\) 得 \(y=\sqrt{3}\)，即上顶点为 \((0,\sqrt{3})\)，故 \(b=\sqrt{3}\)； \(A\) 为上顶点时 \(\triangle AF_{1}F_{2}\) 的周长\(=|AF_{1}|+|AF_{2}|+|F_{1}F_{2}|=2a+2c=6\)，即 \(a+c=3\)；又 \(a^{2}=b^{2}+c^{2}=3+c^{2}\)， 联立 \(a=3-c\)：\((3-c)^{2}=3+c^{2} \Rightarrow  9-6c=3 \Rightarrow  c=1\)、\(a=2\)，\(C\)：\(x^{2}/4+y^{2}/3=1\)． \((2)\) 直线 \(l\)：\(x=ty+1\) 恒过 \((1,0)=F_{2}\)，且不为 \(x\) 轴（\(x\) 轴与椭圆交于两顶点，此时 \(AM\)、\(BN\) 退化），设 \(A(x_{1},y_{1})\)、\(B(x_{2},y_{2})\)， 代入 \(3x^{2}+4y^{2}=12\)：\((3t^{2}+4)y^{2}+6ty-9=0\)（\(\Delta =36t^{2}+36(3t^{2}+4)>0\) 恒成立），\(y_{1}+y_{2}=-6t/(3t^{2}+4)\)，\(y_{1}y_{2}=-9/(3t^{2}+4)\)． \(M(-2,0)\)、\(N(2,0)\)，直线 \(AM\)：\(y=y_{1}/(x_{1}+2)\cdot (x+2)\)，直线 \(BN\)：\(y=y_{2}/(x_{2}-2)\cdot (x-2)\)，交点 \(Q(x,y)\) 满足 \((x+2)/(x-2)=y_{2}(x_{1}+2)/[y_{1}(x_{2}-2)]=y_{2}(ty_{1}+3)/[y_{1}(ty_{2}-1)]=(ty_{1}y_{2}+3y_{2})/(ty_{1}y_{2}-y_{1})\)， 由 \(ty_{1}y_{2}=-9t/(3t^{2}+4)=(3/2)(y_{1}+y_{2})\)，得分子\(=(3/2)(y_{1}+y_{2})+3y_{2}=(3/2)(y_{1}+3y_{2})\)、分母\(=(3/2)(y_{1}+y_{2})-y_{1}=(1/2)(y_{1}+3y_{2})\)， （\(y_{1}+3y_{2}=0\) 时同 \(12-\)拓\(40\) 由 \(y_{1}y_{2}\) 与 \(y_{1}+y_{2}\) 之比会推出矛盾，此处 \(3t^{2}=3t^{2}+4\) 不可能 ） 故 \((x+2)/(x-2)=3\)，解得 \(x=4\)，即点 \(Q\) 恒在定直线 \(x=4\) 上．}
\fi
\end{ansblock}

\begin{ansblock}[2章-拓-课时12-42]
% ans:2章-拓-课时12-42
\ansitem{42}{\((1) x^{2}/4+y^{2}/3=1\)；\((2)\) 证明见详解，定直线为 \(x=1\)}
\ifansbookdetail
\ansnote{详解}{\((1)\) 设半焦距为 \(c\)，\(A(0,b)\)、\(F_{2}(c,0)\)、\(F_{1}(-c,0)\)、\(Q(q,0)\)（\(q<0\)）． \(AQ\perp AF_{2} \Rightarrow \) 向量\(AQ\cdot \)向量\(AF_{2}=(q,-b)\cdot (c,-b)=qc+b^{2}=0 \Rightarrow  q=-b^{2}/c\)； \(F_{1}\) 为 \(QF_{2}\) 的中点 \(\Rightarrow  (q+c)/2=-c \Rightarrow  q=-3c\)，故 \(b^{2}=3c^{2}\)，\(a^{2}=b^{2}+c^{2}=4c^{2}\)，即 \(a=2c\)，且 \(|QF_{2}|=q\) 到 \(c\) 的距离\(=4c\)， \(\triangle QAF_{2}\) 中 \(\angle QAF_{2}=90^{\circ}\)，故过 \(A\)、\(Q\)、\(F_{2}\) 三点的圆以 \(QF_{2}\) 为直径，圆心为 \(F_{1}(-c,0)\)、半径 \(2c\)； 该圆与 \(l\)：\(x-\sqrt{3}y-3=0\) 相切 \(\Rightarrow  |-c-3|/\sqrt{1+3}=2c \Rightarrow  (c+3)/2=2c\)（\(c>0\)）\(\Rightarrow  c=1\)， 于是 \(a=2\)、\(b^{2}=3\)，\(C\)：\(x^{2}/4+y^{2}/3=1\)． \((2)\) 直线 \(m\) 过 \(P(4,0)\) 且与椭圆交于两点，不能为 \(x\) 轴（否则 \(G\)、\(H\) 即长轴两端点，四边形退化），设 \(m\)：\(x=my+4\)， 代入 \(3x^{2}+4y^{2}=12\)：\((3m^{2}+4)y^{2}+24my+36=0\)，须 \(\Delta =576m^{2}-144(3m^{2}+4)=144(m^{2}-4)>0\)，即 \(m^{2}>4\)， \(y_{1}+y_{2}=-24m/(3m^{2}+4)\)，\(y_{1}y_{2}=36/(3m^{2}+4)\)，由此得 \(m y_{1}y_{2}=-(3/2)(y_{1}+y_{2})\)． 设 \(H(x_{1},y_{1})\)、\(G(x_{2},y_{2})\)，\(M(-2,0)\)、\(N(2,0)\)，四边形 \(MNHG\) 的两条对角线为 \(MH\) 与 \(NG\)： \(MH\)：\(y=y_{1}/(x_{1}+2)\cdot (x+2)\)，\(NG\)：\(y=y_{2}/(x_{2}-2)\cdot (x-2)\)，交点横坐标 \(x=2\cdot [y_{1}(x_{2}-2)+y_{2}(x_{1}+2)]/[y_{2}(x_{1}+2)-y_{1}(x_{2}-2)]\)， 由 \(x_{1}+2=my_{1}+6\)、\(x_{2}-2=my_{2}+2\)，分子方括号\(=y_{1}(my_{2}+2)+y_{2}(my_{1}+6)=2my_{1}y_{2}+2y_{1}+6y_{2}\)， 分母方括号\(=y_{2}(my_{1}+6)-y_{1}(my_{2}+2)=6y_{2}-2y_{1}=2(3y_{2}-y_{1})\)， 故 \(x=2\cdot (2my_{1}y_{2}+2y_{1}+6y_{2})/[2(3y_{2}-y_{1})]=2\cdot (my_{1}y_{2}+y_{1}+3y_{2})/(3y_{2}-y_{1})=2\cdot [-(3/2)(y_{1}+y_{2})+y_{1}+3y_{2}]/(3y_{2}-y_{1})=2\cdot [(3/2)y_{2}-(1/2)y_{1}]/(3y_{2}-y_{1})=2\cdot (1/2)=1\)． （分母 \(3y_{2}-y_{1}=0\) 不可能：若 \(y_{1}=3y_{2}\)，由 \(y_{1}+y_{2}=4y_{2}=-24m/(3m^{2}+4)\) 得 \(y_{2}=-6m/(3m^{2}+4)\)，而 \(y_{1}y_{2}=3y_{2}^{2}=36/(3m^{2}+4)\) 得 \(y_{2}^{2}=12/(3m^{2}+4)\)，两式平方相较得 \(36m^{2}=12(3m^{2}+4)\)，即 \(0=48\)，矛盾，故两对角线恒相交、交点总存在 ） 所以对角线交点的横坐标恒为 \(1\)，即四边形 \(MNHG\) 的对角线交点在定直线 \(x=1\) 上．}
\fi
\end{ansblock}

\begin{ansblock}[2章-拓-课时13-01]
% ans:2章-拓-课时13-01
\ansitem{01}{见详解（按半径大小分三类：等半径\(\rightarrow \)线段 \(O_{1}O_{2}\) 的垂直平分线；不等\(\rightarrow \)以 \(O_{1}\)、\(O_{2}\) 为焦点的双曲线的一支）}
\ifansbookdetail
\ansnote{详解}{设 \(\odot O_{1}\)、\(\odot O_{2}\) 的半径为 \(R_{1}\)、\(R_{2}\)，动圆 \(M\) 的半径为 \(r\)． 动圆与两定圆都外切 \(\Rightarrow  |MO_{1}|=r+R_{1}\)、\(|MO_{2}|=r+R_{2}\)，两式相减消去动半径 \(r\)： 若 \(R_{1}=R_{2}\)，则 \(|MO_{1}|=|MO_{2}|\)，\(M\) 到两定点等距，轨迹是线段 \(O_{1}O_{2}\) 的垂直平分线； 若 \(R_{1}>R_{2}\)，则 \(|MO_{1}|-|MO_{2}|=R_{1}-R_{2}>0\)，且 \(R_{1}-R_{2}<R_{1}+R_{2}<|O_{1}O_{2}|\)（外离条件）， 故 \(M\) 的轨迹是以 \(O_{1}\)、\(O_{2}\) 为焦点的双曲线中靠近 \(O_{2}\) 的一支（即远离大圆一侧的支），因差恒为正，只含一支； 若 \(R_{2}>R_{1}\)，同理 \(|MO_{2}|-|MO_{1}|=R_{2}-R_{1}\in (0,|O_{1}O_{2}|)\)，轨迹是以 \(O_{1}\)、\(O_{2}\) 为焦点的双曲线的另一支． （三支并论：后两类合起来即\("\)以 \(O_{1}\)、\(O_{2}\) 为焦点的双曲线\("\)，但每种半径关系下只取其中一支，故本题须按 \(R_{1}\)、\(R_{2}\) 大小分类作答，不能直接答\("\)双曲线\("\)。）}
\fi
\end{ansblock}

\begin{ansblock}[2章-拓-课时13-02]
% ans:2章-拓-课时13-02
\ansitem{02}{\(B\)}
\ifansbookdetail
\ansnote{详解}{设公共半焦距为 \(c\)（焦距 \(2c\)），\(|F_{1}F_{2}|=2c\)． \(\triangle PF_{1}F_{2}\) 以 \(PF_{1}\) 为底边，故两腰相等：\(|PF_{2}|=|F_{1}F_{2}|=2c\)． \(P\) 在椭圆上：\(|PF_{1}|+|PF_{2}|=2a_{1} \Rightarrow  |PF_{1}|=2a_{1}-2c\)； \(P\) 又在第一象限（位于双曲线右支，靠 \(F_{2}\) 一侧）：\(|PF_{1}|-|PF_{2}|=2a_{2} \Rightarrow  |PF_{1}|=2a_{2}+2c\)， 两式相等：\(2a_{1}-2c=2a_{2}+2c \Rightarrow  a_{1}-a_{2}=2c\)， 两边同除 \(c\)：\(a_{1}/c - a_{2}/c=2\)，而 \(a_{1}/c=1/e_{1}\)、\(a_{2}/c=1/e_{2}\)，故 \(1/e_{1} - 1/e_{2}=2\)． （相容性核验：取 \(e_{2}=2\)（双曲线 \(e_{2}>1\) ），则 \(1/e_{1}=2+1/2=5/2\)、\(e_{1}=2/5<1\) ；此时 \(a_{1}=5c/2\)、\(a_{2}=c/2\)，\(|PF_{1}|=2a_{2}+2c=3c\)、\(|PF_{2}|=2c\)、\(|F_{1}F_{2}|=2c\)，三边 \(2c\)、\(2c\)、\(3c\) 满足三角不等式  等腰三角形存在．）故选：\(B\)．}
\fi
\end{ansblock}

\begin{ansblock}[2章-拓-课时13-03]
% ans:2章-拓-课时13-03
\ansitem{03}{\(D\)}
\ifansbookdetail
\ansnote{详解}{由 向量\(AB=\lambda \)向量\(OB\) 得 向量\(OB-\)向量\(OA=\lambda \)向量\(OB\)，即 向量\(OA=(1-\lambda )\)向量\(OB\)，故 \(O\)、\(A\)、\(B\) 三点共线，直线 \(l\) 过原点 \(O\)． 双曲线 \(x^{2}-y^{2}/2=1\) 关于原点中心对称，过原点的直线与它交于 \(A\)、\(B\) 两点时，\(O\) 恰为弦 \(AB\) 的中点． 对任意点 \(M\)，由中线（阿波罗尼奥斯）恒等式（或平行四边形恒等式）：\(|MA|^{2}+|MB|^{2}=2|MO|^{2}+2|AO|^{2}=2|MO|^{2}+(1/2)|AB|^{2}\)． 两项可各自取最小： ① \(|MO|\) 的最小值为原点 \(O\) 到直线 \(x-\sqrt{2}y-3=0\) 的距离 \(d=|-3|/\sqrt{1+2}=3/\sqrt{3}=\sqrt{3}\)，此时 \(2|MO|^{2}=2\times 3=6\)； ② \(|AB|=2|AO|\)，而 \(A\) 在双曲线上、\(O\) 为对称中心，\(|AO|\) 的最小值为实半轴 \(a=1\)（\(A\) 取顶点 \((\pm 1,0)\)，此时直线 \(l\) 为 \(x\) 轴，与双曲线确实交于 \((\pm 1,0)\) 两点 ），故 \(|AB|min=2\)，\((1/2)|AB|^{2}=(1/2)\times 4=2\)； 两最小可同时达到（\(M\) 取 \(O\) 到已知直线的垂足，\(l\) 取 \(x\) 轴，二者互不牵制）， 所以 \(|MA|^{2}+|MB|^{2}\) 的最小值为 \(6+2=8\)．故选：\(D\)．}
\fi
\end{ansblock}

\begin{ansblock}[2章-拓-课时13-04]
% ans:2章-拓-课时13-04
\ansitem{04}{\(C\)}
\ifansbookdetail
\ansnote{详解}{\(a^{2}=9 \Rightarrow  a=3\)，设右焦点为 \(F_{2}\)． 双曲线关于 \(y\) 轴对称，且 \(F_{1}\) 与 \(F_{2}\) 关于 \(y\) 轴对称，故把 \(P_{1}\) 关于 \(y\) 轴反射到 \(P_{2}\) 时，其到左焦点的距离变为到右焦点的距离：\(|P_{1}F_{1}|=|P_{2}F_{2}|\)． 于是 \(|P_{2}F_{1}|-|P_{1}F_{1}|=|P_{2}F_{1}|-|P_{2}F_{2}|\)， 由图（源图 \(P_{2}\) 在右支上、\(P_{1}\) 在左支上），\(P_{2}\) 在右支 \(\Rightarrow  |P_{2}F_{1}|-|P_{2}F_{2}|=2a=6\)， 故所求值为 \(6\)．故选：\(C\)． （若图中 \(P_{2}\) 落在左支，则差为 \(-2a=-6\)；本题选项只含正值，且源图 \(P_{2}\) 标在右支，故取 \(6\)。）}
\fi
\end{ansblock}

\begin{ansblock}[2章-拓-课时13-05]
% ans:2章-拓-课时13-05
\ansitem{05}{\(BD\)}
\ifansbookdetail
\ansnote{详解}{\(A\)：\(1<t<5\) 时两分母同正，但当 \(t=3\) 时 \(5-t=t-1=2\)，方程化为 \(x^{2}+y^{2}=2\)，是圆不是椭圆，\(A\) 错误． \(B\)：\(t<1\) 时 \(5-t>0\)、\(t-1<0\)，方程即 \(x^{2}/(5-t) - y^{2}/(1-t)=1\)，两分母均为正，表示焦点在 \(x\) 轴上的双曲线，\(B\) 正确． \(C\)：\(C\) 为双曲线 \(\Leftrightarrow  (5-t)(t-1)<0 \Leftrightarrow  t<1\) 或 \(t>5\)；此时实、虚半轴平方为 \(|5-t|\) 与 \(|t-1|\)，半焦距平方 \(c^{2}=|5-t|+|t-1|\)，随 \(t\) 变化：\(t<1\) 时 \(c^{2}=6-2t\)、\(t>5\) 时 \(c^{2}=2t-6\)，均非定值。反例 \(t=0\)：\(x^{2}/5 - y^{2}/1=1\)，\(c^{2}=6\)、焦距 \(2c=2\sqrt{6}\neq 4\)，\(C\) 错误． \(D\)：焦点在 \(y\) 轴上的椭圆 \(\Leftrightarrow  t-1>5-t>0 \Leftrightarrow  t>3\) 且 \(t<5 \Leftrightarrow  3<t<5\)，\(D\) 正确． 故选：\(BD\)．}
\fi
\end{ansblock}

\begin{ansblock}[2章-拓-课时13-06]
% ans:2章-拓-课时13-06
\ansitem{06}{\(m=0\)：直线 \(x=0\)；\(0<m<6\)：双曲线右支，\(x^{2}/(m^{2}/4) - y^{2}/(9-m^{2}/4)=1 (x\geqslant m/2)\)；\(m=6\)：射线 \(y=0 (x\geqslant 3)\)；\(m>6\)：点 \(P\) 不存在}
\ifansbookdetail
\ansnote{详解}{\(|F_{1}F_{2}|=2c=6\)． \((1) m=0\)：\(|PF_{1}|=|PF_{2}|\)，\(P\) 在 \(F_{1}F_{2}\) 的垂直平分线上，轨迹是直线，方程 \(x=0\)； \((2) 0<m<6\)：\(|PF_{1}|-|PF_{2}|=m<|F_{1}F_{2}|\) 且差为正，故 \(P\) 到 \(F_{2}\) 较近，轨迹是以 \(F_{1}\)、\(F_{2}\) 为焦点的双曲线中靠近 \(F_{2}\) 的一支，即右支； 实半轴 \(a=m/2\)、\(c=3\)、\(b^{2}=c^{2}-a^{2}=9-m^{2}/4\)（\(0<m<6\) 保证 \(b^{2}>0\) ），轨迹方程 \(x^{2}/(m^{2}/4) - y^{2}/(9-m^{2}/4)=1 (x\geqslant m/2)\)； \((3) m=6=|F_{1}F_{2}|\)：\(\|PF_{1}|-|PF_{2}|\)｜\(=|F_{1}F_{2}|\) 须 \(P\)、\(F_{1}\)、\(F_{2}\) 共线且 \(P\) 不在两焦点之间，再由 \(|PF_{1}|>|PF_{2}|\) 定 \(P\) 在 \(F_{2}\) 外侧，轨迹是射线，方程 \(y=0 (x\geqslant 3)\)； \((4) m>6\)：与 \(\|PF_{1}|-|PF_{2}|\)｜\(\leqslant |F_{1}F_{2}|=6\) 矛盾，满足条件的点 \(P\) 不存在．}
\fi
\end{ansblock}

\par\Needspace*{3\baselineskip}\noindent\biaoqian{【下册】}\par

\begin{ansblock}[2章-拓-课时14-01]
% ans:2章-拓-课时14-01
\ansitem{01}{\(B\)}
\ifansbookdetail
\ansnote{详解}{\(PQ\) 为通径，\(P(c,\pm b^{2}/a)\)，\(|PQ|=2b^{2}/a\)。 \(\triangle PF_{1}Q\) 中 \(\angle PF_{1}Q=90^{\circ}\)，\(F_{2}\) 为斜边 \(PQ\) 的中点，由直角三角形斜边中线定理：\(|F_{1}F_{2}|=|PQ|/2\)，即 \(2c=b^{2}/a\)，\(2ac=b^{2}=c^{2}-a^{2}\)。 两边除以 \(a^{2}\)：\(e^{2}-2e-1=0\)，\(e>1\)，\(e=1+\sqrt{2}\)，故选\(B\)。 验算：\(e=1+\sqrt{2}\) 时 \(e^{2}=3+2\sqrt{2}\)，\(e^{2}-2e-1=(3+2\sqrt{2})-2(1+\sqrt{2})-1=0\) 。}
\fi
\end{ansblock}

\begin{ansblock}[2章-拓-课时14-02]
% ans:2章-拓-课时14-02
\ansitem{02}{\(2\)}
\ifansbookdetail
\ansnote{详解}{\(A\) 在线段 \(F_{1}B\) 上且 \(|F_{1}A|=|AB|\)，故 \(A\) 为 \(F_{1}B\) 的中点；\(O\) 为 \(F_{1}F_{2}\) 的中点，故 \(OA\) 是 \(\triangle F_{1}F_{2}B\) 的中位线：\(OA\parallel F_{2}B\)，\(|F_{2}B|=2|OA|\)。 由 \(F_{1}B\cdot F_{2}B=0\) 得 \(F_{1}B\perp F_{2}B\)，结合 \(OA\parallel F_{2}B\) 得 \(OA\perp F_{1}B\)。 于是 \(\triangle F_{1}OB\) 中，\(OA\) 既是边 \(F_{1}B\) 上的中线又是高，故 \(\triangle F_{1}OB\) 等腰：\(|OB|=|OF_{1}|=c\)，且 \(\angle AOF_{1}=\angle AOB\)。 又 \(A\)、\(B\) 分别在两条渐近线上，两渐近线关于 \(x\) 轴对称，故 \(\angle BOF_{2}=\angle AOF_{1}\)；而 \(\angle AOF_{1}+\angle AOB+\angle BOF_{2}=180^{\circ}\)，三者在上述等量关系下相等，各为 \(60^{\circ}\)。 故 \(b/a=tan60^{\circ}=\sqrt{3}\)，\(e=\sqrt{}(1+(b/a)^{2})=\sqrt{1+3}=2\)。 验算：三内角和 \(3\times 60^{\circ}=180^{\circ}\) ；\(b/a=\sqrt{3}\) 时 \(e=2\) 。}
\fi
\end{ansblock}

\begin{ansblock}[2章-拓-课时14-03]
% ans:2章-拓-课时14-03
\ansitem{03}{\(C\)}
\ifansbookdetail
\ansnote{详解}{六个顶点内接于以 \(F_{1}F_{2}\) 为直径的圆，由直径所对圆周角 \(\angle F_{1}PF_{2}=90^{\circ}\)；正六边形各边对圆心角 \(60^{\circ}\)，得 \(\angle PF_{2}F_{1}=30^{\circ}\)。 解三角形：\(|PF_{2}|=2c\cdot cos30^{\circ}=\sqrt{3}c\)，\(|PF_{1}|=2c\cdot sin30^{\circ}=c\)。 \(P\) 在双曲线上，\(|PF_{2}|-|PF_{1}|=2a \rightarrow  \sqrt{3}c-c=2a \rightarrow  e=c/a=2/(\sqrt{3}-1)=\sqrt{3}+1\)，故选\(C\)。 验算：\(2/(\sqrt{3}-1)\) 有理化 \(=2(\sqrt{3}+1)/2=\sqrt{3}+1\) ；\(P\) 在右支，长焦半径减短焦半径 。 （卷③自带详解照录，验算一致。）}
\fi
\end{ansblock}

\begin{ansblock}[2章-拓-课时14-04]
% ans:2章-拓-课时14-04
\ansitem{04}{\(D\)}
\ifansbookdetail
\ansnote{详解}{圆 \(x^{2}+y^{2}=a^{2}+b^{2}\) 的半径为 \(c\)，其直径恰为 \(F_{1}F_{2}\)，故 \(P\) 在圆上 \(\Rightarrow  \angle F_{1}PF_{2}=90^{\circ}\)（直径所对圆周角）。 设 \(|PF_{1}|=n\)，由定义 \(|PF_{2}|=2a+n\)。\(Rt\triangle PF_{1}F_{2}\) 中：\(n^{2}+(2a+n)^{2}=4c^{2} \rightarrow  2n^{2}+4an+4a^{2}=4(a^{2}+b^{2}) \rightarrow  n^{2}=2b^{2}-2an\)。 由 \(F_{1}P=(1/2)F_{2}Q\)（同向向量）得 \(|F_{2}Q|=2n\) 且 \(F_{1}P\parallel F_{2}Q\)，故 \(\angle QF_{2}F_{1}=180^{\circ}-\angle PF_{1}F_{2}\)，\(cos\angle QF_{2}F_{1}=-n/(2c)\)（\(Rt\triangle \) 中 \(cos\angle PF_{1}F_{2}=n/(2c)\)）。 由定义 \(|F_{1}Q|=2a+2n\)。\(\triangle F_{1}F_{2}Q\) 中余弦定理： \((2a+2n)^{2}=(2n)^{2}+(2c)^{2}+4n^{2} \rightarrow  4a^{2}+8an+4n^{2}=4n^{2}+4c^{2}+4n^{2} \rightarrow  4a^{2}+8an=4c^{2}+4n^{2} \rightarrow  a^{2}+2an=c^{2}+n^{2}=(a^{2}+b^{2})+n^{2} \rightarrow  2an=b^{2}+n^{2}\)，即 \(b^{2}=2an-n^{2}\)。 与 \(n^{2}=2b^{2}-2an\) 联立：\(n^{2}=2(2an-n^{2})-2an=2an-2n^{2} \rightarrow  3n^{2}=2an \rightarrow  a=(3/2)n\)； \(b^{2}=2an-n^{2}=3n^{2}-n^{2}=2n^{2}\)；\(c^{2}=a^{2}+b^{2}=(9/4+2)n^{2}=(17/4)n^{2}\)。 \(e=c/a=[(\sqrt{17}/2)n]/[(3/2)n]=\sqrt{17}/3\)，故选\(D\)。 验算：\(n=1\) 时 \(a=1.5\)、\(b=\sqrt{2}\)、\(c=\sqrt{17}/2\approx 2.06\)，\(c^{2}=a^{2}+b^{2}=2.25+2=4.25\) 。 （卷③自带详解照录校订，读数一致。）}
\fi
\end{ansblock}

\begin{ansblock}[2章-拓-课时14-05]
% ans:2章-拓-课时14-05
\ansitem{05}{\(\sqrt{3}\)}
\ifansbookdetail
\ansnote{详解}{\(O(0,0)\)，\(A(-a,0)\)，\(B(a,0)\)，圆 \(P\) 的圆心为 \(P(0,a)\)，半径 \(\sqrt{2}a\)。 \(|OA|=|OB|=|OP|=a\)，故 \(O\)、\(A\)、\(P\)、\(B\) 四点共圆（圆心 \(O\)、半径 \(a\)），其中 \(AB\) 为该圆的直径（\(A\)、\(B\) 关于 \(O\) 对称），\(P\) 在圆上，由直径所对圆周角得 \(\angle APB=90^{\circ}\)。 \(M\) 在圆上且在第一象限，位于弦 \(AB\) 所对优弧，故圆周角 \(\angle AMB=(1/2)\angle APB=45^{\circ}\)。 夹角公式：\(tan\angle AMB=(n-m)/(1+nm)=1\)，\(n-m=3 \rightarrow  nm=2\)。 设 \(M(x,y)\)：\(m=y/(x+a)\)，\(n=y/(x-a)\)，\(nm=y^{2}/(x^{2}-a^{2})\)；由 \(y^{2}=b^{2}(x^{2}/a^{2}-1)\) 得 \(nm=b^{2}/a^{2}=2\)。 \(e^{2}=1+b^{2}/a^{2}=3\)，\(e=\sqrt{3}\)。 验算：\(mn=(y^{2})/((x-a)(x+a))\)，\((x-a)(x+a)=x^{2}-a^{2}=a^{2}(y^{2}/b^{2})-a^{2}=(a^{2}/b^{2})y^{2} \rightarrow  mn=b^{2}/a^{2}\) 。 （卷③自带详解照录校订：\(|OA|=|OB|=|OP|\) 与圆心角推导等价，此处按向量验。）}
\fi
\end{ansblock}

\begin{ansblock}[2章-拓-课时14-06]
% ans:2章-拓-课时14-06
\ansitem{06}{\(2\sqrt{3}/3\)}
\ifansbookdetail
\ansnote{详解}{设 \(A(a,0)\)，\(B(-a,0)\)，\(P(c,y)\)（\(l\)：\(x=c\)）。\(k\_PA=y/(c-a)\)，\(k\_PB=y/(c+a)\)。 由夹角公式 \(tan\angle APB=|[k\_PA-k\_PB]/[1+k\_PA\cdot k\_PB]|=\sqrt{3}\)： \(|(2ay)/(c^{2}-a^{2})| / |1+y^{2}/(c^{2}-a^{2})|=\sqrt{3} \rightarrow \) 整理为关于 \(y\) 的方程 \(\sqrt{3}y^{2}-2ay+\sqrt{3}(c^{2}-a^{2})=0\) 有实根。 判别式 \(\Delta =4a^{2}-12(c^{2}-a^{2})=4(4a^{2}-3c^{2})\geqslant 0 \rightarrow  4a^{2}\geqslant 3c^{2} \rightarrow  e^{2}\leqslant 4/3\)，\(e\leqslant 2\sqrt{3}/3\)。 离心率最大值 \(2\sqrt{3}/3\)。 验算：\(e=2\sqrt{3}/3\) 时 \(\Delta =0\)，\(P\) 取对称轴位置 \(y=a/\sqrt{3}\) 。 （卷③自带详解照录，验算一致。）}
\fi
\end{ansblock}

\begin{ansblock}[2章-拓-课时14-07]
% ans:2章-拓-课时14-07
\ansitem{07}{\(B\)}
\ifansbookdetail
\ansnote{详解}{共渐近线设 \(x^{2}-2y^{2}=m\)（\(m\neq 0\)）。焦点 \((0,-6)\) 在 \(y\) 轴上，故焦点在 \(y\) 轴，\(m<0\)。 标准形 \(y^{2}/(-m/2)-x^{2}/(-m)=1\)：\(a^{2}=-m/2\)，\(b^{2}=-m\)，\(c^{2}=a^{2}+b^{2}=-3m/2=36 \rightarrow  m=-24\)。 方程 \(y^{2}/12-x^{2}/24=1\)，选\(B\)。 验算：\(c^{2}=12+24=36\)，\(c=\pm 6\) ，焦点含 \((0,-6)\) 。 （卷③自带详解照录，验算一致。）}
\fi
\end{ansblock}

\begin{ansblock}[2章-拓-课时14-08]
% ans:2章-拓-课时14-08
\ansitem{08}{证明见解析。}
\ifansbookdetail
\ansnote{详解}{证明：设 \(C\)：\(x^{2}/m^{2}-y^{2}/m^{2}=1\)（\(m>0\)），则 \(F_{1}(-\sqrt{2}m,0)\)，\(F_{2}(\sqrt{2}m,0)\)。设 \(P(x,y)\)，满足 \(x^{2}-y^{2}=m^{2}\)。 \(|PO|^{2}=x^{2}+y^{2}\)。 \(|PF_{1}|\cdot |PF_{2}|=\sqrt{(x+\sqrt\{2\}m)^{2}+y^{2}}\cdot \sqrt{(x-\sqrt\{2\}m)^{2}+y^{2}} \allowbreak =\sqrt{}\{[(x^{2}+2m^{2}+y^{2})+2\sqrt{2}mx]\cdot [(x^{2}+2m^{2}+y^{2})-2\sqrt{2}mx]\} \allowbreak =\sqrt{(x^{2}+2m^{2}+y^{2})^{2}-8m^{2}x^{2}}\)。 根号内用 \(x^{2}-y^{2}=m^{2}\)（即 \(y^{2}=x^{2}-m^{2}\)）化简： \(x^{2}+2m^{2}+y^{2}=2x^{2}+m^{2}\)，\((2x^{2}+m^{2})^{2}-8m^{2}x^{2}=4x^{4}+4m^{2}x^{2}+m^{4}-8m^{2}x^{2}=4x^{4}-4m^{2}x^{2}+m^{4}=(2x^{2}-m^{2})^{2}\)。 故 \(|PF_{1}|\cdot |PF_{2}|=|2x^{2}-m^{2}|=2x^{2}-m^{2}\)（\(P\) 在双曲线上 \(|x|\geqslant m\)，\(2x^{2}-m^{2}\geqslant m^{2}>0\)）\(=2x^{2}-(x^{2}-y^{2})=x^{2}+y^{2}=|PO|^{2}\)。证毕。 验算：\(m=1\)，\(P(\sqrt{2},1)\)：\(|PO|^{2}=3\)；\(|PF_{1}|=\sqrt{(\sqrt\{2\}+\sqrt\{2\})^{2}+1}=\sqrt{9}=3\)，\(|PF_{2}|=\sqrt{(\sqrt\{2\}-\sqrt\{2\})^{2}+1}=1\)，积\(=3\) 。 （卷③自带详解照录校订：源详解根号内配方路径等价，此处以 \(y^{2}=x^{2}-m^{2}\) 直接配方；关键步 \(2m^{2}=x^{2}-y^{2}\) 使 \(x^{2}+2m^{2}+y^{2}=3x^{2}-y^{2}\)，积\(=(x^{2}+y^{2})^{2}\) 与源一致。）}
\fi
\end{ansblock}

\begin{ansblock}[2章-拓-课时14-09]
% ans:2章-拓-课时14-09
\ansitem{09}{\(18\)}
\ifansbookdetail
\ansnote{详解}{\(PQ\) 与 \(F_{1}F_{2}\) 互相平分（\(P\)、\(Q\) 关于 \(O\) 对称，\(F_{1}\)、\(F_{2}\) 关于 \(O\) 对称）且 \(|PQ|=|F_{1}F_{2}|\)，对角线相等又互相平分 \(\rightarrow \) 四边形 \(PF_{1}QF_{2}\) 为矩形，其面积\(=|PF_{1}|\cdot |PF_{2}|\)。 矩形中 \(|PF_{1}|^{2}+|PF_{2}|^{2}=|F_{1}F_{2}|^{2}=4c^{2}=100\)；又 \(||PF_{1}|-|PF_{2}||=2a=8\)。 \(2|PF_{1}||PF_{2}|=100-64=36 \rightarrow  |PF_{1}||PF_{2}|=18\)，面积\(18\)。 验算：\(a=4\)，\(b=3\)，\(c=5\)，\(4c^{2}=100\) ；\(2|PF_{1}||PF_{2}|=4c^{2}-(2a)^{2}=100-64=36 \rightarrow  18\) 。 （卷③自带详解照录，验算一致。）}
\fi
\end{ansblock}

\begin{ansblock}[2章-拓-课时14-10]
% ans:2章-拓-课时14-10
\ansitem{10}{\(3+2\sqrt{3}\)}
\ifansbookdetail
\ansnote{详解}{左焦点 \((-\sqrt{a^{2}+1},0)=(-2,0) \rightarrow  a^{2}=3\)。 设 \(P(x_{0},y_{0})\)（右支 \(x_{0}\geqslant \sqrt{3}\)），\(y_{0}^{2}=x_{0}^{2}/3-1\)。 \(OP\cdot FP=(x_{0},y_{0})\cdot (x_{0}+2,y_{0})=x_{0}^{2}+2x_{0}+y_{0}^{2}=(4/3)x_{0}^{2}+2x_{0}-1\)。 在 \([\sqrt{3},+\infty )\) 上导数 \((8/3)x_{0}+2>0\)，单调递增，\(x_{0}=\sqrt{3}\) 时取最小： \((4/3)\times 3+2\sqrt{3}-1=3+2\sqrt{3}\)。 验算：\(4+2\sqrt{3}-1=3+2\sqrt{3}\) ；对称轴 \(x_{0}=-3/4\) 在区间外 。 （卷③自带详解照录，验算一致。）}
\fi
\end{ansblock}

\begin{ansblock}[2章-拓-课时14-11]
% ans:2章-拓-课时14-11
\ansitem{11}{\((1) x^{2}/400-y^{2}/500=1\)（\(x<0\)）；\((2) 20\sqrt{2} km\)。}
\ifansbookdetail
\ansnote{详解}{\((1)\) 以 \(O\) 为原点、\(AB\) 所在直线为 \(x\) 轴建系：\(A(-30,0)\)，\(B(30,0)\)，\(C(0,30)\)。 设观察员位置 \(P(x,y)\)：\(|PB|-|PA|=(40/V_{0})\cdot V_{0}=40<|AB|=60\)，轨迹为双曲线左支。 \(2a=40\)，\(2c=60\)，\(a=20\)，\(c=30\)，\(b^{2}=c^{2}-a^{2}=500\)。 轨迹方程 \(x^{2}/400-y^{2}/500=1\)（\(x<0\)）。 \((2)\) 设轨迹上一点 \(M(x,y)\)：\(|MC|^{2}=x^{2}+(y-30)^{2}=x^{2}+y^{2}-60y+900\)。 由 \(x^{2}=(4/5)y^{2}+400\)：\(|MC|^{2}=(9/5)y^{2}-60y+1300=(9/5)(y-50/3)^{2}+800\geqslant 800\)。 当 \(y=50/3\) 时 \(|MC|min=20\sqrt{2}\)，扫描半径 \(r\) 至少 \(20\sqrt{2} km\)。 验算：\(y=50/3\) 时 \((9/5)y^{2}-60y+1300=(9/5)(2500/9)-1000+1300=500-1000+1300=800\) 。 （卷③自带详解照录，验算一致。）}
\fi
\end{ansblock}

\begin{ansblock}[2章-拓-课时14-12]
% ans:2章-拓-课时14-12
\ansitem{12}{\(CD\)}
\ifansbookdetail
\ansnote{详解}{\(A\)：按定义实虚轴互换，共轭双曲线应为 \(y^{2}/b^{2}-x^{2}/a^{2}=1\)，\(A\) 错； \(B\)：原曲线渐近线 \(y=\pm (b/a)x\)，共轭曲线渐近线 \(y=\pm (b/a)x\)（\(a/b\) 与 \(b/a\) 互换后方程相同），渐近线相同，\(B\) 错； \(C\)：\(e_{1}=c/a\)，\(e_{2}=c/b\)，\(e_{1}e_{2}=c^{2}/(ab)=(a^{2}+b^{2})/(ab)=a/b+b/a\geqslant 2\)（当且仅当 \(a=b\) 取等），\(C\) 对； \(D\)：焦点 \((\pm c,0)\) 与 \((0,\pm c)\) 均在圆 \(x^{2}+y^{2}=c^{2}\) 上，\(D\) 对。故选\(CD\)。 验算：\(C\) 取等条件 \(a=b\)（等轴双曲线 \(e_{1}=e_{2}=\sqrt{2}\)，积\(2\)）；\(D\)：\(c^{2}=a^{2}+b^{2}\) 。 （卷③自带详解照录，验算一致。）}
\fi
\end{ansblock}

\begin{ansblock}[2章-拓-课时14-14]
% ans:2章-拓-课时14-14
\ansitem{14}{实轴长\(2\)；虚轴长\(4\sqrt{6}\)；焦点\((\pm 5,0)\)；渐近线方程 \(y=\pm 2\sqrt{6}x\)。}
\ifansbookdetail
\ansnote{详解}{\(a^{2}=1\)、\(b^{2}=24\)，\(a=1\)、\(b=2\sqrt{6}\)，\(c^{2}=1+24=25\)，\(c=5\)。 实轴长\(2a=2\)；虚轴长\(2b=4\sqrt{6}\)；焦点\((\pm 5,0)\)；渐近线 \(y=\pm (b/a)x=\pm 2\sqrt{6}x\)。 （本块自撰亲算：\(25=1+24\) ；教材页底答案条乱码，读数自验。）}
\fi
\end{ansblock}

\begin{ansblock}[2章-拓-课时14-15]
% ans:2章-拓-课时14-15
\ansitem{15}{\(x^{2}/16-y^{2}/9=1\)；\(e=5/4\)。}
\ifansbookdetail
\ansnote{详解}{焦点 \((5,0)\) 在 \(x\) 轴：\(b/a=3/4\)（渐近线 \(y=\pm (b/a)x\)）。 \(c=5\)，\(a^{2}+b^{2}=25\)，\(b^{2}=(9/16)a^{2} \rightarrow  a^{2}(1+9/16)=25 \rightarrow  a^{2}=16\)，\(b^{2}=9\)。 标准方程 \(x^{2}/16-y^{2}/9=1\)，\(e=c/a=5/4\)。 （本块自撰亲算：\(16+9=25\) ；渐近线 \(y=\pm (3/4)x\) 。）}
\fi
\end{ansblock}

\begin{ansblock}[2章-拓-课时14-16]
% ans:2章-拓-课时14-16
\ansitem{16}{它们的渐近线相同，均为 \(y=\pm (b/a)x\)。}
\ifansbookdetail
\ansnote{详解}{\(\lambda >0\) 时方程即 \(x^{2}/(a^{2}\lambda )-y^{2}/(b^{2}\lambda )=1\)，焦点在 \(x\) 轴，渐近线由 \(x^{2}/(a^{2}\lambda )-y^{2}/(b^{2}\lambda )=0\) 得 \(y=\pm (b/a)x\)； \(\lambda <0\) 时方程即 \(y^{2}/(-b^{2}\lambda )-x^{2}/(-a^{2}\lambda )=1\)，焦点在 \(y\) 轴，渐近线由同式得 \(y=\pm (b/a)x\)。 故所有曲线渐近线相同：\(y=\pm (b/a)x\)。 （本块自撰亲算；渐近线只取决于 \(a:b\) 之比，与 \(\lambda \) 无关。）}
\fi
\end{ansblock}

\begin{ansblock}[2章-拓-课时14-17]
% ans:2章-拓-课时14-17
\ansitem{17}{\(x^{2}/9-y^{2}/27=1\)（焦点在 \(x\) 轴；焦点在 \(y\) 轴时为 \(y^{2}/9-x^{2}/27=1\)）。}
\ifansbookdetail
\ansnote{详解}{顶点距 \(6 \rightarrow  2a=6\)，\(a=3\)。 焦点段 \([-c,c]\) 上的分点 \(-a,0,a\) 将其四等分，等分距 \(a=c/2 \rightarrow  c=6\)。 \(b^{2}=c^{2}-a^{2}=36-9=27\)。 焦点在 \(x\) 轴：\(x^{2}/9-y^{2}/27=1\)。 （本块自撰亲算：题面未指明焦点位置，教材按 \(x\) 轴书写；\(y\) 轴情形同法。验算 \(3\)、\(6\) 间距均为\(3\) 。）}
\fi
\end{ansblock}

\begin{ansblock}[2章-拓-课时14-18]
% ans:2章-拓-课时14-18
\ansitem{18}{曲线是双曲线（两支）：以 \(AB\) 所在直线为 \(x\) 轴、\(AB\) 中点为原点，其方程为 \(x^{2}/9-y^{2}/16=1\)。}
\ifansbookdetail
\ansnote{详解}{曲线上点 \(P\) 满足 \(||PA|-|PB||=6<|AB|=10\)，由定义轨迹为双曲线，\(2a=6\)、\(2c=10\)，\(a=3\)、\(c=5\)，\(b^{2}=25-9=16\)。 以 \(AB\) 中点为原点、\(AB\) 为 \(x\) 轴建系，方程 \(x^{2}/9-y^{2}/16=1\)。 作图观察：交点连成两支光滑曲线，与双曲线吻合。 （本块自撰亲算；作图探究题，图形观察与定义互证。）}
\fi
\end{ansblock}

\begin{ansblock}[2章-拓-课时14-19]
% ans:2章-拓-课时14-19
\ansitem{19}{\(x^{2}/(32/5)-y^{2}/(18/5)=1\)。}
\ifansbookdetail
\ansnote{详解}{椭圆 \(a^{2}=10\)、\(b^{2}=5\)，\(c^{2}=5\)，长轴端点 \((\pm \sqrt{10},0)\)，即双曲线焦点在 \(x\) 轴，\(c=\sqrt{10}\)。 渐近线 \(y=\pm (3/4)x \rightarrow  b/a=3/4\)。 \(a^{2}+b^{2}=c^{2}=10\)，\(a^{2}(1+9/16)=10 \rightarrow  a^{2}\times 25/16=10 \rightarrow  a^{2}=32/5\)，\(b^{2}=10-32/5=18/5\)。 标准方程 \(x^{2}/(32/5)-y^{2}/(18/5)=1\)。 （本块自撰亲算：\(32/5+18/5=10\) ；\(b/a=\sqrt{18/5}\div \sqrt{32/5}=\sqrt{18/32}=3/4\) 。）}
\fi
\end{ansblock}

\begin{ansblock}[2章-拓-课时14-20]
% ans:2章-拓-课时14-20
\ansitem{20}{\(x^{2}/(64/13)-y^{2}/(144/13)=1\)。}
\ifansbookdetail
\ansnote{详解}{焦点在 \(x\) 轴，\(c=4\)；渐近线 \(y=(3/2)x \rightarrow  b/a=3/2\)。 \(a^{2}+b^{2}=16\)，\(b^{2}=(9/4)a^{2} \rightarrow  a^{2}(13/4)=16 \rightarrow  a^{2}=64/13\)，\(b^{2}=16-64/13=144/13\)。 标准方程 \(x^{2}/(64/13)-y^{2}/(144/13)=1\)。 （本块自撰亲算：\(64/13+144/13=208/13=16\) ；\((144/64)=9/4\) 。）}
\fi
\end{ansblock}

\begin{ansblock}[2章-拓-课时14-21]
% ans:2章-拓-课时14-21
\ansitem{21}{\(x^{2}/12-y^{2}/8=1\)。}
\ifansbookdetail
\ansnote{详解}{原双曲线 \(c^{2}=16+4=20\)，公共焦点 \(\rightarrow \) 所求双曲线 \(c^{2}=20\)，焦点在 \(x\) 轴。 半实轴长 \(a=2\sqrt{3}\)，\(a^{2}=12\)；\(b^{2}=c^{2}-a^{2}=20-12=8\)。 标准方程 \(x^{2}/12-y^{2}/8=1\)。 （本块自撰亲算：\(12+8=20\) 。）}
\fi
\end{ansblock}

\begin{ansblock}[2章-拓-课时14-22]
% ans:2章-拓-课时14-22
\ansitem{22}{\(2\)}
\ifansbookdetail
\ansnote{详解}{设 \(P(x,y)\)（\(|x|\geqslant 1\)）：\(|PA|^{2}=(x-3)^{2}+y^{2}=(x-3)^{2}+3(x^{2}-1)=4x^{2}-6x+6\)。 对称轴 \(x=3/4<1\)，故 \(x=1\)（右顶点）时取最小：\(4-6+6=4\)，\(|PA|min=2\)。 （本块自撰亲算：\(y^{2}=3(x^{2}-1)\) 代入消元 ；对称轴在 \(|x|\geqslant 1\) 左侧，端点取值。）}
\fi
\end{ansblock}

\begin{ansblock}[2章-拓-课时14-23]
% ans:2章-拓-课时14-23
\ansitem{23}{\(x^{2}/4-y^{2}/5=1\)。}
\ifansbookdetail
\ansnote{详解}{（直接法，不引第二定义）设 \(M(x,y)\)： \(\sqrt{(x-3)^{2}+y^{2}}=(3/2)|x-4/3|\)。 两边平方：\(x^{2}-6x+9+y^{2}=(9/4)(x^{2}-(8/3)x+16/9)=(9/4)x^{2}-6x+4\)。 \(-6x\) 抵消：\(y^{2}=(9/4)x^{2}+4-x^{2}-9=(5/4)x^{2}-5\)，即 \((5/4)x^{2}-y^{2}=5\)， 标准形 \(x^{2}/4-y^{2}/5=1\)。 验算：\(F(3,0)\) 与准线位置 \(x=a^{2}/c=4/3\)、比值 \(c/a=3/2\) 相容（\(a=2\)，\(c=3\) ）。}
\fi
\end{ansblock}

\begin{ansblock}[2章-拓-课时14-24]
% ans:2章-拓-课时14-24
\ansitem{24}{比值等于 \(c/a\)（离心率 \(e\)）。}
\ifansbookdetail
\ansnote{详解}{（直接法）\(P(x,y)\) 在双曲线上：\(y^{2}=b^{2}(x^{2}/a^{2}-1)\)。 \(d_{1}=|PF|=\sqrt{(x+c)^{2}+y^{2}}=\sqrt{x^{2}+2cx+c^{2}+b^{2}x^{2}/a^{2}-b^{2}}\)，以 \(b^{2}=c^{2}-a^{2}\) 代入： \(=x^{2}+2cx+c^{2}+(c^{2}-a^{2})(x^{2}/a^{2})-(c^{2}-a^{2})=(c^{2}/a^{2})x^{2}+2cx+a^{2}=[(c/a)x+a]^{2}\)。 （\(x\geqslant a\) 时 \((c/a)x+a>0\)；\(x\leqslant -a\) 时亦有 \((c/a)x+a<0\)，平方去根号后 \(|PF|=(c/a)|x+a^{2}/c|\)。） \(d_{2}=|x+a^{2}/c|\)（\(P\) 到 \(l\)：\(x=-a^{2}/c\) 的距离）。 故 \(d_{1}/d_{2}=(c/a)|x+a^{2}/c|/|x+a^{2}/c|=c/a\)。 （本块自撰亲算：配方 \([(c/a)x+a]^{2}=c^{2}x^{2}/a^{2}+2cx+a^{2}\) 展开核对 。）}
\fi
\end{ansblock}

\begin{ansblock}[2章-拓-课时14-25]
% ans:2章-拓-课时14-25
\ansitem{25}{\(16\)}
\ifansbookdetail
\ansnote{详解}{\(a=3\)，\(b=4\)，\(c=5\)。设两焦半径 \(r_{1}=|MF_{1}|\)，\(r_{2}=|MF_{2}|\)（\(r_{1}>r_{2}\)）。 \(MF_{1}\perp MF_{2}\)：\(r_{1}^{2}+r_{2}^{2}=|F_{1}F_{2}|^{2}=100\)；定义：\(r_{1}-r_{2}=2a=6\)。 \((r_{1}-r_{2})^{2}=100-2r_{1}r_{2}=36 \rightarrow  r_{1}r_{2}=32\)。 \(S=(1/2)r_{1}r_{2}\cdot sin90^{\circ}=16\)。 （本块自撰亲算：勾股 \(100-36=64=2r_{1}r_{2} \rightarrow  r_{1}r_{2}=32\) 。）}
\fi
\end{ansblock}

\begin{ansblock}[2章-拓-课时14-26]
% ans:2章-拓-课时14-26
\ansitem{26}{\((1) ax^{2}-y^{2}=36a\)（\(y\neq 0\)）；\((2) a>0\) 时为焦点在 \(x\) 轴的双曲线（除去顶点）；\(a<0\) 时为椭圆：其中 \(a=-1\) 为圆 \(x^{2}+y^{2}=36\)（\(y\neq 0\)），\(-1<a<0\) 为焦点在 \(x\) 轴的椭圆，\(a<-1\) 为焦点在 \(y\) 轴的椭圆（均除去长轴端点——方程中 \(y\neq 0\)）。}
\ifansbookdetail
\ansnote{详解}{\((1)\) 过 \((6,0)\) 斜率 \(k_{1}\)、过 \((-6,0)\) 斜率 \(k_{2}\)，\(k_{1}k_{2}=a\)。 \(M(x,y)\)：\(k_{1}=y/(x-6)\)，\(k_{2}=y/(x+6)\)（\(x\neq \pm 6\)，且 \(M\) 不与 \((\pm 6,0)\) 重合故 \(y\neq 0\)）。 \([y/(x-6)]\cdot [y/(x+6)]=a \rightarrow  y^{2}=a(x^{2}-36) \rightarrow  ax^{2}-y^{2}=36a\)（\(y\neq 0\)）。 \((2) a>0\)：\(x^{2}/36-y^{2}/(36a)=1\)，焦点在 \(x\) 轴的双曲线（除去 \((\pm 6,0)\)）； \(a<0\)：记 \(a=-|a|\)：\(x^{2}/36+y^{2}/(36|a|)=1\)，椭圆； 其中 \(a=-1\)：\(x^{2}+y^{2}=36\)，圆；\(-1<a<0\)：\(36>36|a|\)，焦点在 \(x\) 轴的椭圆；\(a<-1\)：\(36<36|a|\)，焦点在 \(y\) 轴的椭圆（均除去 \(x\) 轴上的点，即 \(y\neq 0\)）。 （本块自撰亲算；分类讨论按 \(|a|\) 与 \(1\) 比较。）}
\fi
\end{ansblock}

\begin{ansblock}[2章-拓-课时14-27]
% ans:2章-拓-课时14-27
\ansitem{27}{\(x^{2}/a^{2}-y^{2}/b^{2}=1\)，其中 \(b^{2}=c^{2}-a^{2}\)。}
\ifansbookdetail
\ansnote{详解}{（直接法）设 \(M(x,y)\)： \(\sqrt{(x+c)^{2}+y^{2}}=(c/a)|x+a^{2}/c|\)。 平方：\(x^{2}+2cx+c^{2}+y^{2}=(c^{2}/a^{2})(x^{2}+2a^{2}x/c+a^{4}/c^{2})=(c^{2}/a^{2})x^{2}+2cx+a^{2}\)。 \(2cx\) 抵消：\(y^{2}=(c^{2}/a^{2})x^{2}-x^{2}+a^{2}-c^{2}=(c^{2}/a^{2}-1)x^{2}-(c^{2}-a^{2})=(b^{2}/a^{2})x^{2}-b^{2}\)（\(b^{2}=c^{2}-a^{2}\)）。 即 \((b^{2}/a^{2})x^{2}-y^{2}=b^{2} \rightarrow  x^{2}/a^{2}-y^{2}/b^{2}=1\)。 解释：\(2.6.2\) 例 \(3\) 中所得轨迹方程正是 \(x^{2}/a^{2}-y^{2}/b^{2}=1\)（\(b^{2}=c^{2}-a^{2}\)），本题表明该双曲线可由\("\)到定点与到定直线距离之比为常数 \(c/a"\)直接刻画。 （本块自撰亲算。）}
\fi
\end{ansblock}

\begin{ansblock}[2章-拓-课时14-28]
% ans:2章-拓-课时14-28
\ansitem{28}{\(C\)}
\ifansbookdetail
\ansnote{详解}{设 \(PQ\)：\(x=my+b\)，联立 \(x^{2}-y^{2}/2=1\)：\((m^{2}-1/2)y^{2}+2bmy+b^{2}-1=0\)。 设 \(P(x_{1},y_{1})\)，\(Q(x_{2},y_{2})\)：\(y_{1}+y_{2}=-2bm/(m^{2}-1/2)\)，\(y_{1}y_{2}=(b^{2}-1)/(m^{2}-1/2)\)。 \(AP\perp AQ\)（\(A(-1,0)\)）：\((x_{1}+1)(x_{2}+1)+y_{1}y_{2}=0\)，以 \(x_{i}=my_{i}+b\) 代入： \((1+m^{2})y_{1}y_{2}+(b+1)m(y_{1}+y_{2})+(b+1)^{2}=0\)。 通分（乘 \(m^{2}-1/2\)）： \(2(b^{2}-1)(1+m^{2})-4bm^{2}(b+1)+(2m^{2}-1)(b+1)^{2}=0\)。 按 \(m\) 整理——\(m^{2}\) 系数：\(2(b^{2}-1)-4b(b+1)+2(b+1)^{2}=2b^{2}-2-4b^{2}-4b+2b^{2}+4b+2=0\)（恒为\(0\)）； 常数项：\(2(b^{2}-1)-(b+1)^{2}=b^{2}-2b-3=0 \rightarrow  b=3\) 或 \(b=-1\)。 \(b=-1\) 时 \(PQ\) 过 \(A(-1,0)\)（\(P\)、\(Q\) 之一与 \(A\) 重合或直线退化），舍去；\(b=3\)，\(PQ\) 恒过 \((3,0)\)。故选\(C\)。 验算（特例）：\(PQ\perp x\) 轴即 \(m=0\)、\(x=3\)：与 \(x^{2}-y^{2}/2=1\) 交于 \((3,\pm 4)\)；\(AP=(4,4)\)、\(AQ=(4,-4)\)，\(AP\cdot AQ=16-16=0\) 。}
\fi
\end{ansblock}

\begin{ansblock}[2章-拓-课时14-29]
% ans:2章-拓-课时14-29
\ansitem{29}{\((1) x^{2}-y^{2}/3=1\)；\((2)\) 证明见解析，\(M(-1,0)\)。}
\ifansbookdetail
\ansnote{详解}{\((1)\) 取渐近线 \(y=(b/a)x\)，与 \(x=a^{2}/c\) 交于 \(P(a^{2}/c, ab/c)\)。 \(|PF|^{2}=(c-a^{2}/c)^{2}+(ab/c)^{2}\)，其中 \(c-a^{2}/c=(c^{2}-a^{2})/c=b^{2}/c\)，故 \(|PF|^{2}=b^{4}/c^{2}+a^{2}b^{2}/c^{2}=b^{2}(a^{2}+b^{2})/c^{2}=b^{2}c^{2}/c^{2}=b^{2}\)。 由 \(|PF|=\sqrt{3}\) 得 \(b^{2}=3\)。 \(e^{2}=c^{2}/a^{2}=4=(a^{2}+b^{2})/a^{2} \rightarrow  a^{2}=1\)。 \(C\)：\(x^{2}-y^{2}/3=1\)。 \((2) F(2,0)\)。设 \(Q(x_{0},y_{0})\)（\(x_{0}\geqslant 1\)），\(3x_{0}^{2}-y_{0}^{2}=3\)。 ① \(x_{0}=2\) 时 \(y_{0}=\pm 3\)：\(\angle QFM=90^{\circ}\) 需 \(\angle QMF=45^{\circ}\)，\(|MF|=|QF|=|y_{0}|=3 \rightarrow  M(-1,0)\)，符合。 ② \(x_{0}\neq 2\) 时设 \(M(t,0)\)（\(t<0\)）。有向正切：\(tan\angle QFM=-y_{0}/(x_{0}-2)\)，\(tan\angle QMF=y_{0}/(x_{0}-t)\)。 \(\angle QFM=2\angle QMF \Rightarrow  -y_{0}/(x_{0}-2)=2[y_{0}/(x_{0}-t)]/(1-[y_{0}/(x_{0}-t)]^{2})\)（倍角公式）。 \(y_{0}\neq 0\) 时约去 \(y_{0}\)、交叉相乘：\(-[(x_{0}-t)^{2}-y_{0}^{2}]=2(x_{0}-2)(x_{0}-t)\)， 展开整理：\(3x_{0}^{2}-y_{0}^{2}-(4+4t)x_{0}+4t+t^{2}=0\)。 以 \(3x_{0}^{2}-y_{0}^{2}=3\) 代入：\(-(4+4t)x_{0}+(t^{2}+4t+3)=0\)，对任意 \(x_{0}\geqslant 1\)（\(x_{0}\neq 2\)）恒成立， 故 \(-(4+4t)=0\) 且 \(t^{2}+4t+3=0 \rightarrow  t=-1\)（两式公共解）。 \(y_{0}=0\)（\(Q\) 为右顶点）时直接验证 \(M(-1,0)\) 亦满足 \(\angle QFM=2\angle QMF\)。 综上，存在定点 \(M(-1,0)\)。 验算：\(t^{2}+4t+3=(t+1)(t+3)\) 与 \(-(4+4t)\) 公共根 \(t=-1\) ；取 \(Q(2,3)\)：\(\angle QFM=90^{\circ}\)、\(\angle QMF=45^{\circ}\)，\(2\times 45^{\circ}=90^{\circ}\) 。}
\fi
\end{ansblock}

\begin{ansblock}[2章-拓-课时14-30]
% ans:2章-拓-课时14-30
\ansitem{30}{\((1) x^{2}-y^{2}/3=1\)（\(x>0\)）；\((2) (-\infty ,-\sqrt{3})\cup (\sqrt{3},+\infty )\)；\((3)\) 存在，\(M(-1,0)\)。}
\ifansbookdetail
\ansnote{详解}{\((1) |PF_{1}|-|PF_{2}|=2<|F_{1}F_{2}|=4\)，轨迹为以 \(F_{1}\)、\(F_{2}\) 为焦点、实轴长 \(2\) 的双曲线右支：\(c=2\)，\(a=1\)，\(b^{2}=3\)，\(\Gamma \)：\(x^{2}-y^{2}/3=1\)（\(x>0\)）。 \((2) l\)：\(y=k(x-2)\)，代入 \(\Gamma \)：\((3-k^{2})x^{2}+4k^{2}x-4k^{2}-3=0\)。 \(A\)、\(B\) 为两交点（右支上 \(x>0\)）：需 \(3-k^{2}\neq 0\)、\(\Delta >0\)、\(x_{1}+x_{2}=-4k^{2}/(3-k^{2})>0\)、\(x_{1}x_{2}=(-4k^{2}-3)/(3-k^{2})>0\)。 由 \(x_{1}+x_{2}>0\) 得 \(k^{2}>3\)；此时 \(x_{1}x_{2}>0\)、\(\Delta =16k^{4}-4(3-k^{2})(-4k^{2}-3)>0\) 亦成立。 故 \(k\in (-\infty ,-\sqrt{3})\cup (\sqrt{3},+\infty )\)。 \((3)\) 存在性：设 \(M(m,0)\)。\(l\) 与 \(\Gamma \) 联立同上，\((3-k^{2})x^{2}+4k^{2}x-4k^{2}-3=0\)（\(k^{2}>3\)）， \(x_{1}+x_{2}=-4k^{2}/(3-k^{2})\)，\(x_{1}x_{2}=(-4k^{2}-3)/(3-k^{2})\)。 \(MA\perp MB \Leftrightarrow  (x_{1}-m)(x_{2}-m)+y_{1}y_{2}=0\)，\(y_{i}=k(x_{i}-2)\)： \((1+k^{2})x_{1}x_{2}-(m+2k^{2})(x_{1}+x_{2})+m^{2}+4k^{2}=0\)。 乘 \((3-k^{2})\) 展开： \((1+k^{2})(-4k^{2}-3)+4k^{2}(2k^{2}+m)+(m^{2}+4k^{2})(3-k^{2})=0 \rightarrow  k^{2}(-m^{2}+4m+5)+3m^{2}-3=0\)，对一切 \(k^{2}>3\) 恒成立。 故 \(-m^{2}+4m+5=0\) 且 \(3m^{2}-3=0 \rightarrow  m=-1\)（\(m=1\) 不满足第一式）。 斜率不存在时 \(l\)：\(x=2\)，\(A(2,3)\)、\(B(2,-3)\)：\(MA\cdot MB=(3,3)\cdot (3,-3)=0\) 。 综上，存在 \(M(-1,0)\)。 验算：\(k^{2}(-m^{2}+4m+5)+3m^{2}-3\) 在 \(m=-1\) 时两项恒为\(0\) ；\((3)\) 的韦达代入式与 \((2)\) 一致 。}
\fi
\end{ansblock}

\begin{ansblock}[2章-拓-课时14-32-1]
% ans:2章-拓-课时14-32-1
\ansitem{32-1}{\(D\)}
\ifansbookdetail
\ansnote{详解}{\(P(x_{0},y_{0})\)、\(Q(-x_{0},-y_{0})\)，\(\triangle APQ\) 三顶点坐标和为 \((a,0)\)，其重心为 \(G(a/3,0)\)。 重心在每条中线上，故 \(G\) 在中线 \(QE\) 上；\(QE\) 与 \(x\) 轴的交点唯一，故 \(M(1,0)\) 即重心：\(a/3=1 \rightarrow  a=3\)。 \(b^{2}=1\)，\(c=\sqrt{9+1}=\sqrt{10}\)，\(e=c/a=\sqrt{10}/3\)，故选\(D\)。 验算：\(P\)、\(Q\)、\(A\) 的纵坐标和 \(y_{0}-y_{0}+0=0\)，重心在 \(x\) 轴 ；中线 \(QE\) 过 \((a/3,0)\) 由重心坐标公式 。}
\fi
\end{ansblock}

\begin{ansblock}[2章-拓-课时14-32-2]
% ans:2章-拓-课时14-32-2
\ansitem{32-2}{\((0, 3\sqrt{3})\)}
\ifansbookdetail
\ansnote{详解}{\(e=2 \rightarrow  c^{2}=4a^{2} \rightarrow  b^{2}=c^{2}-a^{2}=3a^{2}\)，\(b/a=\sqrt{3}\)。 \(k_{1}k_{2}=[y_{0}/(x_{0}+a)]\cdot [y_{0}/(x_{0}-a)]=y_{0}^{2}/(x_{0}^{2}-a^{2})=(b^{2}/a^{2})(x_{0}^{2}-a^{2})/(x_{0}^{2}-a^{2})=b^{2}/a^{2}=3\)。 \(P\) 在第一象限右支：\(k_{3}=k\_PO=y_{0}/x_{0}\in (0, b/a)=(0,\sqrt{3})\)（\(P\) 趋近顶点时 \(k_{3}\rightarrow 0\)，沿渐近线方向 \(k_{3}\rightarrow \sqrt{3}\)）。 \(m=3k_{3}\in (0,3\sqrt{3})\)。 （本块自撰亲算：坐标点差法同 \(14-12\)，\(k_{1}k_{2}=b^{2}/a^{2}\) 。）}
\fi
\end{ansblock}

\begin{ansblock}[2章-拓-课时14-33-1]
% ans:2章-拓-课时14-33-1
\ansitem{33-1}{\([\sqrt{3}, 2\sqrt{3})\)}
\ifansbookdetail
\ansnote{详解}{机器人到 \(P\) 的距离比到 \(Q\) 大 \(2\)：距离差 \(2<|PQ|=4\)，且距 \(P\) 更远，轨迹为双曲线右支：\(a=1\)，\(c=2\)，\(b^{2}=3\)，即 \(x^{2}-y^{2}/3=1\)（\(x\geqslant 1\)）。 直线 \(y-3=k(x-\sqrt{3})\)。联立得 \((3-k^{2})x^{2}+(2\sqrt{3}k^{2}-6k)x+6\sqrt{3}k-3k^{2}-12=0\)。 直线与右支无公共点：\(k=\sqrt{3}\) 时直线恰为渐近线 \(y=\sqrt{3}x\)，符合（接触不到）；\(k=-\sqrt{3}\) 不合。 \(3-k^{2}\neq 0\) 时由 \(\Delta <0\) 解得 \(\sqrt{3}<k<2\sqrt{3}\)。 综上 \(k\in [\sqrt{3},2\sqrt{3})\)。 验算：\(k=\sqrt{3}\) 时直线 \(y-3=\sqrt{3}(x-\sqrt{3})\) 即 \(y=\sqrt{3}x\)，恰为渐近线，永不相触 ；\(k=2\sqrt{3}\) 时 \(\Delta =0\)，直线与右支相切（可触），故右端开 。}
\fi
\end{ansblock}

\begin{ansblock}[2章-拓-课时14-33-2]
% ans:2章-拓-课时14-33-2
\ansitem{33-2}{\(A\)}
\ifansbookdetail
\ansnote{详解}{\(e=2 \rightarrow  b^{2}/a^{2}=3\)，直线 \(MN\)：\(y=\sqrt{3}(x+a)\)。设 \(P(t,\sqrt{3}t+\sqrt{3}a)\)，\(t\in [-a,0]\)。 \(PF_{1}\cdot PF_{2}=(-c-t,-\sqrt{3}t-\sqrt{3}a)\cdot (c-t,-\sqrt{3}t-\sqrt{3}a)=t^{2}-c^{2}+3(t+a)^{2}=4t^{2}+6at-a^{2}=4(t+3a/4)^{2}-13a^{2}/4\)。 \(t=-3a/4\) 取最小（\(y\_P=\sqrt{3}a/4\)），\(t=0\) 取最大（\(y\_P=\sqrt{3}a\)）。 面积比\(=\)纵坐标绝对值比：\(S_{1}/S_{2}=\sqrt{3}a/(\sqrt{3}a/4)=4\)，故选\(A\)。 验算：\(4t^{2}+6at-a^{2}\) 的导数 \(8t+6a=0 \rightarrow  t=-3a/4\)（顶点值 \(-13a^{2}/4\)）；端点 \(t=-a\)：\(-3a^{2}=-12a^{2}/4\)、\(t=0\)：\(-a^{2}=-4a^{2}/4\)，均大于 \(-13a^{2}/4\)，故最小在顶点、最大在 \(t=0\) 。}
\fi
\end{ansblock}

\begin{ansblock}[2章-拓-课时14-33-3]
% ans:2章-拓-课时14-33-3
\ansitem{33-3}{\(2\)}
\ifansbookdetail
\ansnote{详解}{\(D(a,b)\)，\(E(a,-b)\)，\(|DE|=2b\)，\(O\) 到直线 \(x=a\) 距离为 \(a\)：\(S=(1/2)\cdot a\cdot 2b=ab\)。 焦距 \(4 \rightarrow  c=2\)，\(a^{2}+b^{2}=4\geqslant 2ab \rightarrow  ab\leqslant 2\)（\(a=b=\sqrt{2}\) 取等）。 \(\triangle ODE\) 面积最大值 \(2\)。}
\fi
\end{ansblock}

\begin{ansblock}[2章-拓-课时14-33-4]
% ans:2章-拓-课时14-33-4
\ansitem{33-4}{\((1) e=2\)；\((2) 3/2\)。}
\ifansbookdetail
\ansnote{详解}{\((1)\) 右支上 \(|PF_{1}|min=c+a\)（左顶点处），\(|PF_{2}|min=c-a\)（右顶点处）。 \(m_{1}m_{2}=c^{2}-a^{2}=b^{2}=3a^{2} \rightarrow  c^{2}=4a^{2} \rightarrow  e=2\)。 \((2) a=2\)，\(c=4\)，\(b^{2}=12\)，双曲线 \(x^{2}/4-y^{2}/12=1\)，\(F_{1}(-4,0)\)。设 \(AB\)：\(y=k(x+4)\)（\(k\neq \pm \sqrt{3}\)）。 联立：\((3-k^{2})x^{2}-8k^{2}x-16k^{2}-12=0\)，\(x_{1}+x_{2}=8k^{2}/(3-k^{2})\)，\(x_{1}x_{2}=(-16k^{2}-12)/(3-k^{2})\)。 弦长 \(|AB|=\sqrt{1+k^{2}}\cdot \sqrt{(x_{1}+x_{2})^{2}-4x_{1}x_{2}}=12(1+k^{2})/|3-k^{2}|\)。 中点 \((4k^{2}/(3-k^{2}), 12k/(3-k^{2}))\)（\(y_{1}+y_{2}=k(x_{1}+x_{2})+8k=24k/(3-k^{2})\)）。 垂直平分线：\(y-12k/(3-k^{2})=-(1/k)(x-4k^{2}/(3-k^{2}))\)，令 \(x=0\) 得 \(y=16k/(3-k^{2})\)。 \(|OD|=16|k|/|3-k^{2}|\)。 \(|AB|/|OD|=[12(1+k^{2})/|3-k^{2}|]\cdot [|3-k^{2}|/(16|k|)]=(3/4)(|k|+1/|k|)\geqslant (3/2)\)（\(|k|=1\) 取等）。 最小值 \(3/2\)。 验算：\(|AB|\) 的核：\(\Delta \) 下根差\(^{2}=(8k^{2}/(3-k^{2}))^{2}+4(16k^{2}+12)/(3-k^{2})=[64k^{4}+4(16k^{2}+12)(3-k^{2})]/(3-k^{2})^{2}=[64k^{4}+192k^{2}+144-64k^{4}-48k^{2}]/(3-k^{2})^{2}=144(k^{2}+1)/(3-k^{2})^{2} \rightarrow  |AB|=\sqrt{1+k^{2}}\cdot 12\sqrt{k^{2}+1}/|3-k^{2}|\) 。}
\fi
\end{ansblock}

\begin{ansblock}[2章-拓-课时14-33-5]
% ans:2章-拓-课时14-33-5
\ansitem{33-5}{\((1) x^{2}/2-y^{2}=1\)；\((2)\) 过定点 \((0,1)\)。}
\ifansbookdetail
\ansnote{详解}{\((1) e^{2}=3/2=c^{2}/a^{2}=(a^{2}+b^{2})/a^{2} \rightarrow  b^{2}/a^{2}=1/2 \rightarrow  a^{2}=2b^{2}\)，即双曲线方程为 \(x^{2}/(2b^{2})-y^{2}/b^{2}=1\)。 代入 \(P\)：\(3/(2b^{2})-(1/2)/b^{2}=1 \rightarrow  (3-1)/(2b^{2})=1 \rightarrow  b^{2}=1\)，\(a^{2}=2\)。\(C\)：\(x^{2}/2-y^{2}=1\)。 \((2) AB\) 斜率不存在时：\(A(x_{0},y_{0})\)、\(B(x_{0},-y_{0})\)，\(k_{1}+k_{2}=(y_{0}-1)/(x_{0}-2)+(-y_{0}-1)/(x_{0}-2)=-2/(x_{0}-2)=1 \rightarrow  x_{0}=0\)，不在曲线上，舍；故斜率存在。 设 \(AB\)：\(y=kx+t\)，代入：\((2k^{2}-1)x^{2}+4ktx+2t^{2}+2=0\)（\(2k^{2}-1\neq 0\)）。 \(x_{1}+x_{2}=-4kt/(2k^{2}-1)\)，\(x_{1}x_{2}=(2t^{2}+2)/(2k^{2}-1)\)。 \(k_{1}+k_{2}=1\)：\((kx_{1}+t-1)/(x_{1}-2)+(kx_{2}+t-1)/(x_{2}-2)=1 \rightarrow  (2k-1)x_{1}x_{2}+(t-2k+1)(x_{1}+x_{2})-4t=0\)。 代入：\((2k-1)(2t^{2}+2)-4kt(t-2k+1)-4t(2k^{2}-1)=0\)（通分后乘 \(2k^{2}-1\)）， 整理 \(t^{2}+(2k-2)t-1+2k=0\)，即 \((t-1)(t+2k-1)=0\)。 \(t=1\)：\(y=kx+1\) 过 \((0,1)\)；\(t=1-2k\)：\(y=k(x-2)+1\) 过 \(Q(2,1)\)，不合（\(A\)、\(B\) 异于 \(Q\) 且 \(Q\) 为弦端情形）。 故直线 \(AB\) 过定点 \((0,1)\)。 验算：\(t^{2}+(2k-2)t+2k-1=(t-1)(t+2k-1)\) 展开核对 。}
\fi
\end{ansblock}

\begin{ansblock}[2章-拓-课时14-33-6]
% ans:2章-拓-课时14-33-6
\ansitem{33-6}{\((1) x^{2}/4-y^{2}=1\)；\((2)\) 存在，\(Q(23/8,0)\)，\(QM\cdot QN=273/64\)。}
\ifansbookdetail
\ansnote{详解}{\((1) b/a=1/2\) 且 \(16/a^{2}-3/b^{2}=1\)：\(b^{2}=a^{2}/4 \rightarrow  16/a^{2}-12/a^{2}=1 \rightarrow  a^{2}=4\)，\(b^{2}=1\)。\(C\)：\(x^{2}/4-y^{2}=1\)。 \((2)\) 设 \(l\)：\(x=my+1\)，联立：\((m^{2}-4)y^{2}+2my-3=0\)（\(m^{2}\neq 4\)，\(\Delta >0 \rightarrow  m^{2}>3\)）。 \(y_{1}+y_{2}=-2m/(m^{2}-4)\)，\(y_{1}y_{2}=-3/(m^{2}-4)\)； \(x_{1}+x_{2}=m(y_{1}+y_{2})+2=-8/(m^{2}-4)\)，\(x_{1}x_{2}=(my_{1}+1)(my_{2}+1)=-(4m^{2}+4)/(m^{2}-4)\)。 设 \(Q(t,0)\)：\(QM\cdot QN=(x_{1}-t)(x_{2}-t)+y_{1}y_{2}=x_{1}x_{2}-t(x_{1}+x_{2})+t^{2}+y_{1}y_{2} \allowbreak =-4+t^{2}+(8t-23)/(m^{2}-4)\)。 与 \(m\) 无关 \(\Leftrightarrow  8t-23=0 \rightarrow  t=23/8\)，常数\(=-4+t^{2}=-4+529/64=273/64\)。 验算：\(x_{1}x_{2}=(-4m^{2}-4)/(m^{2}-4)=-4-20/(m^{2}-4)\)；三项合并分子 \(-20-3+8t=8t-23 \rightarrow  QM\cdot QN=-4+t^{2}+(8t-23)/(m^{2}-4)\) 。}
\fi
\end{ansblock}

\begin{ansblock}[2章-拓-课时14-33-7]
% ans:2章-拓-课时14-33-7
\ansitem{33-7}{\((1) 15/4\)；\((2)\) 双曲线中 \(k\_AB\cdot k\_OM=b^{2}/a^{2}\)（定值），证明见解析。}
\ifansbookdetail
\ansnote{详解}{\((1)\) 联立得 \(16x^{2}+12x-9=0\)，\(x_{1}+x_{2}=-3/4\)，\(x_{1}x_{2}=-9/16\)。 \(|AB|=\sqrt{1+2^{2}}\cdot \sqrt{(-3/4)^{2}-4\times (-9/16)}=\sqrt{5}\cdot \sqrt{9/16+9/4}=\sqrt{5}\cdot \sqrt{45/16}=\sqrt{5}\cdot (3\sqrt{5})/4=15/4\)。 \((2)\) 性质：双曲线 \(x^{2}/a^{2}-y^{2}/b^{2}=1\) 上任意两点 \(A\)、\(B\)（斜率存在时），\(M\) 为 \(AB\) 中点，则 \(k\_AB\cdot k\_OM=b^{2}/a^{2}\)。 证明（点差法）：\(A(x_{3},y_{3})\)、\(B(x_{4},y_{4})\)、\(M(x_{0},y_{0})\)。两式相减： \((x_{3}+x_{4})(x_{3}-x_{4})/a^{2}=(y_{3}+y_{4})(y_{3}-y_{4})/b^{2} \rightarrow  2x_{0}(x_{3}-x_{4})/a^{2}=2y_{0}(y_{3}-y_{4})/b^{2}\)。 \(k\_AB\cdot k\_OM=[(y_{3}-y_{4})/(x_{3}-x_{4})]\cdot (y_{0}/x_{0})=b^{2}/a^{2}\)。定值证毕。 验算：本件 \(k\_AB\cdot k\_OM=b^{2}/a^{2}\) 与 \(14-12\) 的 \(k_{1}k_{2}=b^{2}/a^{2}\) 相容（中心对称两点为其特例）。}
\fi
\end{ansblock}

\begin{ansblock}[2章-拓-课时14-33-8]
% ans:2章-拓-课时14-33-8
\ansitem{33-8}{\((1) x^{2}-y^{2}/8=1\)；\((2)\) 定值 \(2\sqrt{2}\)。}
\ifansbookdetail
\ansnote{详解}{\((1) c=3\)，\(b/a=2\sqrt{2}\)，\(a^{2}+b^{2}=9 \rightarrow  a^{2}(1+8)=9 \rightarrow  a^{2}=1\)，\(b^{2}=8\)。\(C\)：\(x^{2}-y^{2}/8=1\)。 \((2) l\)：\(y=kx+m\)（\(k\neq \pm 2\sqrt{2}\)，\(m\neq 0\)，斜率存在且不为\(0\)）。\(D(-m/k,0)\)，\(|OD|=|m/k|\)。 相切：联立 \((8-k^{2})x^{2}-2kmx-m^{2}-8=0\)，\(\Delta =0 \rightarrow  4k^{2}m^{2}+4(8-k^{2})(m^{2}+8)=0 \rightarrow  m^{2}=k^{2}-8\)（\(k^{2}>8\)）。 \(M\) 与 \(y=2\sqrt{2}x\) 交：\((m/(2\sqrt{2}-k), 2\sqrt{2}m/(2\sqrt{2}-k))\)；\(N\) 与 \(y=-2\sqrt{2}x\) 交：\((-m/(2\sqrt{2}+k), 2\sqrt{2}m/(2\sqrt{2}+k))\)。 \(S\triangle MON=S\triangle MOD+S\triangle NOD=(1/2)|OD|\cdot |y\_M-y\_N|=(1/2)|m|\cdot |x\_M-x\_N| \allowbreak =(1/2)|m|\cdot |4\sqrt{2}m/(8-k^{2})|=2\sqrt{2}m^{2}/(k^{2}-8)=2\sqrt{2}\)（定值）。 验算：\(m^{2}=k^{2}-8>0\) 与 \(k^{2}>8\) 一致 ；\(S\) 只依赖 \(m^{2}/k^{2}-8\) 化简后为常数 。}
\fi
\end{ansblock}

\begin{ansblock}[2章-拓-课时14-33-9]
% ans:2章-拓-课时14-33-9
\ansitem{33-9}{\((1) y=\pm \sqrt{3}x\)；\((2)\) 定值 \(k_{1}/k_{2}=-1/3\)。}
\ifansbookdetail
\ansnote{详解}{\((1) a=1\)，\(c=2\)，\(b^{2}=3\)：\(\Gamma \)：\(x^{2}-y^{2}/3=1\)，渐近线 \(y=\pm \sqrt{3}x\)。 \((2) l\) 斜率不存在：\(P(2,3)\)、\(Q(2,-3)\)，\(k_{1}=3/3=1\)、\(k_{2}=-3/1=-3\)，比值 \(-1/3\)。 \(l\)：\(y=k(x-2)\)（\(k\neq 0\)）：\((k^{2}-3)x^{2}-4k^{2}x+4k^{2}+3=0\)，\(x_{1}+x_{2}=4k^{2}/(k^{2}-3)\)，\(x_{1}x_{2}=(4k^{2}+3)/(k^{2}-3)\)。 \(3k_{1}+k_{2}=3y_{1}/(x_{1}+1)+y_{2}/(x_{2}-1)=k[4x_{1}x_{2}-5(x_{1}+x_{2})+4]/[(x_{1}+1)(x_{2}-1)]\)。 \(4x_{1}x_{2}-5(x_{1}+x_{2})+4=[4(4k^{2}+3)-20k^{2}+4(k^{2}-3)]/(k^{2}-3)=0\)。 故 \(3k_{1}+k_{2}=0\)，\(k_{1}/k_{2}=-1/3\)（定值）。综上得证。 验算：分子 \(16k^{2}+12-20k^{2}+4k^{2}-12=0\) 。}
\fi
\end{ansblock}

\begin{ansblock}[2章-拓-课时14-33-10]
% ans:2章-拓-课时14-33-10
\ansitem{33-10}{\((1) 2x^{2}-y^{2}/2=1\)；\((2)\) 定直线 \(12x-y-2=0\)。}
\ifansbookdetail
\ansnote{详解}{\((1) 2b=2\sqrt{2} \rightarrow  b^{2}=2\)；\(e^{2}=5=1+b^{2}/a^{2} \rightarrow  a^{2}=1/2\)。\(C\)：\(2x^{2}-y^{2}/2=1\)。 \((2)\) 设 \(M(x,y)\)、\(A(x_{1},y_{1})\)、\(B(x_{2},y_{2})\)（\(x_{1}<x_{2}<3\)）。\(|AM|/|MB|=|AP|/|PB|\) 结合共线定比： \((3-x_{1})/(x_{2}-3)=-(x-x_{1})/(x_{2}-x)\)，即 \([6-(x_{1}+x_{2})]x=3(x_{1}+x_{2})-2x_{1}x_{2}\)。① 设 \(l\)：\(y-1=k(x-3)\)，代入 \(C\)：\((4-k^{2})x^{2}+(6k^{2}-2k)x-9k^{2}+6k-3=0\)， \(x_{1}+x_{2}=(2k-6k^{2})/(4-k^{2})\)，\(x_{1}x_{2}=(-9k^{2}+6k-3)/(4-k^{2})\)。代入①整理： \(12x-3=k(x-3)\)，与 ②（\(y-1=k(x-3)\)）联立消 \(k\)：\(12x-3=y-1\)，即 \(12x-y-2=0\)。 点 \(M\) 恒落在定直线 \(12x-y-2=0\) 上。 验算：① 的展开：由定比 \(x=(3(x_{1}+x_{2})-2x_{1}x_{2})/(6-(x_{1}+x_{2}))\)（分比不变）；韦达代入化简 \((6-(x_{1}+x_{2}))\cdot 12x=36(x_{1}+x_{2})-24x_{1}x_{2}\) 与 \(12x-3=k(x-3)\) 等价核：源详解整理结果 \(12x-3=k(x-3)\)，代回 \(y-1=k(x-3)\) 消 \(k\) 得 \(y=12x-2\) 。}
\fi
\end{ansblock}

\begin{ansblock}[2章-拓-课时14-33-11]
% ans:2章-拓-课时14-33-11
\ansitem{33-11}{\((1) 2\sqrt{5}x-y-4\sqrt{5}=0\) 或 \(2\sqrt{5}x+y-4\sqrt{5}=0\)；\((2)\) 定直线 \(x=1/2\)。}
\ifansbookdetail
\ansnote{详解}{设 \(l\)：\(x=my+2\)，\(C(x_{1},y_{1})\)、\(D(x_{2},y_{2})\)，联立：\((4m^{2}-1)y^{2}+16my+12=0\)， \(y_{1}+y_{2}=-16m/(4m^{2}-1)\)，\(y_{1}y_{2}=12/(4m^{2}-1)\)。 \((1) CN=3ND\)（\(N\) 为 \(F_{2}\)，\(C\)、\(D\) 同向分段）\(\rightarrow  y_{1}=-3y_{2}\)。 \(y_{2}=8m/(4m^{2}-1)\)，代入 \(y_{1}y_{2}=-3y_{2}^{2}=12/(4m^{2}-1)\)： \(-3\cdot 64m^{2}/(4m^{2}-1)^{2}=12/(4m^{2}-1) \rightarrow  -16m^{2}=4m^{2}-1 \rightarrow  m^{2}=1/20\)，\(m=\pm \sqrt{5}/10\)。 直线 \(l\)：\(x=\pm (\sqrt{5}/10)y+2\)，即 \(2\sqrt{5}x-y-4\sqrt{5}=0\) 或 \(2\sqrt{5}x+y-4\sqrt{5}=0\)。 \((2) AC\)：\(y=[y_{1}/(x_{1}+1)](x+1)\)；\(BD\)：\(y=[y_{2}/(x_{2}-1)](x-1)\)。以 \(x_{1}=my_{1}+2\)、\(x_{2}=my_{2}+2\) 代入，得 \(x_{1}+1=my_{1}+3\)、\(x_{2}-1=my_{2}+1\)。 联立：\((x+1)/(x-1)=(my_{1}y_{2}+3y_{2})/(my_{1}y_{2}+y_{1})\)。 一般情形（不求 \(y_{1}\)、\(y_{2}\) 具体值）：由韦达 \(y_{1}y_{2}=12/(4m^{2}-1)\)、\(y_{1}+y_{2}=-16m/(4m^{2}-1)\)，得 \(y_{1}y_{2}=-(3/4m)(y_{1}+y_{2})\)。 于是 \((x+1)/(x-1)=[-(3/4)(y_{1}+y_{2})+3y_{2}]/[-(3/4)(y_{1}+y_{2})+y_{1}]=(-3y_{1}+9y_{2})/(y_{1}-3y_{2})=-3\)。 解得 \(x=1/2\)。点 \(P\) 恒在定直线 \(x=1/2\) 上。 验算：\((-3y_{1}+9y_{2})/(y_{1}-3y_{2})=-3(y_{1}-3y_{2})/(y_{1}-3y_{2})=-3\) ；\((1)\) 中 \(-16m^{2}=4m^{2}-1 \rightarrow  20m^{2}=1\) 。}
\fi
\end{ansblock}

\begin{ansblock}[2章-拓-课时14-33-12]
% ans:2章-拓-课时14-33-12
\ansitem{33-12}{\(\sqrt{5}\)}
\ifansbookdetail
\ansnote{详解}{以最细处直径 \(A_{1}A_{2}\) 为 \(x\) 轴、中垂线为 \(y\) 轴建系。\(2a=16\)，\(a=8\)。 上瓶口点 \(B(10,12)\) 在双曲线 \(x^{2}/64-y^{2}/b^{2}=1\) 上：\(100/64-144/b^{2}=1 \rightarrow  144/b^{2}=100/64-1=36/64 \rightarrow  b^{2}=256\)。 \(c^{2}=a^{2}+b^{2}=64+256=320\)，\(e^{2}=320/64=5\)，\(e=\sqrt{5}\)。}
\fi
\end{ansblock}

\begin{ansblock}[2章-拓-课时14-33-13]
% ans:2章-拓-课时14-33-13
\ansitem{33-13}{在 \(A\) 的北偏东 \(30^{\circ}\) 方向，相距 \(10\) 海里处。}
\ifansbookdetail
\ansnote{详解}{以 \(AB\) 为 \(x\) 轴、\(AB\) 中点为原点建系：\(A(3,0)\)，\(B(-3,0)\)，\(C(-5,2\sqrt{3})\)。 \(|PB|-|PA|=4<6=|AB|\)，\(P\) 在以 \(A\)、\(B\) 为焦点、实轴长 \(4\) 的双曲线右支上：\(x^{2}/4-y^{2}/5=1\)（\(x\geqslant 2\)）。 \(|PB|=|PC|\)（\(B\)、\(C\) 同时收到），\(P\) 在 \(BC\) 中垂线上：\(BC\) 中点 \((-4,\sqrt{3})\)，\(k\_BC=2\sqrt{3}/(-5+3)=-\sqrt{3}\)， 中垂线：\(y-\sqrt{3}=(\sqrt{3}/3)(x+4)\)，即 \(x-\sqrt{3}y+7=0\)。 与 \(5x^{2}-4y^{2}=20\) 联立（\(y=(x+7)/\sqrt{3}\) 代入）：\(11x^{2}-56x-256=0\)，解得 \(x=8\) 或 \(x=-32/11\)（左支交点，不合右支 \(x\geqslant 2\)，舍），\(y=5\sqrt{3}\)。 \(P(8,5\sqrt{3})\)：\(k\_AP=5\sqrt{3}/(8-3)=\sqrt{3}\)，倾斜角 \(60^{\circ}\)，\(|PA|=\sqrt{25+75}=10\)。 故 \(P\) 在 \(A\) 的北偏东 \(30^{\circ}\) 方向，相距 \(10\) 海里处。 验算：\(|PB|=\sqrt{(8+3)^{2}+(5\sqrt\{3\})^{2}}=\sqrt{196}=14\)，\(|PB|-|PA|=14-10=4\) ；\(|PC|^{2}=13^{2}+(3\sqrt{3})^{2}=196=|PB|^{2}\) 。}
\fi
\end{ansblock}

\begin{ansblock}[2章-拓-课时14-34-1]
% ans:2章-拓-课时14-34-1
\ansitem{34-1}{\(2\)}
\ifansbookdetail
\ansnote{详解}{两者焦点均在 \(x\) 轴：\(9-k^{2}=k+3\)（\(k>0\)）\(\rightarrow  k^{2}+k-6=0 \rightarrow  (k+3)(k-2)=0 \rightarrow  k=2\)（\(k=-3\) 舍）。 验算：\(k=2\)：椭圆 \(c^{2}=9-4=5\)，双曲线 \(c^{2}=2+3=5\) 。}
\fi
\end{ansblock}

\begin{ansblock}[2章-拓-课时14-34-2]
% ans:2章-拓-课时14-34-2
\ansitem{34-2}{\(C\)}
\ifansbookdetail
\ansnote{详解}{以 \(O\) 为原点、\(AD\) 所在直线为 \(x\) 轴建系。八等分结合 \(AB=1\)、\(BC=2\)、\(CD=1\)：圆半径 \(R=2\)，双曲线过点 \((\sqrt{2},\sqrt{2})\)，\(a=1\)。 \(2/1-2/b^{2}=1 \rightarrow  b^{2}=2\)，\(c^{2}=1+2=3\)，\(c=\sqrt{3}\)，焦距 \(2\sqrt{3}\)，故选\(C\)。}
\fi
\end{ansblock}

\begin{ansblock}[2章-拓-课时14-34-3]
% ans:2章-拓-课时14-34-3
\ansitem{34-3}{\(B\)}
\ifansbookdetail
\ansnote{详解}{\(P_{0}F_{1}\cdot PF_{1}=0 \rightarrow  P_{0}F_{1}\perp PF_{1}\)；\(P_{0}\) 与 \(P\) 关于原点对称 \(\rightarrow  |P_{0}F_{2}|=|PF_{1}|=a\) 且 \(P_{0}F_{1}\perp P_{0}F_{2}\)（\(P_{0}F_{2}\parallel PF_{1}\)）。 \(P_{0}\) 在右支：\(|P_{0}F_{1}|-|P_{0}F_{2}|=2a \rightarrow  |P_{0}F_{1}|=3a\)。 \(Rt\triangle P_{0}F_{1}F_{2}\)（直角在 \(P_{0}\)）：\((3a)^{2}+a^{2}=4c^{2} \rightarrow  10a^{2}=4c^{2}=4(a^{2}+b^{2}) \rightarrow  b^{2}/a^{2}=3/2\)。 渐近线 \(y=\pm (\sqrt{6}/2)x\)，故选\(B\)。}
\fi
\end{ansblock}

\begin{ansblock}[2章-拓-课时14-34-4]
% ans:2章-拓-课时14-34-4
\ansitem{34-4}{\((-\infty ,-2]\)}
\ifansbookdetail
\ansnote{详解}{设 \(P(x_{0},y_{0})\)（\(|x_{0}|\geqslant \sqrt{3}\)），\(Q(-x_{0},-y_{0})\)。 \(MP=(x_{0},y_{0}-1)\)，\(MQ=(-x_{0},-y_{0}-1)\)，\(MP\cdot MQ=-x_{0}^{2}-y_{0}^{2}+1\)。 由 \(y_{0}^{2}=x_{0}^{2}/3-1\)：\(MP\cdot MQ=-x_{0}^{2}-x_{0}^{2}/3+1+1=2-(4/3)x_{0}^{2}\leqslant 2-(4/3)\times 3=-2\)。 取值范围 \((-\infty ,-2]\)。}
\fi
\end{ansblock}

\begin{ansblock}[2章-拓-课时14-34-5]
% ans:2章-拓-课时14-34-5
\ansitem{34-5}{最大值为 \(5/3\)。}
\ifansbookdetail
\ansnote{详解}{\(|PF_{1}|-|PF_{2}|=3|PF_{2}|=2a \rightarrow  |PF_{2}|=(2/3)a\)。 右支焦半径 \(|PF_{2}|\geqslant c-a\)：\((2/3)a\geqslant c-a \rightarrow  c\leqslant (5/3)a \rightarrow  e\leqslant 5/3\)；又 \(e>1\)。 故 \(e\) 的最大值 \(5/3\)。}
\fi
\end{ansblock}

\begin{ansblock}[2章-拓-课时14-34-6]
% ans:2章-拓-课时14-34-6
\ansitem{34-6}{\(A\)}
\ifansbookdetail
\ansnote{详解}{\(e^{2}=3/2=c^{2}/a^{2}=(a^{2}+b^{2})/a^{2} \rightarrow  b^{2}=a^{2}/2\)。 过 \((1,-2)\)：\(4/a^{2}-1/b^{2}=1 \rightarrow  4/a^{2}-2/a^{2}=1 \rightarrow  a^{2}=2\)，\(b^{2}=1\)。 方程 \(y^{2}/2-x^{2}=1\)，故选\(A\)。}
\fi
\end{ansblock}

\begin{ansblock}[2章-拓-课时15-01]
% ans:2章-拓-课时15-01
\ansitem{01}{桥拱所在抛物线 \(y=-(1/80)x^{2}\)；溢流孔所在抛物线分别为 \(y=-(5/36)(x+14)^{2}\)、\(y=-(5/36)(x+7)^{2}\)、\(y=-(5/36)(x-7)^{2}\)、\(y=-(5/36)(x-14)^{2}\)；交点 \(A((140/13),-(245/169))\)、\(B(10,-5/4)\)、\(C((70/13),-(245/676))\)。}
\ifansbookdetail
\ansnote{详解}{设桥拱 \(y=ax^{2}\)。图注（转置说明）：题面与源文按图代入的是点 \((20,-5)\)（源详解 \(-5=a\times 20^{2}\)，\(dump\) 文字「过点\((-5,20)\)」系转置，标注尺寸以源图为准）： \(-5=a\times 20^{2} \rightarrow  a=-1/80\)，桥拱 \(y=-(1/80)x^{2}\)。① 设 \(A\) 所在溢流孔 \(y=a_{1}(x-14)^{2}\)（顶点 \((14,0)\)），亦过 \((20,-5)\)：\(-5=a_{1}\times 36 \rightarrow  a_{1}=-5/36\)。 \(4\) 孔轮廓相同、顶点 \((\pm 7,0)\)、\((\pm 14,0)\)：\(y=-(5/36)(x\mp 7)^{2}\)、\(y=-(5/36)(x\mp 14)^{2}\)。②③ 联立①②：\(400(x-14)^{2}=36x^{2}\)（两边乘 \(2880\)）\(\rightarrow  91x^{2}-2800x+19600=0\)，\(\Delta =705600\)，\(\sqrt{}\Delta =840 \rightarrow  x=20\)（\(y=-5\)）或 \(x=140/13\)（\(y=-245/169\)），取交点 \(A(140/13,-245/169)\)； 联立①③：\(400(x-7)^{2}=36x^{2} \rightarrow  91x^{2}-1400x+4900=0\)，\(\Delta =176400\)，\(\sqrt{}\Delta =420 \rightarrow  x=10\)（\(y=-5/4\)）或 \(x=70/13\)（\(y=-245/676\)），得 \(C(70/13,-245/676)\)、\(B(10,-5/4)\)。 验算：\(x=140/13\)：\(y=-(1/80)(140/13)^{2}=-19600/(80\times 169)=-245/169\) ；\(x=10\)：\(-100/80=-5/4\) ；\(x=70/13\)：\(-4900/(80\times 169)=-245/676\) 。}
\fi
\end{ansblock}

\begin{ansblock}[2章-拓-课时15-02]
% ans:2章-拓-课时15-02
\ansitem{02}{\(A\)}
\ifansbookdetail
\ansnote{详解}{\(y^{2}=8x\)：\(p=4\)，准线 \(x=-2\)，焦点 \(F(2,0)\)。 \(P\) 在准线 \(x=-2\) 上，又在 \(x-y+6=0\) 上：\(P(-2,4)\)。 \(k\_PF=(0-4)/(2+2)=-1\)；由特征\((3) PF\perp AB \rightarrow  k\_AB=1\)。 \(AB\) 过 \(F(2,0)\)：\(y=x-2\)，即 \(x-y-2=0\)，故选\(A\)。 验算：\(P(-2,4)\) 代入 \(x-y+6=0\)：\(-2-4+6=0\) 。 （卷③自带详解照录，验算一致。）}
\fi
\end{ansblock}

\begin{ansblock}[2章-拓-课时15-03]
% ans:2章-拓-课时15-03
\ansitem{03}{\(BCD\)}
\ifansbookdetail
\ansnote{详解}{\(A\)：\(DD_{1}\perp \)底面 \(\rightarrow  DN=\sqrt{MN^{2}-MD^{2}}=\sqrt{4-1}=\sqrt{3}\)，\(N\) 的轨迹是以 \(D\) 为圆心、\(\sqrt{3}\) 为半径的圆。 设 \(MN\) 中点 \(P\)、\(MD\) 中点 \(Q\)，则 \(PQ\parallel ND\)、\(PQ=\sqrt{3}/2\)，\(P\) 的轨迹是以 \(Q\) 为圆心 \(\sqrt{3}/2\) 为半径的圆，面积 \(\pi \times 3/4=3\pi /4\neq \pi \)，\(A\) 错； \(B\)：\(N\) 到 \(BB_{1}\) 的距离\(=NB\)（\(BB_{1}\perp \)底面），即 \(N\) 到定点 \(B\) 与定直线 \(DC\) 距离相等（\(B\) 不在 \(DC\) 上），轨迹为抛物线，\(B\) 对； \(C\)：\(AB\parallel D_{1}C_{1}\)，\(D_{1}N\) 与 \(D_{1}C_{1}\) 所成角 \(\pi /3 \rightarrow  D_{1}N\) 为以 \(D_{1}C_{1}\) 为轴、轴截面顶角 \(2\pi /3\) 的圆锥母线，底面 \(ABCD\) 与轴 \(D_{1}C_{1}\) 平行 \(\rightarrow \) 交线为双曲线，\(C\) 对； \(D\)：\(D_{1}N\) 为以 \(D_{1}B\) 为轴、顶角 \(\pi /3\) 的圆锥母线，平面 \(ABCD\) 与轴 \(D_{1}B\) 斜交 \(\rightarrow \) 交线为椭圆，\(D\) 对。 故选 \(BCD\)。 （卷③自带详解照录校订，逐项判定一致。）}
\fi
\end{ansblock}

\begin{ansblock}[2章-拓-课时15-04]
% ans:2章-拓-课时15-04
\ansitem{04}{\(BCD\)}
\ifansbookdetail
\ansnote{详解}{以 \(D\) 为原点、\(DA\)、\(DC\)、\(DD_{1}\) 为 \(x\)、\(y\)、\(z\) 轴建系：\(A(2,0,0)\)，\(C(0,2,0)\)，\(B(2,2,0)\)，\(M(1,0,2)\)。 \(A\)：\(AC=(-2,2,0)\)，\(BM=(-1,-2,2)\)，\(|cos\langle AC,BM\rangle |=|AC\cdot BM|/(|AC||BM|)=|2-4+0|/(2\sqrt{2}\times 3)=2/(6\sqrt{2})=\sqrt{2}/6\neq \sqrt{2}/3\)，\(A\) 错； \(B\)：\(S=\)½\(|AC_{1}|\cdot h=\sqrt{3}\)，\(|AC_{1}|=2\sqrt{3} \rightarrow  h=1\)：\(P\) 在以 \(AC_{1}\) 为轴、半径 \(1\) 的圆柱面上，该圆柱面与侧面 \(BCC_{1}B_{1}\)（与轴斜交）的交线为椭圆的一部分，\(B\) 对； \(C\)：\(P\) 在侧面 \(BCC_{1}B_{1}\) 上时，到 \(C_{1}D_{1}\) 的距离\(=|PC_{1}|\)，条件即 \(|PC_{1}|=P\) 到直线 \(BC\) 的距离——到定点与定直线等距且 \(C_{1}\) 不在 \(BC\) 上 \(\rightarrow \) 抛物线的一部分，\(C\) 对； \(D\)：取 \(AD\) 中点 \(H\)，\(MH=2\) 且 \(MH\perp \)底面；当交线与 \(AC\) 平行（\(HN\perp AC\)）时二面角最小，截面为过 \(M\)、\(N\) 及四棱中点的正六边形，边长 \(\sqrt{2}\)，面积 \(6\times (\sqrt{3}/4)\times 2=3\sqrt{3}\)，\(D\) 对。 故选 \(BCD\)。 （卷③自带详解照录校订，逐项判定一致；\(A\) 项 \(|AC|=2\sqrt{2}\)、\(|BM|=3\)，余弦 \(2/(6\sqrt{2})=\sqrt{2}/6\) 。）}
\fi
\end{ansblock}

\begin{ansblock}[2章-拓-课时15-05]
% ans:2章-拓-课时15-05
\ansitem{05}{\(C\)}
\ifansbookdetail
\ansnote{详解}{折叠不变量：\(PB=AB\)，\(PD=AD \rightarrow  PB+PD=10>BD=8\)，\(P\) 在以 \(B\)、\(D\) 为焦点、长轴 \(10\)、焦距 \(8\) 的椭圆上（不在长轴端点）。 \(\triangle BPD\cong \triangle BCD\)（三边相等）\(\rightarrow  \angle PBD=\angle CBD\)。过 \(P\) 作 \(PE\perp BD\) 于 \(E\)，连 \(CE\)：\(\triangle BPE\cong \triangle BCE \rightarrow  CE\perp BD\) 且 \(PE=CE \rightarrow  BD\perp \)平面 \(PCE\)。 \(V(P-BCD)=V(B-PCE)+V(D-PCE)=(1/3)S\triangle PCE\cdot BD=(8/3)S\triangle PCE\)。 取 \(PC\) 中点 \(F\)，\(EF\perp PC\)（\(PE=CE\)）：\(S\triangle PCE=(1/2)\cdot PC\cdot EF=1\cdot \sqrt{PE^{2}-1}\)（\(PC=2 \rightarrow  CF=1\)）。 \(PE\) 的 \(max=\)椭圆短半轴\(=\sqrt{5^{2}-4^{2}}=3 \rightarrow  S\triangle PCE max=\sqrt{9-1}=2\sqrt{2}\)。 \(V max=(8/3)\times 2\sqrt{2}=16\sqrt{2}/3\)，故选\(C\)。 验算：\(S\triangle PCE=\sqrt{PE^{2}-1}\)：\(EF=\sqrt{PE^{2}-CF^{2}}=\sqrt{PE^{2}-1}\) ；\(PE=3 \rightarrow  EF=2\sqrt{2}\)，\(S=\)½\(\cdot 2\cdot 2\sqrt{2}=2\sqrt{2}\) 。 （卷③自带详解照录，验算一致。）}
\fi
\end{ansblock}

\begin{ansblock}[2章-拓-课时15-06]
% ans:2章-拓-课时15-06
\ansitem{06}{\(BCD\)}
\ifansbookdetail
\ansnote{详解}{过 \(P\) 作 \(PO\perp \)平面 \(ABCD\) 于 \(O\)，\(OH\perp AD\) 于 \(H\)。 \(AD\perp \)平面 \(POH \rightarrow  AD\perp PH \rightarrow  \angle PHO=\alpha \)；又 \(\angle PBO=\beta \)，\(\alpha =\beta  \rightarrow  OH=OB\)。 \(O\) 的轨迹：底面内到定点 \(B\) 与定直线 \(AD\) 等距的点集——以 \(B\) 为焦点、\(AD\) 为准线的抛物线在四边形 \(ABCD\) 内（含边界）的部分；\(P\) 的轨迹为其上平移（以 \(B_{1}\) 为焦点、\(A_{1}D_{1}\) 为准线），是抛物线的一部分而非整条，\(A\) 错； \(B\)：\(|PB|^{2}=(2\sqrt{2})^{2}+|OB|^{2}\)，\(O\) 为 \(AB\) 中点时 \(|OB|min=1 \rightarrow  |PB|min=\sqrt{8+1}=3\)，\(B\) 对； \(C\)：\(PA_{1}\) 与 \(CD\) 所成角即 \(OA\) 与 \(AB\) 所成角 \(\angle OAB\)。以 \(AB\) 中点 \(M\) 为原点建系，抛物线 \(y^{2}=4x\)，\(A(-1,0)\)。设切线 \(x=ty-1\)：代入得 \(y^{2}-4ty+4=0\)，\(\Delta =16t^{2}-16=0 \rightarrow  t=\pm 1\)，取 \(t=1\)（\(O\) 在四边形内）\(\rightarrow  \angle OAB=\pi /4\)，最大值 \(\pi /4\)，\(C\) 对； \(D\)：\(V(P-A_{1}BC_{1})=V(B-PA_{1}C_{1})=(1/3)S\triangle PA_{1}C_{1}\cdot BB_{1}\)，\(A_{1}C_{1}=2\sqrt{2}\)。\(P\) 到 \(A_{1}C_{1}\) 距离最大 \(\Leftrightarrow  O\) 到 \(AC\) 距离最大：\(O\) 为 \(AB\) 中点时最大，值 \((1/4)|BD|=\sqrt{2}/2\)。 \(Vmax=(2\sqrt{2}/3)\times (1/2)\times 2\sqrt{2}\times (\sqrt{2}/2)=2\sqrt{2}/3\)，\(D\) 对。 故选 \(BCD\)。 验算：\(y^{2}=4x\)（焦点 \(B(1,0)\)，准线 \(x=-1=AD\) 所在直线）：\(|OB|\) 与 \(|OH|\) 分别为到焦点、准线距离，等距即轨迹 ；切线 \(y=x+1\) 过 \(A(-1,0)\)，斜率\(1 \rightarrow  \angle OAB=45^{\circ}\) 。 （卷③自带详解照录校订，逐项判定一致。）}
\fi
\end{ansblock}

\begin{ansblock}[2章-拓-课时15-07]
% ans:2章-拓-课时15-07
\ansitem{07}{\((a,0)\)（\(a\neq 0\)）。}
\ifansbookdetail
\ansnote{详解}{按标准形 \(y^{2}=2px\) 对照：\(2p=4a \rightarrow  p=2a\)，焦点 \((p/2,0)=(a,0)\)。 \(a>0\)：开口向右，焦点 \((a,0)\)；\(a<0\)：开口向左，焦点仍为 \((a,0)\)（\(a<0\) 时即 \(x\) 轴负半轴）。 故焦点 \((a,0)\)，\(a\neq 0\)。 验算：\(a=-1\)：\(y^{2}=-4x\)，焦点 \((-1,0)\) 。 （本块自撰亲算；含参开壳，教材未限 \(a\) 正负，按两种开口统一写 \((a,0)\)。）}
\fi
\end{ansblock}

\begin{ansblock}[2章-拓-课时15-08]
% ans:2章-拓-课时15-08
\ansitem{08}{\((1)\) 焦点 \((0,1/8)\)，准线 \(y=-1/8\)；\((2)\) 焦点 \((0,1/(4a))\)，准线 \(y=-1/(4a)\)（\(a\neq 0\)）。}
\ifansbookdetail
\ansnote{详解}{\((1) y=2x^{2} \rightarrow  x^{2}=(1/2)y=2\cdot (1/4)y\)：\(2p=1/2\)，\(p=1/4\)。 开口向上：焦点 \((0,1/8)\)，准线 \(y=-1/8\)。 \((2) y=ax^{2} \rightarrow  x^{2}=(1/a)y=2\cdot (1/(2a))y\)：\(2p=1/a\)，\(p=1/(2a)\)。 \(a>0\)：开口向上，焦点 \((0,1/(4a))\)，准线 \(y=-1/(4a)\)；\(a<0\)：开口向下，焦点、准线表达式不变（\(1/(4a)<0\)，准线在 \(x\) 轴上方）。 验算：\(a=2\) 时与 \((1)\) 一致：\(1/(4\times 2)=1/8\) 。 （本块自撰亲算。）}
\fi
\end{ansblock}

\begin{ansblock}[2章-拓-课时15-09]
% ans:2章-拓-课时15-09
\ansitem{09}{\((1) y=-2x+6\)；\((2) 6m\)。}
\ifansbookdetail
\ansnote{详解}{\((1) |BF|=p/2=1/2\)，灯柱 \(CD\) 过焦点 \(F(1/2,0)\) 且 \(A\) 到灯柱距离 \(1.5 \rightarrow  x\_A=1/2+3/2=2\)。 代入 \(y^{2}=2x\)：\(y\_A=2\)（取上半支），\(A(2,2)\)。 设 \(A\) 处切线 \(y-2=k(x-2)\)：与 \(y^{2}=2x\) 联立消 \(x\) 得 \(ky^{2}-2y+4-4k=0\)，相切 \(\Delta =4-4k(4-4k)=0 \rightarrow  16k^{2}-16k+4=0 \rightarrow  k=1/2\)。 灯罩轴线\(\perp \)切线：斜率 \(-2\)，方程 \(y-2=-2(x-2)\)，即 \(y=-2x+6\)。 \((2)\) 路宽 \(|DH|=10\)，灯罩轴线过路面中线 \(\rightarrow \) 轴线与路面交点到 \(y\) 轴距离 \(1/2+5=11/2\)。 对 \(y=-2x+6\)，\(x=11/2\) 时 \(y=-5 \rightarrow \) 交点 \((11/2,-5)\)，\(|FD|=5\)。 又 \(x=1/2\) 代入 \(y^{2}=2x \rightarrow  y=\pm 1\)，\(|CF|=1\)。 灯柱高 \(|CD|=|DF|+|FC|=5+1=6(m)\)。 验算：切线斜率方程 \(16k^{2}-16k+4=4(2k-1)^{2}=0\) ；\(|CD|=6\) 。}
\fi
\end{ansblock}

\begin{ansblock}[2章-拓-课时15-10]
% ans:2章-拓-课时15-10
\ansitem{10}{\((0,-1/3)\)；\(y=1/3\)。}
\ifansbookdetail
\ansnote{详解}{\(3x^{2}+4y=0 \rightarrow  x^{2}=-(4/3)y=-2\cdot (2/3)y\)：\(2p=4/3\)，\(p=2/3\)。 开口向下：焦点 \((0,-p/2)=(0,-1/3)\)，准线 \(y=p/2=1/3\)。 验算：顶点到焦点、到准线等距 \(1/3\) 。}
\fi
\end{ansblock}

\begin{ansblock}[2章-拓-课时15-11]
% ans:2章-拓-课时15-11
\ansitem{11}{\(x^{2}=8y\)}
\ifansbookdetail
\ansnote{详解}{准线 \(y=-2\) 在 \(x\) 轴下方 \(\rightarrow \) 开口向上：\(x^{2}=2py\)（\(p>0\)），\(-p/2=-2 \rightarrow  p=4\)，\(2p=8\)。 方程 \(x^{2}=8y\)。 验算：焦点 \((0,4)\)、准线 \(y=-2\) 。}
\fi
\end{ansblock}

\begin{ansblock}[2章-拓-课时15-12]
% ans:2章-拓-课时15-12
\ansitem{12}{\(D\)}
\ifansbookdetail
\ansnote{详解}{椭圆 \(x^{2}/2+y^{2}=1\)：\(a^{2}=2\)、\(b^{2}=1 \rightarrow  c^{2}=1\)，焦点 \((\pm 1,0)\)，对称中心 \(O(0,0)\)。 以 \(O\) 为顶点、焦点为 \((1,0)\) 或 \((-1,0)\)：\(p/2=1 \rightarrow  2p=4 \rightarrow  y^{2}=4x\) 或 \(y^{2}=-4x\)，故选\(D\)。 验算：\(y^{2}=\pm 4x\) 的焦点 \((\pm 1,0)\) 。}
\fi
\end{ansblock}

\begin{ansblock}[2章-拓-课时15-13]
% ans:2章-拓-课时15-13
\ansitem{13}{\(D\)}
\ifansbookdetail
\ansnote{详解}{\((2\sqrt{2})^{2}=2px_{0} \rightarrow  x_{0}=4/p\)。 \(M\) 到准线距离\(=x_{0}+p/2=4/p+p/2=3 \rightarrow  p^{2}-6p+8=0 \rightarrow  p=2\) 或 \(p=4\)。 抛物线 \(y^{2}=4x\) 或 \(y^{2}=8x\)，故选\(D\)。 验算：\(p=2\)：\(x_{0}=2\)，\(M(2,2\sqrt{2})\)，到准线 \(x=-1\) 距离 \(3\) ；\(p=4\)：\(x_{0}=1\)，\(M(1,2\sqrt{2})\)，到准线 \(x=-2\) 距离 \(3\) 。}
\fi
\end{ansblock}

\begin{ansblock}[2章-拓-课时15-14]
% ans:2章-拓-课时15-14
\ansitem{14}{答案见解析：\(x\) 的系数为正数且越大时，抛物线的开口向右且开口越大。}
\ifansbookdetail
\ansnote{详解}{三条抛物线同以 \(x\) 轴为对称轴、顶点在原点、开口向右。作图观察（或取同一 \(y\) 值比较横坐标 \(x=y^{2}/(2p)\)）：系数 \(1\)、\(2\)、\(4\) 递增时，同一高度处横坐标 \(1/(2p)\cdot y^{2}\) 递减，曲线横向收拢更快、开口更大。 故：\(x\) 的系数为正且越大，开口向右且开口越大。 （本块自撰亲算补足说理；作图观察题，答案见解析。）}
\fi
\end{ansblock}

\begin{ansblock}[2章-拓-课时15-15]
% ans:2章-拓-课时15-15
\ansitem{15}{\(D\)}
\ifansbookdetail
\ansnote{详解}{由定义，点到焦点距离\(=\)点到准线 \(x=-p/2\) 的距离，其最小值在顶点处取得，为 \(p/2\)。 恒大于 \(1 \Leftrightarrow  p/2>1 \Leftrightarrow  p>2\)，故选\(D\)。 验算：\(p=2\) 时顶点到焦点距离\(=1\)，不满足\("\)恒大于\("\)  开区间正确。}
\fi
\end{ansblock}

\begin{ansblock}[2章-拓-课时16-01]
% ans:2章-拓-课时16-01
\ansitem{01}{\(B\)}
\ifansbookdetail
\ansnote{详解}{点在 \(y^{2}=4x\) 上，故 \(x\geqslant 0\) 且 \(y^{2}=4x\)。代入得 \(z=x^{2}+(1/2)\cdot 4x+3=x^{2}+2x+3=(x+1)^{2}+2\)。由 \(x\geqslant 0\)，\(x=0\) 时 \(z\) 最小，最小值为 \(3\)。故选：\(B\)。〔卷③自带详解照录校订〕}
\fi
\end{ansblock}

\begin{ansblock}[2章-拓-课时16-02]
% ans:2章-拓-课时16-02
\ansitem{02}{\(42\) 或 \(22\)}
\ifansbookdetail
\ansnote{详解}{按 \(M\) 与抛物线的位置分类。\((1) M\) 在抛物线内部或上：\(40^{2}\leqslant 2p\cdot 20\)，即 \(p\geqslant 40\)。设 \(P\) 在准线 \(x=-p/2\) 上的射影为 \(D\)，\(|PF|=|PD|\)，\(M\)、\(P\)、\(D\) 共线时取最小，最小值为 \(M\) 到准线的距离 \(20+p/2=41\)，得 \(p=42\)（\(\geqslant 40\) ）。\((2) M\) 在抛物线外部：\(p<40\)，\(P\)、\(M\)、\(F\) 共线时取最小，最小值为 \(|MF|=\sqrt{}(40^{2}+(20-p/2)^{2})=41\)，\((20-p/2)^{2}=81\)、\(p/2=11\) 或 \(29\)，\(p=22\) 或 \(p=58\)（\(58>40\) 与外部前提矛盾，舍）。综上 \(p=42\) 或 \(22\)。〔卷③自带详解照录校订；内外分类口径随源〕}
\fi
\end{ansblock}

\begin{ansblock}[2章-拓-课时16-03]
% ans:2章-拓-课时16-03
\ansitem{03}{\(C\)}
\ifansbookdetail
\ansnote{详解}{\(x^{2}=4y\) 的焦点 \(F(0,1)\)，准线 \(y=-1\)。延长 \(QP\) 交准线于 \(M\)，则 \(|PM|=y\_P+1=|PQ|+1\)；由定义 \(|PM|=|PF|\)，故 \(|PQ|=|PF|-1\)。于是 \(|PA|+|PQ|=|PA|+|PF|-1\geqslant |AF|-1=\sqrt{}(8^{2}+(7-1)^{2})-1=10-1=9\)，当且仅当 \(A\)、\(P\)、\(F\) 三点共线时取等（\(A(8,7)\) 在抛物线外：\(64>4\cdot 7=28\)，线段 \(AF\) 与抛物线相交，取等可行）。最小值为 \(9\)，选 \(C\)。〔卷③自带详解照录校订〕}
\fi
\end{ansblock}

\begin{ansblock}[2章-拓-课时16-04]
% ans:2章-拓-课时16-04
\ansitem{04}{\(D\)}
\ifansbookdetail
\ansnote{详解}{化标准形 \(x^{2}=-4y\)：焦点 \(F(0,-1)\)，准线 \(l\)：\(y=1\)。第一根式是 \(P\) 到 \(F(0,-1)\) 的距离 \(|PF|\)，第二根式是 \(P\) 到 \(A(4,-5)\) 的距离 \(|PA|\)。由定义 \(|PF|\) 等于 \(P\) 到准线 \(y=1\) 的距离（\(n\leqslant 0\)，即 \(1-n\)），故 \(|PF|+|PA|\) 不小于 \(A\) 到准线的垂线段长 \(1-(-5)=6\)（\(A\) 在抛物线内部：\(4^{2}=16<-4\cdot (-5)=20\) ，\(P\) 取垂线段与抛物线交点即可取等）。最小值 \(6\)，选 \(D\)。〔卷③自带详解照录校订〕}
\fi
\end{ansblock}

\begin{ansblock}[2章-拓-课时16-05]
% ans:2章-拓-课时16-05
\ansitem{05}{\(D\)}
\ifansbookdetail
\ansnote{详解}{\(M(x,y)\) 在以 \(B(6,0)\) 为圆心、\(2\) 为半径的圆上：\((x-6)^{2}+y^{2}=4\)；\(N(x_{0},y_{0})\) 满足 \(y_{0}^{2}=8x_{0}\)（\(x_{0}\geqslant 0\)）。在圆约束下加权配凑：\((1/4)|MA|^{2}=(1/4)[(x-2)^{2}+y^{2}]=(1/4)[(x-2)^{2}+y^{2}]+(3/4)[(x-6)^{2}+y^{2}-4]=(x-5)^{2}+y^{2}\)（所加项在圆上为 \(0\)），即 \((1/2)|MA|=|MC|\)，其中 \(C(5,0)\)。于是 \(|MN|+(1/2)|MA|=|MN|+|MC|\geqslant |CN|=\sqrt{}((x_{0}-5)^{2}+y_{0}^{2})=\sqrt{}((x_{0}-5)^{2}+8x_{0})=\sqrt{}((x_{0}-1)^{2}+24)\geqslant 2\sqrt{6}\)，当且仅当 \(x_{0}=1\)（\(N(1,\pm 2\sqrt{2})\)）且 \(C\)、\(M\)、\(N\) 三点共线（\(M\) 为线段 \(CN\) 与圆的交点）时取等——取等可行：\(C(5,0)\) 在圆内（\(|CB|=1<2\)）而 \(N\) 在圆外（\(|NB|^{2}=(x_{0}-6)^{2}+8x_{0}=(x_{0}-2)^{2}+32>4\) 恒成立），线段 \(CN\) 必与圆相交。最小值 \(2\sqrt{6}\)，选 \(D\)。〔卷③自带详解照录校订；配凑恒等式本件复核：¼\([(x-2)^{2}+y^{2}]-[(x-5)^{2}+y^{2}]=-\)¾\([(x-6)^{2}+y^{2}-4]=0\) 于圆上成立 〕}
\fi
\end{ansblock}

\begin{ansblock}[2章-拓-课时16-06]
% ans:2章-拓-课时16-06
\ansitem{06}{\(B\)}
\ifansbookdetail
\ansnote{详解}{由对称性不妨设 \(P(x,y)\)（\(y>0\)）。特殊位：\(x=1\) 时 \(P(1,2)\)，\(PA\perp x\) 轴，\(tan\angle APB=|AB|/|PA|=2/2=1\)，\(\angle APB=\pi /4\)；\(x=3\) 时 \(P(3,2\sqrt{3})\)，\(PB\perp x\) 轴，\(tan\angle APB=|AB|/|PB|=2/(2\sqrt{3})=\sqrt{3}/3\)，\(\angle APB=\pi /6\)。一般位 \(x\neq 1\) 且 \(x\neq 3\)：\(tan\angle APB=(k\_PB-k\_PA)/(1+k\_PA\cdot k\_PB)\)，其中 \(k\_PA=y/(x-1)\)、\(k\_PB=y/(x-3)\)，代入化简得 \(tan\angle APB=2y/(x^{2}-4x+3+y^{2})\)；以 \(y^{2}=4x\) 代 \(x=y^{2}/4\)：\(tan\angle APB=2y/(y^{4}/16+3)=2/((1/16)y^{3}+3/y)=2/((1/16)y^{3}+1/y+1/y+1/y)\)。分母四项用均值不等式：\((1/16)y^{3}+1/y+1/y+1/y \allowbreak \geqslant  4\cdot ((1/16)y^{3}\cdot (1/y)^{3})^{1/4}=4\cdot (1/16)^{1/4}=2\)，故 \(tan\angle APB\leqslant 2/2=1\)；等号要求 \((1/16)y^{3}=1/y\) 即 \(y=2\)（\(P(1,2)\)，落在 \(x=1\) 特殊位，不在一般位内），故一般位 \(0<tan\angle APB<1\)、\(0<\angle APB<\pi /4\)。综上 \(\angle APB\) 最大值 \(\pi /4\)，选 \(B\)。〔卷③自带详解照录校订；装配注：源解用两直线到角公式，教师版可改向量夹角 \(cos\angle APB=(PA\cdot PB)/(|PA||PB|)\) 路径，结论不变〕}
\fi
\end{ansblock}

\begin{ansblock}[2章-拓-课时16-07]
% ans:2章-拓-课时16-07
\ansitem{07}{\(y^{2}=3x\) 或 \(y^{2}=-3x\)}
\ifansbookdetail
\ansnote{详解}{设方程为 \(y^{2}=2px\) 或 \(y^{2}=-2px\)（\(p>0\)），交点 \(A(x_{1},y_{1})\)、\(B(x_{2},y_{2})\)（\(y_{1}>0>y_{2}\)）。两曲线均关于 \(x\) 轴对称，公共弦被 \(x\) 轴垂直平分：\(x_{1}=x_{2}\)、\(y_{2}=-y_{1}\)，弦长 \(y_{1}-y_{2}=2y_{1}=2\sqrt{3} \Rightarrow  y_{1}=\sqrt{3}\)。代入圆：\(x_{1}^{2}=4-3=1\)、\(x_{1}=\pm 1\)。点 \((1,\sqrt{3})\) 在 \(y^{2}=2px\) 上：\(3=2p\)、\(p=3/2 \Rightarrow  y^{2}=3x\)；点 \((-1,\sqrt{3})\) 在 \(y^{2}=-2px\) 上：同理 \(p=3/2 \Rightarrow  y^{2}=-3x\)。故方程为 \(y^{2}=3x\) 或 \(y^{2}=-3x\)。〔卷③自带详解照录校订〕}
\fi
\end{ansblock}

\begin{ansblock}[2章-拓-课时16-08]
% ans:2章-拓-课时16-08
\ansitem{08}{\(BCD\)}
\ifansbookdetail
\ansnote{详解}{圆 \(C\) 圆心 \((0,1)\)、半径 \(4\)；抛物线 \(x^{2}=4y\) 的焦点恰为 \(F(0,1)\)、准线 \(y=-1\)，圆心与焦点重合。联立 \(x^{2}=4y\) 与 \(x^{2}+(y-1)^{2}=16\)：\(4y+(y-1)^{2}=16 \Rightarrow  y^{2}+2y-15=0 \Rightarrow  y=3\)（\(y=-5\) 舍，\(x^{2}=-20\) 无解），交点 \((\pm 2\sqrt{3},3)\)。劣弧 \(AB\) 为弦 \(y=3\) 上方含点 \((0,5)\) 的短弧（圆心到弦距离 \(2<\)半径 \(4\)），其上动点 \(P\) 纵坐标 \(y\_P\in (3,5)\)，\(A\) 项误以横坐标 \(2\sqrt{3}\) 为下界，错。\(B\)：\(P\)、\(N\) 同横坐标且 \(P\) 在 \(N\) 上方，\(|PN|=y\_P-y\_N\)；由定义 \(|NF|=y\_N+1\)（\(N\) 到准线距离），故 \(|PN|+|NF|=y\_P+1\)，恰为 \(P\) 到准线 \(y=-1\) 的距离，对（特位 \(x=0\)：\(P(0,4)\)、\(N(0,0)\)，\(4+1=5=y\_P+1\) ）。\(C\)：圆心 \((0,1)\) 到准线 \(y=-1\) 距离为 \(2\)，对。\(D\)：\(|PF|=\)半径 \(4\)，\(\triangle PFN\) 周长\(=|PF|+|PN|+|NF|=4+(y\_P-y\_N)+(y\_N+1)=y\_P+5\in (8,10)\)，对。故选 \(BCD\)。〔卷③自带详解照录校订；装配图位：源含简图〕}
\fi
\end{ansblock}

\begin{ansblock}[2章-拓-课时16-09]
% ans:2章-拓-课时16-09
\ansitem{09}{\((0,0)\)（最小值 \(8\)）}
\ifansbookdetail
\ansnote{详解}{设 \(P(x_{0},y_{0})\)，\(y_{0}^{2}=-4x_{0}\)（\(x_{0}\leqslant 0\)）。\(AP\cdot BP=(x_{0}-2)(x_{0}-4)+y_{0}^{2}=x_{0}^{2}-6x_{0}+8-4x_{0}=x_{0}^{2}-10x_{0}+8=(x_{0}-5)^{2}-17\)。由 \(x_{0}\leqslant 0\)，\((x_{0}-5)^{2}\geqslant 25\) 且随 \(x_{0}\) 增大而减小，故 \(x_{0}=0\) 时 \(AP\cdot BP\) 最小，最小值 \(8\)，此时 \(y_{0}=0\)，点 \(P(0,0)\)。〔卷③自带详解照录校订〕}
\fi
\end{ansblock}

\begin{ansblock}[2章-拓-课时16-10]
% ans:2章-拓-课时16-10
\ansitem{10}{\(ABD\)}
\ifansbookdetail
\ansnote{详解}{\(A\)：由定义，\(|PF|\) 等于 \(P\) 到准线 \(l\) 的距离，而 \(PE\perp l\) 于 \(E\)，故 \(|PE|=|PF|\)，对。\(B\)：\(PQ\) 平分 \(\angle EPF\) 的外角 \(\Rightarrow  \angle FPQ=\angle NPQ\)；又 \(EN\parallel x\) 轴（同垂直于准线）\(\Rightarrow  \angle NPQ=\angle PQF\)（内错角），得 \(\angle FPQ=\angle PQF\)，\(\triangle FPQ\) 中 \(|PF|=|QF|\)，对。\(C\)：若 \(|PN|=|MF|\)，又 \(|QM|=|QN|\)（\(Q\) 在外角平分线上，到角两边等距）、\(\angle QMF=\angle QNP=90^{\circ}\)，则 \(\triangle FMQ\cong \triangle PNQ\)，得 \(|FQ|=|PQ|\)，结合 \(B\) 得 \(\triangle PFQ\) 为正三角形、\(\angle PFQ=\pi /3\)，则 \(P\) 被锁定为定点，与「\(P\) 为异于 \(O\) 的任意一点」矛盾，\(C\) 不恒成立。\(D\)：四边形 \(KQNE\) 为矩形（\(\angle KEQ=\angle EQN=\angle K=90^{\circ}\)）\(\Rightarrow  |EN|=|KQ|\)；\(|PN|=|EN|-|EP|=|KQ|-|PF|=|KF|+|FQ|-|PF|=|KF|\)（用 \(|EP|=|PF|=|FQ|\)），对。故选 \(ABD\)。〔卷③自带详解照录校订；装配图位：题面图嵌于源；\(D\) 的线段加减链已按源补全〕}
\fi
\end{ansblock}

\begin{ansblock}[2章-拓-课时16-11]
% ans:2章-拓-课时16-11
\ansitem{11}{\(ABC\)}
\ifansbookdetail
\ansnote{详解}{准线 \(m\) 交 \(x\) 轴于 \(P\)，\(|PF|=p\)。分别过 \(A\)、\(B\) 作准线的垂线，垂足 \(E\)、\(M\)。倾斜角 \(60^{\circ}\)：\(AE\parallel x\) 轴 \(\Rightarrow  \angle EAF=60^{\circ}\)；由定义 \(|AE|=|AF| \Rightarrow  \triangle AEF\) 为等边三角形 \(\Rightarrow  \angle EFP=\angle AEF=60^{\circ}\)、\(\angle FEP=30^{\circ}\)（直角\(\triangle EPF\) 中），于是 \(|AF|=|EF|=2|PF|=2p=8\)，\(p=4\)，\(A\) 对。又 \(|AE|=|EF|=2|PF|\) 且 \(PF\parallel AE \Rightarrow  F\) 为 \(AD\) 中点，\(|DF|=|FA|\)，\(B\) 对。\(\angle DAE=60^{\circ}\) 的补算：\(\triangle ADE\) 中 \(\angle AED=90^{\circ}\)、\(\angle ADE=30^{\circ} \Rightarrow  |BD|=2|BM|=2|BF|\)（末步用定义 \(|BM|=|BF|\)），\(C\) 对。由 \(|BD|=2|BF|\) 得 \(|BF|=(1/3)|DF|=(1/3)|AF|=8/3\neq 4\)，\(D\) 错。故选 \(ABC\)。〔卷③自带详解照录校订；装配图位：源含简图〕}
\fi
\end{ansblock}

\begin{ansblock}[2章-拓-课时16-12]
% ans:2章-拓-课时16-12
\ansitem{12}{\(A\)}
\ifansbookdetail
\ansnote{详解}{圆化标准形：\((x-1)^{2}+y^{2}=1\)，圆心 \(F(1,0)\)、半径 \(1\)，恰为抛物线 \(y^{2}=4x\) 的焦点。设 \(A(x_{1},y_{1})\)、\(D(x_{2},y_{2})\)（\(A\)、\(D\) 为直线与抛物线的两交点，\(B\)、\(C\) 为与圆的两交点，自上而下顺次），由焦半径 \(|AF|=x_{1}+1\)、\(|DF|=x_{2}+1\)，得 \(|AB|=|AF|-|BF|=x_{1}+1-1=x_{1}\)，\(|CD|=|DF|-|CF|=x_{2}\)。设 \(l\)：\(x=my+1\) 代入 \(y^{2}=4x\) 得 \(y^{2}-4my-4=0\)，\(y_{1}y_{2}=-4\)，故 \(|AB|\cdot |CD|=x_{1}x_{2}=(y_{1}^{2}/4)(y_{2}^{2}/4)=(y_{1}y_{2})^{2}/16=16/16=1\)，为定值 \(1\)（与 \(m\) 无关，最值之说随之落空）。选 \(A\)。〔卷③自带详解照录校订；源用「平均性质」名目，按 §五 口径已改写为常规联立韦达，答案不变〕}
\fi
\end{ansblock}

\begin{ansblock}[2章-拓-课时16-13]
% ans:2章-拓-课时16-13
\ansitem{13}{\(B\)（\(-4\)）}
\ifansbookdetail
\ansnote{详解}{联立 \(x^{2}=4(kx+2)\)：\(x^{2}-4kx-8=0\)（\(\Delta =16k^{2}+32>0\) 恒成立），\(x_{1}+x_{2}=4k\)、\(x_{1}x_{2}=-8\)。\(y_{1}y_{2}=(kx_{1}+2)(kx_{2}+2)=k^{2}x_{1}x_{2}+2k(x_{1}+x_{2})+4=-8k^{2}+8k^{2}+4=4\)（即过 \((0,2)\) 的弦两端纵坐标之积为定值 \(4\)，\(A\)、\(B\) 分居 \(y\) 轴两侧故 \(x_{1}x_{2}<0\) 与 \(-8\) 相符）。\(OA\cdot OB=x_{1}x_{2}+y_{1}y_{2}=-8+4=-4\)，与 \(k\) 无关。值 \(-4\) 对应选项 \(B\)，选 \(B\)。〔校订注：卷③此题面之直线、抛物线方程及选项 \(A\)、\(B\) 数值在原 \(dump\) 中均为图嵌（【图】占位，图内文本未入对象层），题方按源【编注】【分析】【详解】读数 \(y=kx+2\)、\(x^{2}=4y\) 与 \(OA\cdot OB=-4\) 补全（联立恒等于源算路复核 ，见本块推导）；选项 \(A\) 数值未能回，装配回源图或改选项；「平均性质」名目已按 §五 改常规联立〕}
\fi
\end{ansblock}

\begin{ansblock}[2章-拓-课时16-14]
% ans:2章-拓-课时16-14
\ansitem{14}{\(AD\)}
\ifansbookdetail
\ansnote{详解}{抛物线 \(x^{2}=2y\) 的焦点 \((0,1/2)\)、准线 \(y=-1/2\)。设 \(l\)：\(y=kx+1/2\)，代入 \(x^{2}=2y\) 得 \(x^{2}-2kx-1=0\)，\(x_{1}+x_{2}=2k\)、\(x_{1}x_{2}=-1\)；\(y_{1}+y_{2}=k(x_{1}+x_{2})+1=2k^{2}+1\)、\(y_{1}y_{2}=(kx_{1}+1/2)(kx_{2}+1/2)=k^{2}x_{1}x_{2}+(k/2)(x_{1}+x_{2})+1/4=-k^{2}+k^{2}+1/4=1/4\)。\(A\)：过点 \((1,1/2)\)、斜率 \(t\) 的直线 \(y=t(x-1)+1/2\) 代入 \(x^{2}=2y\) 得 \(x^{2}-2tx+2t-1=0\)，即 \((x-1)(x-(2t-1))=0\)，另一交点横标 \(2t-1\)；相切 \(\Leftrightarrow  2t-1=1 \Leftrightarrow  t=1\)，切线 \(y=x-1/2\) 即 \(2x-2y-1=0\)，\(A\) 对。\(B\)：\(AM\cdot BM=(-x_{1},-1/2-y_{1})\cdot (-x_{2},-1/2-y_{2})=x_{1}x_{2}+(1/2+y_{1})(1/2+y_{2})=x_{1}x_{2}+(1/4)+(1/2)(y_{1}+y_{2})+y_{1}y_{2}=-1+1/4+(1/2)(2k^{2}+1)+1/4=k^{2}\)，仅 \(k=0\)（通径位）时为 \(0\)，不恒成立，\(B\) 错。\(C\)：取 \(l\parallel x\) 轴（\(k=0\)，即通径），\(A(1,1/2)\)、\(B(-1,1/2)\)，\(|OA|=|OB|=\sqrt{1+1/4}=\sqrt{5}/2\)，\(|OA|+|OB|=\sqrt{5}\)，严格不等式不成立，\(C\) 错。\(D\)：\(BN\perp \)准线 \(\Rightarrow  N(x_{2},-1/2)\)；直线 \(OA\)：\(y=(y_{1}/x_{1})x=(x_{1}/2)x\)（用 \(y_{1}=x_{1}^{2}/2\)），代入 \(x=x_{2}\) 得纵标 \((x_{1}x_{2})/2=-1/2\)，即直线 \(OA\) 恰过 \(N(x_{2},-1/2)\)，\(A\)、\(O\)、\(N\) 共线，\(D\) 对。故选 \(AD\)。〔卷③自带详解照录校订；\(A\) 项源解用导数 \(y'=x\)（系选修\(2\) 工具），已消前引改「联立重根」路径，结论一致；\(C\) 项源 \(dump\) 排印「\(B(1,-1/2)\)」系讹（该点不在抛物线上），依 \(k=0\) 弦 \(x^{2}=1\) 校正为 \(B(-1,1/2)\)〕}
\fi
\end{ansblock}

\begin{ansblock}[2章-拓-课时16-15]
% ans:2章-拓-课时16-15
\ansitem{15}{\(y^{2}=(p/2)x\)（\(x\geqslant 0\)），轨迹是顶点在原点、对称轴为 \(x\) 轴、开口向右的抛物线（\(2\tilde p \allowbreak =p/2\)，焦点 \((p/8,0)\)）}
\ifansbookdetail
\ansnote{详解}{设抛物线上点 \((x_{0},y_{0})\)，\(y_{0}^{2}=2px_{0}\)（\(x_{0}\geqslant 0\)），垂线段中点 \((x,y)=(x_{0}, y_{0}/2)\)。消参：\(x_{0}=x\)、\(y_{0}=2y\) 代入得 \(4y^{2}=2px\)，即 \(y^{2}=(p/2)x\)（\(x\geqslant 0\)）。这是顶点在原点、对称轴为 \(x\) 轴、开口向右的抛物线，其 \(2\tilde p \allowbreak =p/2\)（\(\tilde p \allowbreak =p/4\)），焦点 \((p/8,0)\)。（原点处垂线段退化为点 \(O\)，中点仍为 \(O\)，轨迹含顶点 。）〔亲算印证：消参直推，读数自验〕}
\fi
\end{ansblock}

\begin{ansblock}[2章-拓-课时16-16]
% ans:2章-拓-课时16-16
\ansitem{16}{\(y^{2}=8x\)}
\ifansbookdetail
\ansnote{详解}{点 \(A\) 横坐标 \(2>0\) 在 \(y\) 轴右侧，故抛物线开口向右，设 \(y^{2}=2px\)（\(p>0\)），\(F(p/2,0)\)。设 \(A(2,y\_A)\)，\(y\_A^{2}=4p\)。\(FA\rightarrow =(2-p/2, y\_A)\)、\(OA\rightarrow =(2, y\_A)\)，\(FA\cdot OA=2(2-p/2)+y\_A^{2}=4-p+4p=4+3p=16\)，解得 \(p=4\)，抛物线为 \(y^{2}=8x\)。（回验：\(A(2,\pm 4)\)，\(FA\cdot OA=2(2-2)+16=16\) 。）〔亲算印证：自验读数 \(y^{2}=8x\)；教材答案条乱码段（\(raw\) 行\(4554\)起）未取，与 §三\(.1\) 口径一致〕}
\fi
\end{ansblock}

\begin{ansblock}[2章-拓-课时16-17]
% ans:2章-拓-课时16-17
\ansitem{17}{证明见详解；类似结论：\(y^{2}=-2px\) 上 \(|MF|=p/2-x_{0}\)（\(x_{0}\leqslant 0\)）；\(x^{2}=2py\) 上 \(|MF|=y_{0}+p/2\)；\(x^{2}=-2py\) 上 \(|MF|=p/2-y_{0}\)}
\ifansbookdetail
\ansnote{详解}{证明：准线 \(l\)：\(x=-p/2\)。由抛物线定义，\(|MF|\) 等于 \(M\) 到 \(l\) 的距离；\(M\) 横坐标 \(x_{0}\geqslant 0\)，故距离为 \(x_{0}-(-p/2)=x_{0}+p/2\)。\(\rule{1.05ex}{1.05ex}\) 类似结论（同法：焦半径\(=\)到准线距离，注意坐标符号范围）：①\(y^{2}=-2px\)（\(x_{0}\leqslant 0\)，准线 \(x=p/2\)）：\(|MF|=p/2-x_{0}\)；②\(x^{2}=2py\)（\(y_{0}\geqslant 0\)，准线 \(y=-p/2\)）：\(|MF|=y_{0}+p/2\)；③\(x^{2}=-2py\)（\(y_{0}\leqslant 0\)，准线 \(y=p/2\)）：\(|MF|=p/2-y_{0}\)。统一记法：焦半径\(=p/2+\)（沿开口方向的点到顶点有向距离）。〔亲算印证：定义直推，与 §三\(.1\)「\(|MF|=|y_{0}|+p/2\)」型读数同族〕}
\fi
\end{ansblock}

\begin{ansblock}[2章-拓-课时16-19]
% ans:2章-拓-课时16-19
\ansitem{19}{\(B\)}
\ifansbookdetail
\ansnote{详解}{连接 \(PF\)。由抛物线定义，\(P\) 到焦点的距离等于 \(P\) 到准线的距离，即 \(|PF|=|PQ|\)，\(\triangle QPF\) 为等腰三角形，\(P\) 到线段 \(FQ\) 两端点等距，故 \(P\) 在线段 \(FQ\) 的垂直平分线上，选 \(B\)。（反例排 \(A\)：取 \(y^{2}=4x\)、\(P(1,2)\)，则 \(F(1,0)\)、\(Q(-1,2)\)，\(|OF|=1\) 而 \(|OQ|=\sqrt{5}\)，\(O\) 不在垂直平分线上；\(C\)、\(D\) 随之不恒成立。）〔亲算印证：定义直推，与轮\(2\)亲算④抽样档「过 \(P\)，选\(B\)」一致〕}
\fi
\end{ansblock}

\begin{ansblock}[2章-拓-课时16-20]
% ans:2章-拓-课时16-20
\ansitem{20}{\(2\sqrt{3}\)}
\ifansbookdetail
\ansnote{详解}{椭圆左顶点 \((-4,0)\)。设 \(P(-y^{2}/4,y)\)，则 \(d^{2}=(-y^{2}/4+4)^{2}+y^{2}=(1/16)y^{4}-y^{2}+16=(1/16)(y^{2}-8)^{2}+12\geqslant 12\)，当 \(y^{2}=8\)（\(y=\pm 2\sqrt{2}\)，\(x=-2\)，点 \(P(-2,\pm 2\sqrt{2})\) 在抛物线上 ）时取等，\(d\) 的最小值为 \(\sqrt{12}=2\sqrt{3}\)。〔亲算印证：轮\(2\)亲算④抽样档 \(2\sqrt{3}\) 与 \(3\)章件\(14-\#10\) 源答案一致〕}
\fi
\end{ansblock}

\begin{ansblock}[2章-拓-课时16-21]
% ans:2章-拓-课时16-21
\ansitem{21}{\(7/2\)；\(P(2,2)\)}
\ifansbookdetail
\ansnote{详解}{\(A(3,2)\)：\(2^{2}=4<2\cdot 3=6\)，\(A\) 在抛物线内部。准线 \(l\)：\(x=-1/2\)。设 \(P\) 在 \(l\) 上的射影为 \(Q\)，由定义 \(|PF|=|PQ|\)，\(|PA|+|PF|=|PA|+|PQ|\geqslant A\) 到准线的水平距离\(=3-(-1/2)=7/2\)，当 \(A\)、\(P\)、\(Q\) 共线（\(y\_P=y\_A=2\)）时取等。代入 \(y^{2}=2x\) 得 \(x=2\)，即 \(P(2,2)\)，最小值 \(7/2\)。〔亲算印证：轮\(2\)亲算④抽样档「\(7/2\)、\(P(2,2)\)」与源答案一致〕}
\fi
\end{ansblock}

\begin{ansblock}[2章-拓-课时16-22]
% ans:2章-拓-课时16-22
\ansitem{22}{\(D\)}
\ifansbookdetail
\ansnote{详解}{\(P\) 在抛物线上：\(16=4m\)、\(m=4\)，\(P(4,-4)\)。焦半径（定义）：准线 \(x=-1\)，\(|PF|=x\_P+1=4+1=5\)。（两点距复核：\(F(1,0)\)，\(|PF|=\sqrt{9+16}=5\) 。）选 \(D\)。〔亲算印证：轮\(2\)亲算④抽样档「\(5\)，选\(D\)」与源答案一致〕}
\fi
\end{ansblock}

\begin{ansblock}[2章-拓-课时16-23]
% ans:2章-拓-课时16-23
\ansitem{23}{\(y^{2}=12x\)}
\ifansbookdetail
\ansnote{详解}{设圆 \(M\) 半径为 \(R\)。\(\odot C\) 圆心 \(C(3,0)\)、半径 \(1\)，外切 \(\Rightarrow  |MC|=R+1\)；\(\odot M\) 与直线 \(x=-2\) 相切 \(\Rightarrow \) 圆心 \(M\) 到 \(x=-2\) 的距离为 \(R\)，从而 \(M\) 到直线 \(x=-3\) 的距离为 \(R+1\)。于是动点 \(M\) 到定点 \(C(3,0)\) 的距离等于它到定直线 \(x=-3\) 的距离，轨迹是以 \(C\) 为焦点、\(x=-3\) 为准线的抛物线（定义）：\(p/2=3\)、\(p=6\)，标准方程 \(y^{2}=12x\)。（核验 \(x\geqslant 0\)：\(M\) 到 \(x=-3\) 距离\(=\)横坐标\(+3=R+1>0\) 且 \(|MC|=R+1\geqslant 1>0\) 自洽。）〔亲算印证：轮\(2\)亲算④抽样档 \(y^{2}=12x\) 与源答案一致〕}
\fi
\end{ansblock}

\begin{ansblock}[2章-拓-课时16-24]
% ans:2章-拓-课时16-24
\ansitem{24}{\(C\)}
\ifansbookdetail
\ansnote{详解}{\(|OA|=|OB|\)：\(x_{1}^{2}+y_{1}^{2}=x_{2}^{2}+y_{2}^{2}\)，而 \(y\_i=x\_i^{2}/2\)，\((x_{1}^{2}-x_{2}^{2})(1+(x_{1}^{2}+x_{2}^{2})/4)=0 \Rightarrow  x_{1}^{2}=x_{2}^{2}\)，\(A\)、\(B\) 异点故 \(x_{2}=-x_{1}\)，两点关于 \(y\) 轴对称。设 \(B(a,a^{2}/2)\)、\(A(-a,a^{2}/2)\)（\(a>0\)），\(S=\)½\(\cdot 2a\cdot (a^{2}/2)=a^{3}/2=12\sqrt{3} \Rightarrow  a^{3}=24\sqrt{3}=(2\sqrt{3})^{3} \Rightarrow  a=2\sqrt{3}\)，\(B(2\sqrt{3},6)\)。记 \(OB\) 与 \(y\) 轴夹角 \(\theta \)：\(tan\theta =a/(a^{2}/2)=2\sqrt{3}/6=\sqrt{3}/3 \Rightarrow  \theta =30^{\circ}\)，\(\angle AOB=2\theta =60^{\circ}\)。选 \(C\)。〔亲算印证：轮\(2\)亲算④抽样档「\(60^{\circ}\)，选\(C\)」与源答案一致〕}
\fi
\end{ansblock}

\begin{ansblock}[2章-拓-课时16-25]
% ans:2章-拓-课时16-25
\ansitem{25}{\(7\sqrt{3}/12\)}
\ifansbookdetail
\ansnote{详解}{等腰梯形两底 \(AB\parallel CD\)，而抛物线上平行于 \(x\) 轴方向的弦（两交点同纵坐标）不存在（\(y^{2}=2px\) 单值），故两底必为垂直于 \(x\) 轴的弦：\(A\)、\(B\) 关于 \(x\) 轴对称，\(C\)、\(D\) 关于 \(x\) 轴对称。\(|AB|=2 \Rightarrow  y\_A=1\)（取上半），\(|CD|=4 \Rightarrow  y\_D=2\)。设 \(A(m,1)\)；腰与底夹角 \(\angle ADC=60^{\circ}\)，底铅直，故 \(DA\) 与铅直线夹 \(60^{\circ}\)，纵差 \(2-1=1 \Rightarrow \) 横差\(=tan60^{\circ}=\sqrt{3}\)，\(D(m+\sqrt{3},2)\)。两点在抛物线上：\(1=2pm\)、\(4=2p(m+\sqrt{3})\)，两式相减 \(3=2\sqrt{3} p \Rightarrow  p=\sqrt{3}/2\)、\(m=1/(2p)=\sqrt{3}/3\)。由焦半径（拓\(16-17\) 结论）\(|AF|=m+p/2=\sqrt{3}/3+\sqrt{3}/4=7\sqrt{3}/12\)。〔亲算印证：轮\(2\)亲算②\(-1\) 关键步「\(p=\sqrt{3}/2\)、\(m=\sqrt{3}/3\)、\(|AF|=7\sqrt{3}/12\)」与 \(3\)章件\(15-\#14\) 源答案一致〕}
\fi
\end{ansblock}

\begin{ansblock}[2章-拓-课时16-26]
% ans:2章-拓-课时16-26
\ansitem{26}{\(A\)}
\ifansbookdetail
\ansnote{详解}{\(F(1,0)\)。设 \(l\)：\(x=ty+1\)（含斜率不存在），代入 \(y^{2}=4x\)：\(y^{2}-4ty-4=0\)，\(y_{1}+y_{2}=4t\)、\(y_{1}y_{2}=-4\)，\(x_{1}+x_{2}=t(y_{1}+y_{2})+2=4t^{2}+2\)、\(x_{1}x_{2}=(y_{1}y_{2})^{2}/16=1\)。\(\angle AMB\) 为锐角 \(\Leftrightarrow  MA\cdot MB>0\)（且 \(A\)、\(M\)、\(B\) 不共线，\(M\) 在 \(x\) 轴负半轴而弦过 \(F\) 恒成立）：\(MA\cdot MB=(x_{1}-m)(x_{2}-m)+y_{1}y_{2}=x_{1}x_{2}-m(x_{1}+x_{2})+m^{2}-4=m^{2}-2m-3-4mt^{2}\)。题源解答按「对一切过 \(F\) 的弦（\(\forall t\in R\)）成立」求解：\(m<0\) 时 \(-4mt^{2}\geqslant 0\) 且随 \(t^{2}\) 递增，最严处 \(t=0\)，须 \(m^{2}-2m-3>0 \Rightarrow  m<-1\) 或 \(m>3\)，配 \(m<0\) 得 \(m\in (-\infty ,-1)\)，选 \(A\)。〔措辞存疑：题面「存在 \(M\) 使 \(\angle AMB\) 为锐角」若按「对某条弦成立」解读，则一切 \(m<0\) 皆可（\(t\) 足够大即成立），与选项无合——故题面取「对任意焦点弦恒可为锐角」义，与源答案 \(A\) 一致，已登本件存疑清单〕〔亲算印证：轮\(2\)亲算②\(-1\)「\((-\infty ,-1)\)，选\(A\)」与 \(3\)章件\(18-\#1\) 源答案一致〕}
\fi
\end{ansblock}

\begin{ansblock}[2章-拓-课时16-27]
% ans:2章-拓-课时16-27
\ansitem{27}{\(B\)}
\ifansbookdetail
\ansnote{详解}{\(p=4\)，准线 \(x=-2\)。分别过 \(A\)、\(M\)、\(B\) 作准线的垂线，垂足 \(C\)、\(D\)、\(E\)，则 \(MD\) 为梯形（退化情形为中点公式仍合）\(ACEB\) 的中位线，\(|MD|=(|AC|+|BE|)/2=(|AF|+|BF|)/2=|AB|/2\)（定义）。设 \(l\)：\(x=my+2\) 代入 \(y^{2}=8x\)：\(y^{2}-8my-16=0\)，\(|AB|=\sqrt{1+m^{2}}|y_{1}-y_{2}|=\sqrt{1+m^{2}}\cdot \sqrt{64m^{2}+64}=8(1+m^{2})\geqslant 8\)（\(m=0\) 即通径取等）。故 \(|MD|=|AB|/2\geqslant 4\)，最小值 \(4\)，选 \(B\)。〔亲算印证：轮\(2\)亲算②\(-1\)「\(\geqslant \)½通径\(=4\)，选\(B\)」与 \(3\)章件\(18-\#2\) 源答案一致〕}
\fi
\end{ansblock}

\begin{ansblock}[2章-拓-课时16-28]
% ans:2章-拓-课时16-28
\ansitem{28}{\(A\)}
\ifansbookdetail
\ansnote{详解}{由定义 \(d_{1}=|MF|\)；设 \(MN\perp l'\) 于 \(N\)，则 \(d_{2}=|MN|\)，\(d_{1}+d_{2}=|MF|+|MN|\geqslant d(F,l')\)（\(F\) 到 \(l'\) 的垂线段长），当 \(M\) 落在「过 \(F\) 垂直于 \(l'\) 的直线」与抛物线的交点上（该垂线与 \(l'\) 的垂足 \(N\) 在 \(F\) 同侧射影之间，可行）时取最小。\(d(F,l')=|1-0+2|/\sqrt{2}=3\sqrt{2}/2\)。设 \(l'\) 与 \(x\) 轴交于 \(D(-2,0)\)，\(|FD|=3\)；在 \(Rt\triangle DNF\)（\(\angle DNF=90^{\circ}\)）中 \(cos\angle NFD=|FN|/|FD|=(3\sqrt{2}/2)/3=\sqrt{2}/2\)；取最小值时 \(M\) 在线段 \(FN\) 上、射线 \(FO\) 与 \(FD\) 同向，\(cos\angle MFO=\sqrt{2}/2\)，选 \(A\)。（\(M\) 坐标核：\(M(t^{2},2t)\) 且 \(FM\parallel (1,-1)\) 法向 \(\rightarrow  t=\sqrt{2}-1\)，\(M(3-2\sqrt{2},2\sqrt{2}-2)\) 在 \(FN\) 段内 。）〔亲算印证：轮\(2\)亲算②\(-1\)「\(cos\angle MFO=\sqrt{2}/2\)，选\(A\)」与 \(3\)章件\(18-\#7\) 源答案一致；装配图位：源解含图〕}
\fi
\end{ansblock}

\begin{ansblock}[2章-拓-课时16-29]
% ans:2章-拓-课时16-29
\ansitem{29}{\((1) y^{2}=4x\)；\((2)\) 直线 \(AB\) 恒过定点 \((-1,0)\)}
\ifansbookdetail
\ansnote{详解}{\((1)\) 设圆心 \(M(x,y)\)，圆过 \((2,0) \Rightarrow \) 半径\(^{2}=(x-2)^{2}+y^{2}\)；被 \(y\) 轴截弦长 \(4 \Rightarrow \) 半径\(^{2}=x^{2}+2^{2}\)（半弦 \(2\)、弦心距 \(|x|\)）。两式相等：\((x-2)^{2}+y^{2}=x^{2}+4 \Rightarrow  y^{2}=4x\)。 \((2)\) 顶点在轨迹 \(y^{2}=4x\) 上。设 \(B(x_{1},y_{1})\)、\(C(x_{2},y_{2})\)，\(BC\) 过 \(F(1,0)\) 且不垂直于坐标轴，设 \(BC\)：\(x=ty+1\)：\(y^{2}-4ty-4=0 \Rightarrow  y_{1}y_{2}=-4\)。由 \(A\)、\(C\) 关于 \(x\) 轴对称，\(A(x_{2},-y_{2})\)。设 \(AB\)：\(y=kx+m\)（\(k\neq 0\)）代入 \(ky^{2}-4y+4m=0\)，其两根为 \(y_{1}\)、\(-y_{2}\)：\(y_{1}\cdot (-y_{2})=4m/k\)，即 \(4=4m/k\)、\(m=k\)。故 \(AB\)：\(y=k(x+1)\)，对一切实有 \(k\) 恒过定点 \((-1,0)\)。\(\rule{1.05ex}{1.05ex}\)（源答案栏「证明见解析」未落定点数值，定点 \((-1,0)\) 系亲算补出，登轮\(2\)亲算⑤\(-5\)。）〔亲算印证：轮\(2\)亲算②\(-1\) 参数点法同得 \((-1,0)\)；本块 \((2)\) 采源韦达路径改写（消「见解析」）〕}
\fi
\end{ansblock}

\begin{ansblock}[2章-拓-课时16-30]
% ans:2章-拓-课时16-30
\ansitem{30}{\((1) x^{2}=4y\)；\((2)\) 点 \(P\) 恒在定直线 \(y=-1\)（即 \(C\) 的准线）上}
\ifansbookdetail
\ansnote{详解}{\((1) A\) 横坐标 \(2\sqrt{2}\)、在圆上：\(y\_A=(12-8)/2=2\)，\(A(2\sqrt{2},2)\)。代入 \(x^{2}=2py\)：\(8=4p\)、\(p=2\)，抛物线为 \(x^{2}=4y\)。 \((2)\) 先给切线式并验证：点 \(M(x_{1},x_{1}^{2}/4)\) 处，直线 \(y=(x_{1}/2)x-x_{1}^{2}/4\) 与 \(x^{2}=4y\) 联立得 \(x^{2}-2x_{1}x+x_{1}^{2}=(x-x_{1})^{2}=0\)，唯一公共点（重根 \(x_{1}\)），即为过 \(M\) 的切线 \(l_{1}\)；同理 \(N(x_{2},x_{2}^{2}/4)\) 处切线 \(l_{2}\)：\(y=(x_{2}/2)x-x_{2}^{2}/4\)。两切线联立：\((x_{1}-x_{2})x/2=(x_{1}^{2}-x_{2}^{2})/4 \Rightarrow  x_{0}=(x_{1}+x_{2})/2\)，\(y_{0}=(x_{1}/2)x_{0}-x_{1}^{2}/4=x_{1}x_{2}/4\)。设 \(l\)：\(y=kx+1\)（过 \(F(0,1)\)；\(l\) 为 \(x=0\) 时与抛物线仅一交点，不合「不同两点」），代入 \(x^{2}=4y\)：\(x^{2}-4kx-4=0 \Rightarrow  x_{1}x_{2}=-4\)，故 \(y_{0}=-1\) 恒成立，点 \(P\) 恒在定直线 \(y=-1\)（准线）上。\(\rule{1.05ex}{1.05ex}\)〔前引注：源解以导数 \(y'=x/2\) 写切线（系选修\(2\) 工具），本区消前引改「重根验证」路径，切线式与结论同源一致；\((2)\) 源答案「证明见解析」，定直线 \(y=-1\) 系亲算②\(-1\) 落值（⑤\(-5\)）〕}
\fi
\end{ansblock}

\begin{ansblock}[2章-拓-课时17-01]
% ans:2章-拓-课时17-01
\ansitem{01}{\(C\)（弦长 \(8/5\)）}
\ifansbookdetail
\ansnote{详解}{\(a^{2}=4\)、\(b^{2}=1\)，\(c^{2}=a^{2}-b^{2}=3\)、\(c=\sqrt{3}\)，右焦点 \((\sqrt{3},0)\)。直线 \(l\)：\(y=x-\sqrt{3}\)。代入 \(x^{2}/4+y^{2}=1\)：\(x^{2}/4+(x-\sqrt{3})^{2}=1\)，乘 \(4\) 得 \(x^{2}+4(x^{2}-2\sqrt{3}x+3)=4\)，即 \(5x^{2}-8\sqrt{3}x+8=0\)。\(\Delta =(8\sqrt{3})^{2}-4\cdot 5\cdot 8=192-160=32>0\)，\(x_{1}+x_{2}=8\sqrt{3}/5\)、\(x_{1}x_{2}=8/5\)。弦长 \(|AB|=\sqrt{1+k^{2}}\cdot |x_{1}-x_{2}|=\sqrt{2}\cdot \sqrt{\Delta }/5=\sqrt{2}\cdot \sqrt{32}/5=\sqrt{2}\cdot 4\sqrt{2}/5=8/5\)。（互核：\(|x_{1}-x_{2}|^{2}=(x_{1}+x_{2})^{2}-4x_{1}x_{2}=192/25-32/5=192/25-160/25=32/25\)，\(|AB|^{2}=2\cdot 32/25=64/25\)，\(|AB|=8/5\) 。）故选 \(C\)。〔卷④自带详解读数校订：源【分析】给 \(a^{2}=4\)、\(b^{2}=1\)、\(c=\sqrt{3}\)、\(y=x-\sqrt{3}\)、\(\Delta =32\)、\(|AB|=\sqrt{2}\cdot \sqrt{}\Delta /5=8/5\)、选\(C\)；源【详解】与【大招指引】用「硬解定理」名目（\(\Delta =4B^{2}mn(A^{2}m+B^{2}n-C^{2})\)、\(|AB|=\sqrt{A^{2}+B^{2}}\cdot \sqrt{}\Delta /(|B||A^{2}m+B^{2}n|)\)），按 §五 名目纪律改为常规联立韦达书写，读数不变；【题后反思】栏之常规法即为上引推导〕〔装配注：题面椭圆方程与选项 \(A\)、\(B\)、\(C\) 数值在原 \(dump\) 为【图】占位，椭圆 \(x^{2}/4+y^{2}=1\) 按源【详解】「由椭圆【图】得 \(a^{2}=4,b^{2}=1\)」补全，选项 \(A\)、\(B\)、\(C\) 数值未回——装配回源图或按读数 \(8/5\) 重排选项（列瑕疵⑤）〕}
\fi
\end{ansblock}

\begin{ansblock}[2章-拓-课时17-02]
% ans:2章-拓-课时17-02
\ansitem{02}{\(\pm 1\)}
\ifansbookdetail
\ansnote{详解}{联立消 \(y\)：\(x^{2}/2+(kx+1)^{2}=1\)，乘 \(2\) 得 \(x^{2}+2(k^{2}x^{2}+2kx+1)=2\)，即 \((1+2k^{2})x^{2}+4kx=0\)。二次项系数 \(1+2k^{2}>0\) 恒成立。设 \(M(x_{1},y_{1})\)、\(N(x_{2},y_{2})\)，则 \(x_{1}+x_{2}=-4k/(1+2k^{2})\)、\(x_{1}x_{2}=0\)（事实上直线过定点 \((0,1)\)，即椭圆上顶点，一交点恒为它）。由 \(|MN|=4\sqrt{2}/3\) 得 \((x_{1}-x_{2})^{2}+(y_{1}-y_{2})^{2}=32/9\)，即 \((1+k^{2})(x_{1}-x_{2})^{2}=32/9\)；而 \((x_{1}-x_{2})^{2}=(x_{1}+x_{2})^{2}-4x_{1}x_{2}=16k^{2}/(1+2k^{2})^{2}\)，代入得 \((1+k^{2})\cdot 16k^{2}/(1+2k^{2})^{2}=32/9\)，约 \(16\) 并交叉相乘：\(9k^{2}(1+k^{2})=2(1+2k^{2})^{2}\)，展开 \(9k^{2}+9k^{4}=2+8k^{2}+8k^{4}\)，即 \(k^{4}+k^{2}-2=0\)，\((k^{2}+2)(k^{2}-1)=0\)，故 \(k^{2}=1\)、\(k=\pm 1\)。（回验 \(k=1\)：\((1+2)x^{2}+4x=0 \rightarrow  x=0\)、\(-4/3\)，\(|MN|=\sqrt{2}\cdot 4/3=4\sqrt{2}/3\) 。）〔卷④自带详解照录校订（源解同路，得 \(k^{4}+k^{2}-2=0\)、\(k=\pm 1\)）〕}
\fi
\end{ansblock}

\begin{ansblock}[2章-拓-课时17-03]
% ans:2章-拓-课时17-03
\ansitem{03}{\(k<-\sqrt{5}/2\) 或 \(k>\sqrt{5}/2\)（\(|k|>\sqrt{5}/2\)）}
\ifansbookdetail
\ansnote{详解}{代入 \(x^{2}-(kx-1)^{2}=4\)，展开 \(x^{2}-k^{2}x^{2}+2kx-1=4\)，整理为 \((1-k^{2})x^{2}+2kx-5=0\)。 ① \(1-k^{2}=0\)（\(k=\pm 1\)）：降为一次方程 \(2kx-5=0\)，有解 \(x=5/(2k)\)，即有公共点，不合「没有公共点」。 ② \(1-k^{2}\neq 0\)：无公共点 \(\Leftrightarrow  \Delta <0\)，\(\Delta =(2k)^{2}-4(1-k^{2})(-5)=4k^{2}+20(1-k^{2})=20-16k^{2}\)。令 \(20-16k^{2}<0\) 得 \(k^{2}>5/4\)，即 \(|k|>\sqrt{5}/2\)。（此时 \(1-k^{2}\neq 0\) 自动满足，因 \(5/4>1\)。） 故 \(k\) 的取值范围为 \(k<-\sqrt{5}/2\) 或 \(k>\sqrt{5}/2\)。（几何校：双曲线 \(x^{2}/4-y^{2}/4=1\) 的渐近线斜率 \(\pm 1\)，直线恒过 \((0,-1)\)；\(|k|>\sqrt{5}/2\approx 1.118>1\) 属比渐近线陡、且与实轴交点在双曲线顶点内侧的情形，无交点自洽。）〔亲算印证：§一 依据格读数 \(|k|>\sqrt{5}/2\)，一致；教材答案条 \(raw\) 乱码段未取〕}
\fi
\end{ansblock}

\begin{ansblock}[2章-拓-课时17-04]
% ans:2章-拓-课时17-04
\ansitem{04}{\((x-1/2)^{2}+(y-1)^{2}=1\) 或 \((x-1/2)^{2}+(y+1)^{2}=1\)（两圆）}
\ifansbookdetail
\ansnote{详解}{\(y^{2}=2x\) 中 \(2p=2\)、\(p=1\)，准线 \(l: x=-1/2\)。设圆心 \(C(a,b)\) 在抛物线上，则 \(b^{2}=2a\)（故 \(a\geqslant 0\)）、半径 \(r\)。与 \(x\) 轴相切 \(\Rightarrow  r=|b|\)；与准线 \(x=-1/2\) 相切 \(\Rightarrow  r=a+1/2\)（\(a\geqslant 0\) 故 \(a-(-1/2)=a+1/2\)）。两式相等：\(|b|=a+1/2\)，以 \(a=b^{2}/2\) 代得 \(|b|=b^{2}/2+1/2\)，即 \(|b|^{2}-2|b|+1=0\)，\((|b|-1)^{2}=0\)，故 \(|b|=1\)、\(b=\pm 1\)，进而 \(a=1/2\)、\(r=1\)。所以圆有两解：\((x-1/2)^{2}+(y-1)^{2}=1\) 与 \((x-1/2)^{2}+(y+1)^{2}=1\)。（回验：圆心 \((1/2,1)\) 在抛物线上 \(1=2\cdot 1/2\) ；到 \(x\) 轴距离 \(1=r\) ；到准线距离 \(1/2+1/2=1=r\) 。）〔亲算印证：双解结构与 §四 表「两圆」读数一致〕}
\fi
\end{ansblock}

\begin{ansblock}[2章-拓-课时17-05]
% ans:2章-拓-课时17-05
\ansitem{05}{\(-3\)}
\ifansbookdetail
\ansnote{详解}{焦点 \(F(1,0)\)，直线 \(l: y=x-1\)。代入 \((x-1)^{2}=4x\) 得 \(x^{2}-2x+1=4x\)，即 \(x^{2}-6x+1=0\)（\(\Delta =36-4=32>0\)）。设 \(A(x_{1},y_{1})\)、\(B(x_{2},y_{2})\)，则 \(x_{1}+x_{2}=6\)、\(x_{1}x_{2}=1\)。又 \(y_{1}y_{2}=(x_{1}-1)(x_{2}-1)=x_{1}x_{2}-(x_{1}+x_{2})+1=1-6+1=-4\)。故 \(OA\cdot OB=x_{1}x_{2}+y_{1}y_{2}=1+(-4)=-3\)。（互核：以 \(y\) 为主元，\(x=y+1\) 代 \(y^{2}=4(y+1)\) 得 \(y^{2}-4y-4=0\)，\(y_{1}y_{2}=-4\) ；\(x_{1}x_{2}=(y_{1}+1)(y_{2}+1)=y_{1}y_{2}+(y_{1}+y_{2})+1=-4+4+1=1\) 。）〔亲算印证：§一 依据格与 §四 表读数 \(-3\)，一致；向量数量积系选修\(1\) 已学内容，不触前引闸〕}
\fi
\end{ansblock}

\begin{ansblock}[2章-拓-课时17-06]
% ans:2章-拓-课时17-06
\ansitem{06}{证明见详解}
\ifansbookdetail
\ansnote{详解}{焦点 \(F(p/2,0)\)。设过 \(F\) 的直线为 \(x=my+p/2\)：此写法已含垂直 \(x\) 轴的弦（\(m=0\)，即通径），而未含的只有 \(y=0\)（抛物线的对称轴），它与抛物线仅一个公共点，与题设「相交（有两个交点）」矛盾，故不损一般性。代入 \(y^{2}=2px\)：\(y^{2}=2p(my+p/2)=2pmy+p^{2}\)，整理为关于 \(y\) 的一元二次方程 \(y^{2}-2pmy-p^{2}=0\)。其二次项系数为 \(1\neq 0\)，且 \(\Delta =(2pm)^{2}+4p^{2}=4p^{2}(m^{2}+1)>0\)（\(p>0\)），故确两交点，与题设相符。由韦达定理，\(y_{1}y_{2}=-p^{2}/1=-p^{2}\)。\(\rule{1.05ex}{1.05ex}\)（特位核验：\(m=0\) 通径时 \(y^{2}=p^{2}\)、\(y=\pm p\)、积 \(-p^{2}\) ；\(m=1\)、\(p=2\) 时 \(y^{2}-4y-4=0\)、积 \(-4=-p^{2}\) ，与拓\(17-05\) 之 \(y_{1}y_{2}=-4\) 同。）〔亲算印证：定义式直推\(+\)韦达，与课时\(16\) 拓\(16-13/14\) 中「焦点弦两端点纵（横）坐标之积为定值」同族；本条原划批\(E\)、批\(E\) 实取复习题未取，回笼本件（§四 表注）〕}
\fi
\end{ansblock}

\begin{ansblock}[2章-拓-课时17-07]
% ans:2章-拓-课时17-07
\ansitem{07}{\(10\)}
\ifansbookdetail
\ansnote{详解}{\(2p=8\)、\(p=4\)，焦点 \(F(2,0)\)、准线 \(x=-2\)。直线 \(l: y=2(x-2)\)。代入 \(y^{2}=8x\) 得 \(4(x-2)^{2}=8x\)，即 \((x-2)^{2}=2x\)，展开 \(x^{2}-4x+4=2x\)，得 \(x^{2}-6x+4=0\)（\(\Delta =36-16=20>0\)）。设 \(A(x_{1},y_{1})\)、\(B(x_{2},y_{2})\)，则 \(x_{1}+x_{2}=6\)。由抛物线定义（焦半径\(=\)到准线距离，准线 \(x=-2\)）\(|AF|=x_{1}+2\)、\(|BF|=x_{2}+2\)，焦点弦 \(|AB|=|AF|+|BF|=x_{1}+x_{2}+4=6+4=10\)。（互核：以 \(y\) 为主元，\(x=y/2+2\) 代 \(y^{2}=8(y/2+2)\) 得 \(y^{2}-4y-16=0\)，\(y_{1}+y_{2}=4\)、\(y_{1}y_{2}=-16\)，\((y_{1}-y_{2})^{2}=16+64=80\)，\(|AB|=\sqrt{}(1+(1/2)^{2})\cdot |y_{1}-y_{2}|=(\sqrt{5}/2)\cdot 4\sqrt{5}=10\) 。）〔亲算印证：两主元路径互核同为 \(10\)；焦点弦公式 \(|AB|=x_{1}+x_{2}+p=6+4=10\) 同〕}
\fi
\end{ansblock}

\begin{ansblock}[2章-拓-课时17-08]
% ans:2章-拓-课时17-08
\ansitem{08}{\((1) |AB|=8\sqrt{5}\)，中点 \((-4,-7)\)；\((2)\) 不成立（\(OA\cdot OB=-15\neq 0\)）}
\ifansbookdetail
\ansnote{详解}{\((1)\) 代入 \(x^{2}=-8(x-3)\)：\(x^{2}=-8x+24\)，即 \(x^{2}+8x-24=0\)（\(\Delta =64+96=160>0\)）。设 \(A(x_{1},y_{1})\)、\(B(x_{2},y_{2})\)，则 \(x_{1}+x_{2}=-8\)、\(x_{1}x_{2}=-24\)。弦长 \(|AB|=\sqrt{1+1^{2}}\cdot |x_{1}-x_{2}|=\sqrt{2}\cdot \sqrt{\Delta }=\sqrt{2}\cdot \sqrt{160}=\sqrt{320}=8\sqrt{5}\)。中点 \(x_{0}=(x_{1}+x_{2})/2=-4\)，\(y_{0}=x_{0}-3=-7\)，即 \((-4,-7)\)。（根之实值 \(x=-4\pm 2\sqrt{10}\)，对应 \(y=-7\pm 2\sqrt{10}\)，代 \(x^{2}=-8y\) 校：\((-4\pm 2\sqrt{10})^{2}=16\mp 16\sqrt{10}+40=56\mp 16\sqrt{10}\)，\(-8(-7\pm 2\sqrt{10})=56\mp 16\sqrt{10}\) 。） \((2) y_{1}y_{2}=(x_{1}-3)(x_{2}-3)=x_{1}x_{2}-3(x_{1}+x_{2})+9=-24+24+9=9\)。故 \(OA\cdot OB=x_{1}x_{2}+y_{1}y_{2}=-24+9=-15\neq 0\)，所以 \(OA\) 与 \(OB\) 不垂直，「\(OA\perp OB\)」不成立。（若按三点式复核：取 \(A(-4+2\sqrt{10}, -7+2\sqrt{10})\)、\(B(-4-2\sqrt{10}, -7-2\sqrt{10})\)，两斜率之积 \(k\_OA\cdot k\_OB=(y_{1}y_{2})/(x_{1}x_{2})=9/(-24)=-3/8\neq -1\)，同得否。）〔亲算印证：\((1) 8\sqrt{5}\)、\((-4,-7)\)；\((2) -15\neq 0\) 不垂直，两式自验一致〕}
\fi
\end{ansblock}

\begin{ansblock}[2章-拓-课时17-09]
% ans:2章-拓-课时17-09
\ansitem{09}{证明见详解}
\ifansbookdetail
\ansnote{详解}{双曲线分两种标准形，证法同，先取 \(C: x^{2}/a^{2}-y^{2}/b^{2}=1\)（\(a>0,b>0\)），其渐近线为 \(y=\pm (b/a)x\)。题设「\(l\) 的斜率与渐近线的斜率相等」即 \(l\) 与某条渐近线平行或重合，不妨设 \(l: y=(b/a)x+m\)（\(m\) 为任意实数，\(m=0\) 时 \(l\) 就是渐近线本身）。代入 \(C\)：\(x^{2}/a^{2}-((b/a)x+m)^{2}/b^{2}=1\)，而 \(((b/a)x+m)^{2}/b^{2}=(x/a+m/b)^{2}=x^{2}/a^{2}+2mx/(ab)+m^{2}/b^{2}\)，故二次项 \(x^{2}/a^{2}\) 相消，得 \(-2mx/(ab)-m^{2}/b^{2}=1\)，即 \(-(2m/(ab))x=1+m^{2}/b^{2}\)。 ① \(m\neq 0\)：一次项系数 \(-2m/(ab)\neq 0\)，方程有唯一解 \(x=-ab(1+m^{2}/b^{2})/(2m)\)，即 \(l\) 与 \(C\) 恰一个公共点。 ② \(m=0\)：等式左端为 \(0\)、右端为 \(1\)，矛盾，无解，即渐近线与 \(C\) 没有公共点。 两种情形公共点均不超过一个，故「最多只有一个公共点」\(\rule{1.05ex}{1.05ex}\)。 焦点在 \(y\) 轴上时（\(C: y^{2}/a^{2}-x^{2}/b^{2}=1\)，渐近线 \(y=\pm (a/b)x\)）同法：设 \(y=(a/b)x+m\) 代入并以 \(x\) 为主元，二次项 \(x^{2}(a^{2}/b^{2})/a^{2}-x^{2}/b^{2}=x^{2}/b^{2}-x^{2}/b^{2}=0\) 相消，化为一次方程或矛盾式，结论同上。\(\rule{1.05ex}{1.05ex}\)（本条与 \(17-11(a)\)、\(17-14\) 之降次情形同一事理，可作其原理说明。）〔亲算印证：消元降次直推，不引后置结论〕}
\fi
\end{ansblock}

\begin{ansblock}[2章-拓-课时17-10]
% ans:2章-拓-课时17-10
\ansitem{10}{三条：\(x=0\)；\(y=p\)；\(y=(1/2)x+p\)（即 \(x-2y+2p=0\)）}
\ifansbookdetail
\ansnote{详解}{分「斜率不存在」「斜率为 \(0\)」「斜率存在且非 \(0\)」三类穷举。 ① 斜率不存在：直线 \(x=0\)（即 \(y\) 轴），代 \(y^{2}=2p\cdot 0=0\) 得 \(y=0\)，唯一公共点 \((0,0)\)（顶点），符合。 ② 斜率为 \(0\)：直线 \(y=p\)（过 \(A(0,p)\) 且平行对称轴 \(x\) 轴），代 \(p^{2}=2px\) 得 \(x=p/2\)，唯一公共点 \((p/2,p)\)，符合——此即 \(17-02/17-05\) 的退化型「一公共点而不切」。 ③ 斜率存在且非 \(0\)：设 \(y=kx+p\)（\(k\neq 0\)），代 \((kx+p)^{2}=2px\)，得 \(k^{2}x^{2}+(2kp-2p)x+p^{2}=0\)（二次项系数 \(k^{2}\neq 0\)）。只有一个公共点 \(\Leftrightarrow  \Delta =0\)：\(\Delta =(2kp-2p)^{2}-4k^{2}p^{2}=4p^{2}(k-1)^{2}-4k^{2}p^{2}=4p^{2}[(k-1)^{2}-k^{2}]=4p^{2}(1-2k)=0\)，故 \(k=1/2\)，切线为 \(y=(1/2)x+p\)，唯一公共点为 \(x=(2p-2kp)/(2k^{2})=(2p-p)/(1/2)=2p\)、\(y=2p\)，即切点 \((2p,2p)\)（校：\((2p)^{2}=4p^{2}=2p\cdot 2p\) ）。 综上满足条件的直线共三条：\(x=0\)、\(y=p\)、\(y=(1/2)x+p\)。（几何校：\(A(0,p)\) 在抛物线外（\(p^{2}>0=2p\cdot 0\)），故恰一条切线；另两条分别由「与对称轴平行」「与准线同向的竖直割线」两个退化方向提供，条数与 §四 表「三条」读数合。）〔亲算印证：三类穷举\(+\Delta =0\)，读数自验；点 \(A(0,p)\) 恰在准线 \(x=-p/2\) 右侧、且在过顶点之切线（\(y\) 轴）上〕}
\fi
\end{ansblock}

\begin{ansblock}[2章-拓-课时17-11]
% ans:2章-拓-课时17-11
\ansitem{11}{\(y=2x-2\)}
\ifansbookdetail
\ansnote{详解}{设 \(l: y=2x+m\)。因抛物线以 \(x=y^{2}/4\) 表出，改以 \(y\) 为主元：由 \(y=2x+m\) 得 \(x=(y-m)/2\)，代入 \(y^{2}=4x=(4(y-m))/2=2(y-m)\)，整理为 \(y^{2}-2y+2m=0\)（二次项系数 \(1\neq 0\)）。相交于两点须 \(\Delta =(-2)^{2}-4\cdot 2m=4-8m>0\)，即 \(m<1/2\)；此时 \(y_{1}+y_{2}=2\)、\(y_{1}y_{2}=2m\)，\((y_{1}-y_{2})^{2}=(y_{1}+y_{2})^{2}-4y_{1}y_{2}=4-8m\)。 弦长用 \(y\) 的差：\(|AB|=\sqrt{}(1+(1/k)^{2})\cdot |y_{1}-y_{2}|=\sqrt{1+1/4}\cdot \sqrt{4-8m}=(\sqrt{5}/2)\sqrt{4-8m}\)。令其等于 \(5\)：\(\sqrt{4-8m}=10/\sqrt{5}=2\sqrt{5}\)，平方得 \(4-8m=20\)，故 \(m=-2\)（\(<1/2\)  合相交条件）。 所以直线 \(l\) 的方程为 \(y=2x-2\)。（回代以 \(x\) 为主元复核：\((2x-2)^{2}=4x \rightarrow  4x^{2}-8x+4=4x \rightarrow  x^{2}-3x+1=0\)，\(|x_{1}-x_{2}|=\sqrt{9-4}=\sqrt{5}\)，\(|AB|=\sqrt{1+4}\cdot \sqrt{5}=\sqrt{5}\cdot \sqrt{5}=5\) 。）〔亲算印证：两主元路径互核同得 \(m=-2\)；单解由 \(\Delta >0\) 与平方可逆性保证，无需舍根外的取舍〕}
\fi
\end{ansblock}

\begin{ansblock}[2章-拓-课时17-12]
% ans:2章-拓-课时17-12
\ansitem{12}{\(25/4\)}
\ifansbookdetail
\ansnote{详解}{\(2p=8\)、\(p=4\)，焦点 \(F(2,0)\)，准线 \(x=-2\)。由 \(A(8,8)\)、\(F(2,0)\) 得直线 \(AB\) 的斜率\(=(8-0)/(8-2)=4/3\)，方程 \(y=(4/3)(x-2)\)。代入 \(y^{2}=8x\)：\((16/9)(x-2)^{2}=8x\)，即 \(2(x-2)^{2}=9x\)，展开 \(2x^{2}-8x+8=9x\)，得 \(2x^{2}-17x+8=0\)（\(\Delta =289-64=225=15^{2}\)）。解得 \(x=(17\pm 15)/4\)，即 \(x=8\)（点 \(A\)）或 \(x=1/2\)，故 \(B(1/2, -2)\)（\(y=(4/3)(1/2-2)=(4/3)(-3/2)=-2\)）。 方法一（中点坐标）：中点横坐标 \(x_{0}=(8+1/2)/2=17/4\)，到准线 \(x=-2\) 的距离\(=17/4+2=25/4\)。 方法二（焦半径\(/\)定义）：\(|AF|=x\_A+p/2=8+2=10\)、\(|BF|=x\_B+p/2=1/2+2=5/2\)，焦点弦上「中点到准线的距离\(=(|AF|+|BF|)/2=|AB|/2\)」（梯形中位线，同 \(16\) 片拓\(16-27\) 型），得 \((10+5/2)/2=(25/2)/2=25/4\) 。 （校：\(y_{1}y_{2}=8\cdot (-2)=-16=-p^{2}\)  与 拓\(17-06\) 结论相符。）〔亲算印证：两法同为 \(25/4\)〕}
\fi
\end{ansblock}

\begin{ansblock}[2章-拓-课时17-13]
% ans:2章-拓-课时17-13
\ansitem{13}{\(x=3\)}
\ifansbookdetail
\ansnote{详解}{设直线为 \(x=x_{0}\)（\(x_{0}\geqslant 0\)，否则与抛物线无交点）。代入 \(y^{2}=4x_{0}\)，得 \(y=\pm 2\sqrt{x}_{0}\)，两交点 \(A(x_{0},2\sqrt{x}_{0})\)、\(B(x_{0},-2\sqrt{x}_{0})\)（关于 \(x\) 轴对称，故该弦被 \(x\) 轴垂直平分）。\(|AB|=2\cdot 2\sqrt{x}_{0}=4\sqrt{x}_{0}=4\sqrt{3}\)，故 \(\sqrt{x}_{0}=\sqrt{3}\)、\(x_{0}=3\)。所以直线 \(AB\) 的方程为 \(x=3\)。（回验：\(y^{2}=12\)、\(y=\pm 2\sqrt{3}\)、\(|AB|=4\sqrt{3}\) ；此即通径型垂直弦，长 \(4\sqrt{3}\) 大于通径 \(2p=4\)，点 \((3,\pm 2\sqrt{3})\) 在曲线上 。）〔亲算印证：读数 \(x=3\) 自验；§一 依据格「\(2|y|=4\sqrt{3} \Rightarrow  y=\pm 2\sqrt{3}\)、\(x=y^{2}/4=3\)」同〕}
\fi
\end{ansblock}

\begin{ansblock}[2章-拓-课时17-14]
% ans:2章-拓-课时17-14
\ansitem{14}{证明见详解}
\ifansbookdetail
\ansnote{详解}{取抛物线标准形 \(y^{2}=2px\)（\(p>0\)），顶点 \(O(0,0)\)、焦点 \(F(p/2,0)\)、准线 \(x=-p/2\)、对称轴为 \(x\) 轴。设 \(P(x_{1},y_{1})\)、\(Q(x_{2},y_{2})\)，\(y_{1}y_{2}\neq 0\)（焦点弦的端点不为顶点，且由 拓\(17-06\) 有 \(y_{1}y_{2}=-p^{2}\neq 0\)）。 直线 \(OP\) 的斜率\(=y_{1}/x_{1}\)，而 \(y_{1}^{2}=2px_{1}\) 故 \(x_{1}=y_{1}^{2}/(2p)\)，斜率\(=y_{1}/(y_{1}^{2}/(2p))=2p/y_{1}\)，即 \(OP: y=(2p/y_{1})x\)。 令 \(x=-p/2\)（准线），得 \(M\) 的纵坐标 \(y\_M=(2p/y_{1})(-p/2)=-p^{2}/y_{1}\)。 由 拓\(17-06\)（焦点弦两端点纵坐标之积 \(y_{1}y_{2}=-p^{2}\)）得 \(y_{2}=-p^{2}/y_{1}\)，于是 \(y\_M=y_{2}\)。 \(M(-p/2, y\_M)\) 与 \(Q(x_{2},y_{2})\) 纵坐标相同，而横坐标不同（\(x_{2}\geqslant 0>-p/2\)），故直线 \(MQ\) 为水平线，即 \(MQ\) 平行于抛物线的对称轴 \(x\) 轴。\(\rule{1.05ex}{1.05ex}\) （通径特位：\(PQ\perp x\) 轴时 \(x_{1}=x_{2}=p/2\)、\(y_{2}=-y_{1}\)，\(y\_M=-p^{2}/y_{1}=-(2p\cdot p/2)/y_{1}=-p^{2}/y_{1}\)，由 \(y_{1}^{2}=2p\cdot (p/2)=p^{2}\) 得 \(y\_M=-y_{1}=y_{2}\)  证法不变；焦点在 \(y\) 轴、开口向左等其他标准形同法（换元即可），结论为「\(MQ\) 平行于该抛物线的对称轴」。）〔亲算印证：以 拓\(17-06\) 之焦点弦纵积定值为引，定义与韦达可直推，不引后置结论〕}
\fi
\end{ansblock}

\begin{ansblock}[2章-拓-课时17-15]
% ans:2章-拓-课时17-15
\ansitem{15}{证明见详解}
\ifansbookdetail
\ansnote{详解}{取标准形 \(y^{2}=2px\)（\(p>0\)），准线 \(l: x=-p/2\)。设 \(P_{1}(x_{1},y_{1})\)、\(P_{2}(x_{2},y_{2})\)，弦 \(P_{1}P_{2}\) 的中点为 \(M\)，则圆心为 \(M\)、半径 \(r=|P_{1}P_{2}|/2\)。 由抛物线定义，焦半径等于点到准线的距离：\(|P_{1}F|=x_{1}+p/2\)、\(|P_{2}F|=x_{2}+p/2\)。因 \(F\) 在线段 \(P_{1}P_{2}\) 上（焦点弦过焦点），\(|P_{1}P_{2}|=|P_{1}F|+|P_{2}F|=x_{1}+x_{2}+p\)，故 \(r=(x_{1}+x_{2}+p)/2\)。 又圆心 \(M\) 的横坐标为 \((x_{1}+x_{2})/2\)，\(M\) 到准线 \(x=-p/2\) 的距离 \(d=(x_{1}+x_{2})/2+p/2=(x_{1}+x_{2}+p)/2\)。 于是 \(d=r\)，即圆心到准线的距离恰等于半径，故以 \(P_{1}P_{2}\) 为直径的圆与准线相切。\(\rule{1.05ex}{1.05ex}\) （另一表述：分别过 \(P_{1}\)、\(P_{2}\)、\(M\) 向准线作垂线，垂足为 \(H_{1}\)、\(H_{2}\)、\(H\)，则 \(|MH|\) 是直角梯形 \(P_{1}H_{1}H_{2}P_{2}\) 的中位线，\(|MH|=(|P_{1}H_{1}|+|P_{2}H_{2}|)/2=(|P_{1}F|+|P_{2}F|)/2=|P_{1}P_{2}|/2=r\)，同得相切——此即 \(17\) 片 \(16-11/\)拓\(16-27\) 反复使用的「中位线\(=\)焦半径和之半」结构。）〔亲算印证：定义\(+\)中位线双表述一致；与 拓\(17-12\) 方法二同族〕}
\fi
\end{ansblock}

\begin{ansblock}[2章-拓-课时17-16]
% ans:2章-拓-课时17-16
\ansitem{16}{\((x-3)^{2}+(y-2)^{2}=16\) 或 \((x-11)^{2}+(y+6)^{2}=144\)（两圆）}
\ifansbookdetail
\ansnote{详解}{焦点 \(F(1,0)\)、准线 \(x=-1\)，直线 \(l: y=x-1\)。代入 \((x-1)^{2}=4x\) 得 \(x^{2}-6x+1=0\)（\(\Delta =32>0\)），设 \(A(x_{1},y_{1})\)、\(B(x_{2},y_{2})\)，则 \(x_{1}+x_{2}=6\)、\(x_{1}x_{2}=1\)，\(x=3\pm 2\sqrt{2}\)。 先定必备数据：弦的中点 \(M=(\)（\(x_{1}+x_{2})/2, (x_{1}+x_{2})/2-1)=(3,2)\)；\(|AB|=\sqrt{2}\cdot |x_{1}-x_{2}|=\sqrt{2}\cdot \sqrt{36-4}=\sqrt{2}\cdot \sqrt{32}=8\)，故 \(|MA|=|MB|=4\)。 所求圆过 \(A\)、\(B\)，其圆心在 \(AB\) 的中垂线上：中垂线斜率 \(-1\)、过 \(M\)，即 \(y=-x+5\)，设圆心 \(C(t, 5-t)\)。 圆的半径按「与准线 \(x=-1\) 相切」为 \(r=|t+1|\)；按「过 \(A\)」为 \(r^{2}=|CA|^{2}=|CM|^{2}+|MA|^{2}=2(t-3)^{2}+16\)（因 \(|CM|^{2}=(t-3)^{2}+(5-t-2)^{2}=2(t-3)^{2}\)）。 令两式相等：\((t+1)^{2}=2(t-3)^{2}+16\)，展开 \(t^{2}+2t+1=2t^{2}-12t+18+16\)，移项得 \(t^{2}-14t+33=0\)，\((t-3)(t-11)=0\)，故 \(t=3\) 或 \(t=11\)（均使 \(r=|t+1|>0\) ，两解皆合）。 \(t=3\)：圆心 \((3,2)=M\)、\(r=4\)，即 \((x-3)^{2}+(y-2)^{2}=16\)——此圆恰为以焦点弦 \(AB\) 为直径的圆，与 拓\(17-15\)（\(B9\) 结论「以焦点弦为直径的圆与准线相切」）相符。 \(t=11\)：圆心 \((11,-6)\)、\(r=12\)，即 \((x-11)^{2}+(y+6)^{2}=144\)。 故满足条件的圆有两个：\((x-3)^{2}+(y-2)^{2}=16\) 与 \((x-11)^{2}+(y+6)^{2}=144\)。（验第二圆：取 \(A(3+2\sqrt{2}, 2+2\sqrt{2})\)，\(|CA|^{2}=(11-3-2\sqrt{2})^{2}+(-6-2-2\sqrt{2})^{2}=(8-2\sqrt{2})^{2}+(-8-2\sqrt{2})^{2}=(72-32\sqrt{2})+(72+32\sqrt{2})=144=r^{2}\) ；圆心到准线距离 \(11+1=12=r\) 。）〔亲算印证：双解 \(t=3\)、\(11\) 系本块自算（教材答案条 \(raw\) 乱码段未取，§三\(.1\)），并以 拓\(17-15\) 之定理与代点计算双向校验〕}
\fi
\end{ansblock}

\begin{ansblock}[2章-拓-课时17-17]
% ans:2章-拓-课时17-17
\ansitem{17}{\(D\)（\(4\)条）}
\ifansbookdetail
\ansnote{详解}{① 斜率不存在：\(l\) 为 \(x=0\)（\(y\) 轴），代入得 \(-y^{2}/3=1\) 无实数解，公共点 \(0\) 个，不合。 ② 斜率存在：设 \(l: y=kx+3\)，代入 \(x^{2}/4-y^{2}/3=1\)（乘 \(12\)）：\(3x^{2}-4(kx+3)^{2}=12\)，展开得 \((3-4k^{2})x^{2}-24kx-48=0\)。 \((a) 3-4k^{2}\neq 0\)（\(l\) 不平行渐近线 \(y=\pm (\sqrt{3}/2)x\)）：一个公共点 \(\Leftrightarrow  \Delta =0\)，\(\Delta =(-24k)^{2}+4(3-4k^{2})\cdot 48=576k^{2}+192(3-4k^{2})=576-192k^{2}\)，令其为 \(0\) 得 \(k^{2}=3\)，\(k=\pm \sqrt{3}\)——两条切线，符合。 \((b) 3-4k^{2}=0\)，即 \(k=\pm \sqrt{3}/2\)：方程降为一次 \(-24kx-48=0\)，一次项系数 \(-24k\neq 0\)（\(k=\pm \sqrt{3}/2\) 时），有唯一解 \(x=-2/k\)，即 \(l\) 平行渐近线时与双曲线恰一个公共点——又两条，符合。 合计 \(2+2=4\) 条，选 \(D\)。（校：\(k=\pm \sqrt{3}\) 时 \(3-4k^{2}=3-12=-9\neq 0\)  两类不重；\(k=\pm \sqrt{3}/2\) 与 \(\pm \sqrt{3}\) 互异  四类无重复。）〔\(3\)章件\(11-\#7\) 自带详解照录校订（源解同路：\((3-4k^{2})x^{2}-24kx-48=0\)、\(\Delta =0\) 得 \(k=\pm \sqrt{3}\)、退化一次得 \(k=\pm \sqrt{3}/2\)，共\(4\)条，选\(D\)）；本块补「斜率不存在」与「一次项系数非零」两步穷举〕〔亲算印证：轮\(2\)亲算③难端档「件\(11-7 4\)条，选\(D\)」与源答案一致〕}
\fi
\end{ansblock}

\begin{ansblock}[2章-拓-课时17-18]
% ans:2章-拓-课时17-18
\ansitem{18}{\(D\)}
\ifansbookdetail
\ansnote{详解}{代入 \(x^{2}-(kx-1)^{2}=1\)，展开 \(x^{2}-k^{2}x^{2}+2kx-1=1\)，整理为 \((1-k^{2})x^{2}+2kx-2=0\)。设两交点横坐标 \(x_{1}\)、\(x_{2}\)，「同在右支且为不同两点」须同时满足： \((i) 1-k^{2}\neq 0\)（保持二次，否则一次方程只一交点）； \((ii) \Delta =(2k)^{2}-4(1-k^{2})(-2)=4k^{2}+8(1-k^{2})=8-4k^{2}>0 \Leftrightarrow  k^{2}<2\)； \((iii)\) 两根之积 \(x_{1}x_{2}=-2/(1-k^{2})>0 \Leftrightarrow  1-k^{2}<0 \Leftrightarrow  k^{2}>1\)（两根同号）； \((iv)\) 两根之和 \(x_{1}+x_{2}=-2k/(1-k^{2})=2k/(k^{2}-1)>0 \Leftrightarrow  k>0\)（结合 \((iii)\) 同为正，即在右支 \(x\geqslant 1\)）。 合 \((ii)(iii)(iv)\)：\(1<k^{2}<2\) 且 \(k>0\)，得 \(1<k<\sqrt{2}\)，即 \((1,\sqrt{2})\)，选 \(D\)。（校：此范围内 \(1-k^{2}\neq 0\) 自动满足；端点 \(k=1\) 使二次项消失（一交点）、\(k=\sqrt{2}\) 使 \(\Delta =0\)（切，一交点），故两端均开，排除 \(B\) 之左闭。）〔\(3\)章件\(11-\#9\) 自带详解照录校订（源解列 \(\{1-k^{2}\neq 0, \Delta >0,\) 和\(>0,\) 积\(>0\}\) 四条件，解得 \(1<k<\sqrt{2}\)，选\(D\)）〕〔亲算印证：轮\(2\)亲算③难端档「件\(11-9 (1,\sqrt{2})\)，选\(D\)」与源答案一致〕}
\fi
\end{ansblock}

\begin{ansblock}[2章-拓-课时17-19]
% ans:2章-拓-课时17-19
\ansitem{19}{\((-2,2)\)}
\ifansbookdetail
\ansnote{详解}{曲线侧：\(2x=\sqrt{4+y^{2}}\) 要求 \(x\geqslant 0\)，两边平方得 \(4x^{2}=4+y^{2}\)，即 \(x^{2}-y^{2}/4=1\) 且 \(x\geqslant 0\)（即该双曲线的右支，含右顶点 \((1,0)\)，渐近线 \(y=\pm 2x\)）。直线 \(y=m(x+1)\) 恒过定点 \((-1,0)\)。代入 \(4x^{2}-y^{2}=4\)：\(4x^{2}-m^{2}(x+1)^{2}=4\)，即 \((4-m^{2})x^{2}-2m^{2}x-(m^{2}+4)=0\)。 ① \(|m|<2\)（\(4-m^{2}>0\)）：两根之积 \(=-(m^{2}+4)/(4-m^{2})<0\)，故有一正根；正根满足 \(x^{2}-y^{2}/4=1\) 且 \(x>0\)，即 \(x\geqslant 1\)，落在右支上 \(\Rightarrow \) 有公共点。 ② \(|m|=2\)：方程降为 \(-2m^{2}x-(m^{2}+4)=0\)，解得 \(x=-(m^{2}+4)/(2m^{2})=-8/8=-1<0\)，不在右支 \(\Rightarrow \) 无公共点。 ③ \(|m|>2\)（\(4-m^{2}<0\)）：积 \(>0\) 且和 \(=2m^{2}/(4-m^{2})<0\)，两根同为负 \(\Rightarrow \) 右支无公共点。 故有公共点 \(\Leftrightarrow  |m|<2\)，即 \(m\in (-2,2)\)。〔\(3\)章件\(11-\#10\) 自带详解照录校订（源解为数形结合：右支、渐近线 \(y=\pm 2x\)、直线过 \((-1,0)\)，斜率须夹在两渐近线之间，得 \(-2<m<2\)；本块补代数穷举①②③以去图依赖，读数同）〕〔亲算印证：轮\(2\)亲算④抽样档「件\(11-10 (-2,2)\)」与源答案一致〕}
\fi
\end{ansblock}

\begin{ansblock}[2章-拓-课时17-20]
% ans:2章-拓-课时17-20
\ansitem{20}{\(D\)}
\ifansbookdetail
\ansnote{详解}{若 \(l\) 斜率不存在（\(x=1\)）：代入得 \(2-y^{2}=2\)，\(y=0\)，只一交点，不合「交于两点且 \(P\) 为中点」；故设 \(l: y=k(x-1)+2\)。代入 \(2x^{2}-y^{2}=2\)：\((2-k^{2})x^{2}+2k(k-2)x-(k^{2}-4k+6)=0\)（其中一次项系数由 \(-2k(2-k)\) 化为 \(2k(k-2)\)、常数项为 \(-(2-k)^{2}-2=-(k^{2}-4k+6)\)）。须 \(2-k^{2}\neq 0\)。 中点条件：\(x_{1}+x_{2}=2 \Rightarrow  -2k(k-2)/(2-k^{2})=2 \Rightarrow  -2k^{2}+4k=4-2k^{2} \Rightarrow  4k=4 \Rightarrow  k=1\)（合 \(2-k^{2}=1\neq 0\) ）。 此时方程为 \(x^{2}-2x-3=0\)，\(x_{1}+x_{2}=2\)、\(x_{1}x_{2}=-3\)，两根 \(3\)、\(-1\)，\(\Delta =4+12=16>0\)  两交点 \(M(3,4)\)、\(N(-1,0)\)（校：\(2\cdot 9-16=2\) ；\(2\cdot 1-0=2\) ）。 \(|MN|=\sqrt{1+k^{2}}\cdot |x_{1}-x_{2}|=\sqrt{2}\cdot \sqrt{}((x_{1}+x_{2})^{2}-4x_{1}x_{2})=\sqrt{2}\cdot \sqrt{4+12}=\sqrt{2}\cdot 4=4\sqrt{2}\)，选 \(D\)。〔\(3\)章件\(12-\#2\) 自带详解照录校订（源解同路：设 \(y-2=k(x-1)\)、\(x_{1}+x_{2}=2x\_P\)、解得 \(k=1\)、\(x_{1}x_{2}=-3\)、\(|MN|=\sqrt{2}\cdot \sqrt{16}=4\sqrt{2}\)，选\(D\)；源 \(dump\) 常数项作「\(-(k^(4)-4k+6)\)」系提取层 \(k^{2}\rightarrow k^(4)\) 讹写，本块按 \(-(k^{2}-4k+6)\) 校正并复核：\(k=1\) 时常数项\(=-3\)，与 \(x^{2}-2x-3=0\) 合 ）〕〔亲算印证：轮\(2\)亲算③难端档「件\(12-2 k=1\)、\(|MN|=4\sqrt{2}\)，选\(D\)」与源答案一致〕}
\fi
\end{ansblock}

\begin{ansblock}[2章-拓-课时17-21]
% ans:2章-拓-课时17-21
\ansitem{21}{\(A\)（\(m=-3\)）}
\ifansbookdetail
\ansnote{详解}{圆 \(C\) 配方法：\((x+1)^{2}+(y+2)^{2}=5-m\)（须 \(m<5\)），圆心 \((-1,-2)\)、半径 \(r=\sqrt{5-m}\)。\(AB\) 是圆 \(C\) 的直径，故圆心即 \(AB\) 的中点：中点 \((-1,-2)\)。 求弦所在直线（按常规联立韦达，不引中点弦公式之名目）：设 \(l\) 斜率为 \(k\)（若 \(l\) 为 \(x=-1\)，代入双曲线得 \(1-y^{2}/2=1\)、\(y=0\)，只一交点，不合），以 \(u=x+1\) 平移（则 \(x=u-1\)、\(y=ku-2\)，中点条件即 \(u_{1}+u_{2}=0\)）。代入 \(2x^{2}-y^{2}=2\)：\(2(u-1)^{2}-(ku-2)^{2}=2 \rightarrow  2u^{2}-4u+2-k^{2}u^{2}+4ku-4=2\)，即 \((2-k^{2})u^{2}+(4k-4)u-4=0\)。由 \(u_{1}+u_{2}=0\) 得 \(-(4k-4)/(2-k^{2})=0\)，故 \(k=1\)（\(2-k^{2}=1\neq 0\) ）。 于是 \(l: y=(x+1)-2=x-1\)，方程化为 \(u^{2}-4=0\)，\(u=\pm 2 \rightarrow  A(1,0)\)、\(B(-3,-4)\)（校：\(1-0=1\) ；\(9-16/2=1\) ），\(|AB|=\sqrt{4^{2}+4^{2}}=4\sqrt{2}\)。 由直径 \(|AB|=\sqrt{D^{2}+E^{2}-4F}\)（此处 \(D=2\)、\(E=4\)、\(F=m\)）：\(|AB|^{2}=4+16-4m=20-4m\)，令 \((4\sqrt{2})^{2}=32=20-4m\)，解得 \(m=-3\)（\(<5\)  圆确存在）。故选 \(A\)。（互核：\(r=|AB|/2=2\sqrt{2}\)、\(r^{2}=8=5-m \Rightarrow  m=-3\) ；又 \(|CA|^{2}=(1+1)^{2}+(0+2)^{2}=8=r^{2}\) 。）〔\(3\)章件\(12-\#19\) 自带详解照录校订（源解：圆心 \(C(-1,-2)\) 为 \(AB\) 中点、「根据双曲线中点差法的结论 \(k\_AB=(b^{2}/a^{2})\cdot (x_{0}/y_{0})=2\cdot (-1)/(-2)=1\)」、\(l: y=x-1\)、解得 \(A(1,0)\)、\(B(-3,-4)\)、\(|AB|=4\sqrt{2}\)、\(\sqrt{4+16-4m}=4\sqrt{2} \Rightarrow  m=-3\)，选\(A\)；名目「中点差法」按 §五 口径改以平移\(+\)韦达写出，\(k\_AB\) 读数同为 \(1\)）〕〔亲算印证：轮\(2\)亲算④抽样档「件\(12-19 m=-3\)，选\(A\)」与源答案一致；\(0913\) 收口轮补「\(3\)章」前缀（§四 表注），与 \(2\)章件\(12\)《圆与圆的位置关系》\(40\)题件无关〕}
\fi
\end{ansblock}

\tailfill[30mm]

\end{multicols}
\end{document}
